% concatenated from main.tex, math_commands.tex

\documentclass[10pt]{article} % For LaTeX2e
%\usepackage{tmlr}
% If accepted, instead use the following line for the camera-ready submission:
\usepackage[accepted]{tmlr}
% To de-anonymize and remove mentions to TMLR (for example for posting to preprint servers), instead use the following:
%\usepackage[preprint]{tmlr}
% Optional math commands from https://github.com/goodfeli/dlbook_notation.
%%%%% NEW MATH DEFINITIONS %%%%%

\usepackage{amsmath,amsfonts,bm}

% Mark sections of captions for referring to divisions of figures
\newcommand{\figleft}{{\em (Left)}}
\newcommand{\figcenter}{{\em (Center)}}
\newcommand{\figright}{{\em (Right)}}
\newcommand{\figtop}{{\em (Top)}}
\newcommand{\figbottom}{{\em (Bottom)}}
\newcommand{\captiona}{{\em (a)}}
\newcommand{\captionb}{{\em (b)}}
\newcommand{\captionc}{{\em (c)}}
\newcommand{\captiond}{{\em (d)}}

% Highlight a newly defined term
\newcommand{\newterm}[1]{{\bf #1}}


% Figure reference, lower-case.
\def\figref#1{figure~\ref{#1}}
% Figure reference, capital. For start of sentence
\def\Figref#1{Figure~\ref{#1}}
\def\twofigref#1#2{figures \ref{#1} and \ref{#2}}
\def\quadfigref#1#2#3#4{figures \ref{#1}, \ref{#2}, \ref{#3} and \ref{#4}}
% Section reference, lower-case.
\def\secref#1{section~\ref{#1}}
% Section reference, capital.
\def\Secref#1{Section~\ref{#1}}
% Reference to two sections.
\def\twosecrefs#1#2{sections \ref{#1} and \ref{#2}}
% Reference to three sections.
\def\secrefs#1#2#3{sections \ref{#1}, \ref{#2} and \ref{#3}}
% Reference to an equation, lower-case.
\def\eqref#1{equation~\ref{#1}}
% Reference to an equation, upper case
\def\Eqref#1{Equation~\ref{#1}}
% A raw reference to an equation---avoid using if possible
\def\plaineqref#1{\ref{#1}}
% Reference to a chapter, lower-case.
\def\chapref#1{chapter~\ref{#1}}
% Reference to an equation, upper case.
\def\Chapref#1{Chapter~\ref{#1}}
% Reference to a range of chapters
\def\rangechapref#1#2{chapters\ref{#1}--\ref{#2}}
% Reference to an algorithm, lower-case.
\def\algref#1{algorithm~\ref{#1}}
% Reference to an algorithm, upper case.
\def\Algref#1{Algorithm~\ref{#1}}
\def\twoalgref#1#2{algorithms \ref{#1} and \ref{#2}}
\def\Twoalgref#1#2{Algorithms \ref{#1} and \ref{#2}}
% Reference to a part, lower case
\def\partref#1{part~\ref{#1}}
% Reference to a part, upper case
\def\Partref#1{Part~\ref{#1}}
\def\twopartref#1#2{parts \ref{#1} and \ref{#2}}

\def\ceil#1{\lceil #1 \rceil}
\def\floor#1{\lfloor #1 \rfloor}
\def\1{\bm{1}}
\newcommand{\train}{\mathcal{D}}
\newcommand{\valid}{\mathcal{D_{\mathrm{valid}}}}
\newcommand{\test}{\mathcal{D_{\mathrm{test}}}}

\def\eps{{\epsilon}}


% Random variables
\def\reta{{\textnormal{$\eta$}}}
\def\ra{{\textnormal{a}}}
\def\rb{{\textnormal{b}}}
\def\rc{{\textnormal{c}}}
\def\rd{{\textnormal{d}}}
\def\re{{\textnormal{e}}}
\def\rf{{\textnormal{f}}}
\def\rg{{\textnormal{g}}}
\def\rh{{\textnormal{h}}}
\def\ri{{\textnormal{i}}}
\def\rj{{\textnormal{j}}}
\def\rk{{\textnormal{k}}}
\def\rl{{\textnormal{l}}}
% rm is already a command, just don't name any random variables m
\def\rn{{\textnormal{n}}}
\def\ro{{\textnormal{o}}}
\def\rp{{\textnormal{p}}}
\def\rq{{\textnormal{q}}}
\def\rr{{\textnormal{r}}}
\def\rs{{\textnormal{s}}}
\def\rt{{\textnormal{t}}}
\def\ru{{\textnormal{u}}}
\def\rv{{\textnormal{v}}}
\def\rw{{\textnormal{w}}}
\def\rx{{\textnormal{x}}}
\def\ry{{\textnormal{y}}}
\def\rz{{\textnormal{z}}}

% Random vectors
\def\rvepsilon{{\mathbf{\epsilon}}}
\def\rvtheta{{\mathbf{\theta}}}
\def\rva{{\mathbf{a}}}
\def\rvb{{\mathbf{b}}}
\def\rvc{{\mathbf{c}}}
\def\rvd{{\mathbf{d}}}
\def\rve{{\mathbf{e}}}
\def\rvf{{\mathbf{f}}}
\def\rvg{{\mathbf{g}}}
\def\rvh{{\mathbf{h}}}
\def\rvu{{\mathbf{i}}}
\def\rvj{{\mathbf{j}}}
\def\rvk{{\mathbf{k}}}
\def\rvl{{\mathbf{l}}}
\def\rvm{{\mathbf{m}}}
\def\rvn{{\mathbf{n}}}
\def\rvo{{\mathbf{o}}}
\def\rvp{{\mathbf{p}}}
\def\rvq{{\mathbf{q}}}
\def\rvr{{\mathbf{r}}}
\def\rvs{{\mathbf{s}}}
\def\rvt{{\mathbf{t}}}
\def\rvu{{\mathbf{u}}}
\def\rvv{{\mathbf{v}}}
\def\rvw{{\mathbf{w}}}
\def\rvx{{\mathbf{x}}}
\def\rvy{{\mathbf{y}}}
\def\rvz{{\mathbf{z}}}

% Elements of random vectors
\def\erva{{\textnormal{a}}}
\def\ervb{{\textnormal{b}}}
\def\ervc{{\textnormal{c}}}
\def\ervd{{\textnormal{d}}}
\def\erve{{\textnormal{e}}}
\def\ervf{{\textnormal{f}}}
\def\ervg{{\textnormal{g}}}
\def\ervh{{\textnormal{h}}}
\def\ervi{{\textnormal{i}}}
\def\ervj{{\textnormal{j}}}
\def\ervk{{\textnormal{k}}}
\def\ervl{{\textnormal{l}}}
\def\ervm{{\textnormal{m}}}
\def\ervn{{\textnormal{n}}}
\def\ervo{{\textnormal{o}}}
\def\ervp{{\textnormal{p}}}
\def\ervq{{\textnormal{q}}}
\def\ervr{{\textnormal{r}}}
\def\ervs{{\textnormal{s}}}
\def\ervt{{\textnormal{t}}}
\def\ervu{{\textnormal{u}}}
\def\ervv{{\textnormal{v}}}
\def\ervw{{\textnormal{w}}}
\def\ervx{{\textnormal{x}}}
\def\ervy{{\textnormal{y}}}
\def\ervz{{\textnormal{z}}}

% Random matrices
\def\rmA{{\mathbf{A}}}
\def\rmB{{\mathbf{B}}}
\def\rmC{{\mathbf{C}}}
\def\rmD{{\mathbf{D}}}
\def\rmE{{\mathbf{E}}}
\def\rmF{{\mathbf{F}}}
\def\rmG{{\mathbf{G}}}
\def\rmH{{\mathbf{H}}}
\def\rmI{{\mathbf{I}}}
\def\rmJ{{\mathbf{J}}}
\def\rmK{{\mathbf{K}}}
\def\rmL{{\mathbf{L}}}
\def\rmM{{\mathbf{M}}}
\def\rmN{{\mathbf{N}}}
\def\rmO{{\mathbf{O}}}
\def\rmP{{\mathbf{P}}}
\def\rmQ{{\mathbf{Q}}}
\def\rmR{{\mathbf{R}}}
\def\rmS{{\mathbf{S}}}
\def\rmT{{\mathbf{T}}}
\def\rmU{{\mathbf{U}}}
\def\rmV{{\mathbf{V}}}
\def\rmW{{\mathbf{W}}}
\def\rmX{{\mathbf{X}}}
\def\rmY{{\mathbf{Y}}}
\def\rmZ{{\mathbf{Z}}}

% Elements of random matrices
\def\ermA{{\textnormal{A}}}
\def\ermB{{\textnormal{B}}}
\def\ermC{{\textnormal{C}}}
\def\ermD{{\textnormal{D}}}
\def\ermE{{\textnormal{E}}}
\def\ermF{{\textnormal{F}}}
\def\ermG{{\textnormal{G}}}
\def\ermH{{\textnormal{H}}}
\def\ermI{{\textnormal{I}}}
\def\ermJ{{\textnormal{J}}}
\def\ermK{{\textnormal{K}}}
\def\ermL{{\textnormal{L}}}
\def\ermM{{\textnormal{M}}}
\def\ermN{{\textnormal{N}}}
\def\ermO{{\textnormal{O}}}
\def\ermP{{\textnormal{P}}}
\def\ermQ{{\textnormal{Q}}}
\def\ermR{{\textnormal{R}}}
\def\ermS{{\textnormal{S}}}
\def\ermT{{\textnormal{T}}}
\def\ermU{{\textnormal{U}}}
\def\ermV{{\textnormal{V}}}
\def\ermW{{\textnormal{W}}}
\def\ermX{{\textnormal{X}}}
\def\ermY{{\textnormal{Y}}}
\def\ermZ{{\textnormal{Z}}}

% Vectors
\def\vzero{{\bm{0}}}
\def\vone{{\bm{1}}}
\def\vmu{{\bm{\mu}}}
\def\vtheta{{\bm{\theta}}}
\def\va{{\bm{a}}}
\def\vb{{\bm{b}}}
\def\vc{{\bm{c}}}
\def\vd{{\bm{d}}}
\def\ve{{\bm{e}}}
\def\vf{{\bm{f}}}
\def\vg{{\bm{g}}}
\def\vh{{\bm{h}}}
\def\vi{{\bm{i}}}
\def\vj{{\bm{j}}}
\def\vk{{\bm{k}}}
\def\vl{{\bm{l}}}
\def\vm{{\bm{m}}}
\def\vn{{\bm{n}}}
\def\vo{{\bm{o}}}
\def\vp{{\bm{p}}}
\def\vq{{\bm{q}}}
\def\vr{{\bm{r}}}
\def\vs{{\bm{s}}}
\def\vt{{\bm{t}}}
\def\vu{{\bm{u}}}
\def\vv{{\bm{v}}}
\def\vw{{\bm{w}}}
\def\vx{{\bm{x}}}
\def\vy{{\bm{y}}}
\def\vz{{\bm{z}}}

% Elements of vectors
\def\evalpha{{\alpha}}
\def\evbeta{{\beta}}
\def\evepsilon{{\epsilon}}
\def\evlambda{{\lambda}}
\def\evomega{{\omega}}
\def\evmu{{\mu}}
\def\evpsi{{\psi}}
\def\evsigma{{\sigma}}
\def\evtheta{{\theta}}
\def\eva{{a}}
\def\evb{{b}}
\def\evc{{c}}
\def\evd{{d}}
\def\eve{{e}}
\def\evf{{f}}
\def\evg{{g}}
\def\evh{{h}}
\def\evi{{i}}
\def\evj{{j}}
\def\evk{{k}}
\def\evl{{l}}
\def\evm{{m}}
\def\evn{{n}}
\def\evo{{o}}
\def\evp{{p}}
\def\evq{{q}}
\def\evr{{r}}
\def\evs{{s}}
\def\evt{{t}}
\def\evu{{u}}
\def\evv{{v}}
\def\evw{{w}}
\def\evx{{x}}
\def\evy{{y}}
\def\evz{{z}}

% Matrix
\def\mA{{\bm{A}}}
\def\mB{{\bm{B}}}
\def\mC{{\bm{C}}}
\def\mD{{\bm{D}}}
\def\mE{{\bm{E}}}
\def\mF{{\bm{F}}}
\def\mG{{\bm{G}}}
\def\mH{{\bm{H}}}
\def\mI{{\bm{I}}}
\def\mJ{{\bm{J}}}
\def\mK{{\bm{K}}}
\def\mL{{\bm{L}}}
\def\mM{{\bm{M}}}
\def\mN{{\bm{N}}}
\def\mO{{\bm{O}}}
\def\mP{{\bm{P}}}
\def\mQ{{\bm{Q}}}
\def\mR{{\bm{R}}}
\def\mS{{\bm{S}}}
\def\mT{{\bm{T}}}
\def\mU{{\bm{U}}}
\def\mV{{\bm{V}}}
\def\mW{{\bm{W}}}
\def\mX{{\bm{X}}}
\def\mY{{\bm{Y}}}
\def\mZ{{\bm{Z}}}
\def\mBeta{{\bm{\beta}}}
\def\mPhi{{\bm{\Phi}}}
\def\mLambda{{\bm{\Lambda}}}
\def\mSigma{{\bm{\Sigma}}}

% Tensor
\DeclareMathAlphabet{\mathsfit}{\encodingdefault}{\sfdefault}{m}{sl}
\SetMathAlphabet{\mathsfit}{bold}{\encodingdefault}{\sfdefault}{bx}{n}
\newcommand{\tens}[1]{\bm{\mathsfit{#1}}}
\def\tA{{\tens{A}}}
\def\tB{{\tens{B}}}
\def\tC{{\tens{C}}}
\def\tD{{\tens{D}}}
\def\tE{{\tens{E}}}
\def\tF{{\tens{F}}}
\def\tG{{\tens{G}}}
\def\tH{{\tens{H}}}
\def\tI{{\tens{I}}}
\def\tJ{{\tens{J}}}
\def\tK{{\tens{K}}}
\def\tL{{\tens{L}}}
\def\tM{{\tens{M}}}
\def\tN{{\tens{N}}}
\def\tO{{\tens{O}}}
\def\tP{{\tens{P}}}
\def\tQ{{\tens{Q}}}
\def\tR{{\tens{R}}}
\def\tS{{\tens{S}}}
\def\tT{{\tens{T}}}
\def\tU{{\tens{U}}}
\def\tV{{\tens{V}}}
\def\tW{{\tens{W}}}
\def\tX{{\tens{X}}}
\def\tY{{\tens{Y}}}
\def\tZ{{\tens{Z}}}


% Graph
\def\gA{{\mathcal{A}}}
\def\gB{{\mathcal{B}}}
\def\gC{{\mathcal{C}}}
\def\gD{{\mathcal{D}}}
\def\gE{{\mathcal{E}}}
\def\gF{{\mathcal{F}}}
\def\gG{{\mathcal{G}}}
\def\gH{{\mathcal{H}}}
\def\gI{{\mathcal{I}}}
\def\gJ{{\mathcal{J}}}
\def\gK{{\mathcal{K}}}
\def\gL{{\mathcal{L}}}
\def\gM{{\mathcal{M}}}
\def\gN{{\mathcal{N}}}
\def\gO{{\mathcal{O}}}
\def\gP{{\mathcal{P}}}
\def\gQ{{\mathcal{Q}}}
\def\gR{{\mathcal{R}}}
\def\gS{{\mathcal{S}}}
\def\gT{{\mathcal{T}}}
\def\gU{{\mathcal{U}}}
\def\gV{{\mathcal{V}}}
\def\gW{{\mathcal{W}}}
\def\gX{{\mathcal{X}}}
\def\gY{{\mathcal{Y}}}
\def\gZ{{\mathcal{Z}}}

% Sets
\def\sA{{\mathbb{A}}}
\def\sB{{\mathbb{B}}}
\def\sC{{\mathbb{C}}}
\def\sD{{\mathbb{D}}}
% Don't use a set called E, because this would be the same as our symbol
% for expectation.
\def\sF{{\mathbb{F}}}
\def\sG{{\mathbb{G}}}
\def\sH{{\mathbb{H}}}
\def\sI{{\mathbb{I}}}
\def\sJ{{\mathbb{J}}}
\def\sK{{\mathbb{K}}}
\def\sL{{\mathbb{L}}}
\def\sM{{\mathbb{M}}}
\def\sN{{\mathbb{N}}}
\def\sO{{\mathbb{O}}}
\def\sP{{\mathbb{P}}}
\def\sQ{{\mathbb{Q}}}
\def\sR{{\mathbb{R}}}
\def\sS{{\mathbb{S}}}
\def\sT{{\mathbb{T}}}
\def\sU{{\mathbb{U}}}
\def\sV{{\mathbb{V}}}
\def\sW{{\mathbb{W}}}
\def\sX{{\mathbb{X}}}
\def\sY{{\mathbb{Y}}}
\def\sZ{{\mathbb{Z}}}

% Entries of a matrix
\def\emLambda{{\Lambda}}
\def\emA{{A}}
\def\emB{{B}}
\def\emC{{C}}
\def\emD{{D}}
\def\emE{{E}}
\def\emF{{F}}
\def\emG{{G}}
\def\emH{{H}}
\def\emI{{I}}
\def\emJ{{J}}
\def\emK{{K}}
\def\emL{{L}}
\def\emM{{M}}
\def\emN{{N}}
\def\emO{{O}}
\def\emP{{P}}
\def\emQ{{Q}}
\def\emR{{R}}
\def\emS{{S}}
\def\emT{{T}}
\def\emU{{U}}
\def\emV{{V}}
\def\emW{{W}}
\def\emX{{X}}
\def\emY{{Y}}
\def\emZ{{Z}}
\def\emSigma{{\Sigma}}

% entries of a tensor
% Same font as tensor, without \bm wrapper
\newcommand{\etens}[1]{\mathsfit{#1}}
\def\etLambda{{\etens{\Lambda}}}
\def\etA{{\etens{A}}}
\def\etB{{\etens{B}}}
\def\etC{{\etens{C}}}
\def\etD{{\etens{D}}}
\def\etE{{\etens{E}}}
\def\etF{{\etens{F}}}
\def\etG{{\etens{G}}}
\def\etH{{\etens{H}}}
\def\etI{{\etens{I}}}
\def\etJ{{\etens{J}}}
\def\etK{{\etens{K}}}
\def\etL{{\etens{L}}}
\def\etM{{\etens{M}}}
\def\etN{{\etens{N}}}
\def\etO{{\etens{O}}}
\def\etP{{\etens{P}}}
\def\etQ{{\etens{Q}}}
\def\etR{{\etens{R}}}
\def\etS{{\etens{S}}}
\def\etT{{\etens{T}}}
\def\etU{{\etens{U}}}
\def\etV{{\etens{V}}}
\def\etW{{\etens{W}}}
\def\etX{{\etens{X}}}
\def\etY{{\etens{Y}}}
\def\etZ{{\etens{Z}}}

% The true underlying data generating distribution
\newcommand{\pdata}{p_{\rm{data}}}
% The empirical distribution defined by the training set
\newcommand{\ptrain}{\hat{p}_{\rm{data}}}
\newcommand{\Ptrain}{\hat{P}_{\rm{data}}}
% The model distribution
\newcommand{\pmodel}{p_{\rm{model}}}
\newcommand{\Pmodel}{P_{\rm{model}}}
\newcommand{\ptildemodel}{\tilde{p}_{\rm{model}}}
% Stochastic autoencoder distributions
\newcommand{\pencode}{p_{\rm{encoder}}}
\newcommand{\pdecode}{p_{\rm{decoder}}}
\newcommand{\precons}{p_{\rm{reconstruct}}}

\newcommand{\laplace}{\mathrm{Laplace}} % Laplace distribution

\newcommand{\E}{\mathbb{E}}
\newcommand{\Ls}{\mathcal{L}}
\newcommand{\R}{\mathbb{R}}
\newcommand{\emp}{\tilde{p}}
\newcommand{\lr}{\alpha}
\newcommand{\reg}{\lambda}
\newcommand{\rect}{\mathrm{rectifier}}
\newcommand{\softmax}{\mathrm{softmax}}
\newcommand{\sigmoid}{\sigma}
\newcommand{\softplus}{\zeta}
\newcommand{\KL}{D_{\mathrm{KL}}}
\newcommand{\Var}{\mathrm{Var}}
\newcommand{\standarderror}{\mathrm{SE}}
\newcommand{\Cov}{\mathrm{Cov}}
% Wolfram Mathworld says $L^2$ is for function spaces and $\ell^2$ is for vectors
% But then they seem to use $L^2$ for vectors throughout the site, and so does
% wikipedia.
\newcommand{\normlzero}{L^0}
\newcommand{\normlone}{L^1}
\newcommand{\normltwo}{L^2}
\newcommand{\normlp}{L^p}
\newcommand{\normmax}{L^\infty}

\newcommand{\parents}{Pa} % See usage in notation.tex. Chosen to match Daphne's book.

\DeclareMathOperator*{\argmax}{arg\,max}
\DeclareMathOperator*{\argmin}{arg\,min}

\DeclareMathOperator{\sign}{sign}
\DeclareMathOperator{\Tr}{Tr}
\let\ab\allowbreak

\usepackage{booktabs} 
\usepackage{wrapfig}
\usepackage{caption}
%\usepackage[table]{xcolor}
\usepackage{subcaption}
\usepackage{enumitem}
\usepackage{hyperref}
\usepackage{url}
\usepackage{amssymb}
%\usepackage{minitoc}
\usepackage{graphicx}
\usepackage{thmtools} 
\usepackage{thm-restate}
\newtheorem{definition}{Definition}
\newtheorem{remark}{Remark}
\newtheorem{lemma}{Lemma}
\newtheorem{assumption}{Assumption}
\newtheorem{proposition}{Proposition }
\newtheorem{condition}{Condition}
\newtheorem{corollary}{Corollary}
\newtheorem{example}{Example}
\newtheorem{theorem}{Theorem}
\newtheorem{proof}{Proof}
\newcommand{\red}[1]{\textcolor{red}{#1}}

\title{Adaptive Gradient Normalization and Independent Sampling for (Stochastic) Generalized-Smooth Optimization}

% Authors must not appear in the submitted version. They should be hidden
% as long as the tmlr package is used without the [accepted] or [preprint] options.
% Non-anonymous submissions will be rejected without review.

\author{\name Yufeng Yang \email yufeng.yang@tamu.edu \\
      \addr Department of Computer Science and Engineering\\
      \addr Texas A\& M University\\
      \addr College Station, TX 77843, USA
      \AND
      \name Erin E. Tripp \email etripp@hamilton.edu \\
      \addr Mathematics and Statistics Department\\
      \addr Hamilton College\\
      \addr Clinton, NY 13323, USA
      \AND
      \name Yifan Sun \email ysun@cs.stonybrook.edu\\
      \addr Department of Computer Science\\
      \addr Stony Brook University \\
      \addr Stony Brook, NY 11794, USA
      \AND
      \name Shaofeng Zou \email 
      zou@asu.edu  \\
      \addr School of Electrical, Computer and Energy Engineering \\
      \addr Arizona State University \\
      \addr Tempe, AZ 85287, USA
      \AND
      \name Yi Zhou \email yi.zhou@tamu.edu \\
      \addr Department of Computer Science and Engineering\\
      \addr Texas A\&M University\\
      \addr College Station, TX 77843, USA
      }

% The \author macro works with any number of authors. Use \AND 
% to separate the names and addresses of multiple authors.

\newcommand{\fix}{\marginpar{FIX}}
\newcommand{\new}{\marginpar{NEW}}

\def\month{07}  % Insert correct month for camera-ready version
\def\year{2025} % Insert correct year for camera-ready version
\def\openreview{\url{https://openreview.net/forum?id=KKSQQMlEfw}} % Insert correct link to OpenReview for camera-ready version


\begin{document}

\maketitle

\begin{abstract}
Recent studies have shown that many nonconvex machine learning problems satisfy a generalized-smooth condition that extends beyond traditional smooth nonconvex optimization. However, the existing algorithms are not fully adapted to such generalized-smooth nonconvex geometry and encounter significant technical limitations on their convergence analysis. In this work, we first analyze the convergence of adaptively normalized gradient descent under function geometries characterized by generalized-smoothness and generalized P{\L} condition, revealing the advantage of adaptive gradient normalization. Our results provide theoretical insights into adaptive normalization across various scenarios.
%introduce  and analyze how $\beta$-normalized gradient descent hyper-parameter choices and learning objective geometry characterized by generalized-smoothness and P{\L} condition impact convergence. 
For stochastic generalized-smooth nonconvex optimization, we propose \textbf{I}ndependent-\textbf{A}daptively \textbf{N}ormalized \textbf{S}tochastic \textbf{G}radient \textbf{D}escent algorithm, which leverages adaptive gradient normalization, independent sampling, and gradient clipping to achieve an $\mathcal{O}(\epsilon^{-4})$ sample complexity under relaxed noise assumptions. Experiments\footnote{Code available at \href{https://github.com/ynyang94/Gensmooth-IAN-SGD}{github.com/ynyang94/Gensmooth-IAN-SGD}} on large-scale nonconvex generalized-smooth problems demonstrate the fast convergence of our algorithm.
\end{abstract}

\section{Introduction}
In modern machine learning, the convergence of gradient-based optimization algorithms has been well studied in the standard smooth nonconvex setting.
However, it has been shown recently that the standard $L$-smoothness fails to hold in many nonconvex machine learning problems, including distributionally-robust optimization(DRO) \citep{levy2020large,jin2021nonconvex}, meta-learning \citep{firstordermeta,metagensmooth} and language models \citep{liu2021pretrain,zhang2020gradientclippingacceleratestraining}. Instead, these problems were shown to satisfy a so-called {\em generalized-smooth} condition \citep{zhang2020gradientclippingacceleratestraining}, in which the smoothness parameter can scale with the
gradient norm in the optimization process.

In the existing literature, various works have proposed different algorithms for solving generalized-smooth nonconvex optimization problems. Specifically, \citet{ li2023convexnonconvexoptimizationgeneralized,zhang2020gradientclippingacceleratestraining,pmlr-v202-chen23ar,petergensmooth,L0L1newstudy,reisizadeh2023variancereducedclippingnonconvexoptimization} have demonstrated that {\em deterministic} first-order algorithms such as gradient descent, normalized gradient descent and clipped gradient descent can achieve $\mathcal{O}(\epsilon^{-2})$ iteration complexity under mild assumptions. These complexity results match the lower bound obtained by classical first-order methods \citep{lowerbound}. 
In particular, \citet{pmlr-v202-chen23ar} empirically demonstrated that proper usage of adaptive gradient normalization can substantially accelerate convergence. However, the formal theoretical justification and understanding of adaptive gradient normalization is still lacking for first-order algorithms in generalized-smooth optimization.

%However, It still lacks theoretical justifications for the role of adaptive normalization on when and how it can make first-order optimization algorithms achieve faster convergence. 


On the other hand, some other works, including \citet{li2023convexnonconvexoptimizationgeneralized,zhang2020gradientclippingacceleratestraining,genralclip} studied first-order {\em stochastic} algorithms in generalized-smooth nonconvex optimization.
Specifically, one line of works \citep{li2023convexnonconvexoptimizationgeneralized} focused on the classic stochastic gradient descent (SGD) algorithm \citep{ghadimi2013stochastic}. However, in generalized-smooth setting, the convergence analysis of SGD either relies on adopting very large batch size or involves large constants \citep{li2023convexnonconvexoptimizationgeneralized}. These theoretical bottlenecks may be a key factor limiting the practical performance of SGD. Empirical observations from \citet{pmlr-v202-chen23ar} indicate that SGD often converges slowly due to ill-conditioned smoothness parameters when the gradient norm is large.
Another line of works \citep{zhang2020gradientclippingacceleratestraining, genralclip,reisizadeh2023variancereducedclippingnonconvexoptimization, usitchclip} focused on clipped SGD, which leverages gradient normalization and clipping to handle the generalized-smooth geometry. Although clipped SGD has demonstrated superior performance in solving large-scale problems,
the existing theoretical analysis has several limitations. First, to establish convergence guarantees in the stochastic setting, existing studies often rely on strong assumptions, such as the stochastic approximation error being almost surely bounded or the use of extremely large batch sizes.
Second, the existing designs of clipped SGD adopt the standard stochastic gradient normalization scheme, which is not fully adapted to the function geometry characterized by the generalized-smooth condition.

Having observed the algorithmic and theoretical limitations discussed above, we aim to advance the algorithm design and analysis for generalized-smooth optimization through investigating the following two fundamental and complementary questions. 
% in this work, we first explore the interplay between adaptive gradient normalization and the convergence of first-order algorithms under nonconvex generalized-smooth geometry. Then, we design stochastic algorithms tailored for large-scale generalized-smooth nonconvex optimization. To achieve this overarching goal, we need to address several fundamental challenges. First, there is limited understanding of how adaptive gradient normalization interplays with function geometry characterized by generalized-smooth on convergence rate.
% Thus, we want to answer the following questions.
\begin{itemize}[leftmargin = *]
    \item \textit{Q1: In deterministic generalized-smooth optimization, how does adaptive gradient normalization affect the convergence rate of first-order algorithm, e.g., under Polyak-{\L}ojasiewicz type conditions?
    }
    %reveal the deep connections among normalization structure, hyper-parameter selection and function landscape geometry characterized by generalized-smooth and other specific condition i.e., generalized Polyak-{\L}ojasiewicz (PL) condition? What is the influence of  hyper-parameter choices and function geometry regarding on convergence rate?
    \item \textit{Q2: In stochastic generalized-smooth optimization, can we design a novel algorithm that guarantees convergence under relaxed noise assumptions while relying on only a small number of samples for computing each stochastic gradient}?
    
\end{itemize}

 %\citep{zhang2020gradientclippingacceleratestraining, genralclip} these studies rely on the strong assumption that the stochastic approximation error is bounded almost surely.

% Furthermore, in stochastic generalized-smooth nonconvex optimization, the existing algorithms are impractical either due to strong assumptions and algorithmic parameters to establish convergence guarantee or poor empirical performance. 
% Therefore, we aim to answer the following question.

% the   for generalized-smooth nonconvex optimization under large-scale setting, it is expected large-batch stochastic samples  may suffer from high computation complexity. Also, it's impractical to introduce and estimate huge approximation error bound when designing algorithms. 
% \begin{itemize}[leftmargin = *]
%     \item \textit{Q2: Can we develop a novel algorithm tailored for stochastic generalized-smooth optimization with convergence guarantee under relaxed nosie assumptions?
%     %not only converges fast under large-scale settings but also keep the small batch size during each iteration to maintain the algorithm efficiency?
%     }
% \end{itemize}
In this work, we provide comprehensive answers to both questions by developing new algorithms and convergence analysis in generalized-smooth nonconvex optimization.
We summarize our contributions as follows.
\subsection{Our Contributions}
    To understand the advantage of using adaptive gradient normalization, we first study the convergence rate of adaptive normalized gradient descent (AN-GD) in deterministic generalized-smooth optimization under the generalized Polyak-{\L}ojasiewicz (P{\L}) condition over a broad spectrum of gradient normalization parameters. Our results reveal the interplay among learning rate, gradient normalization parameter and function geometry parameter, and characterize their impact on the type of convergence rate.
    In particular, our results reveal the advantage of using adaptive gradient normalization and provide theoretical guidance on choosing proper gradient normalization parameter to improve convergence rate.
    
    %of reveals deep connections among learning rate selection, normalization power and  which gives strong intuition in selecting algorithm hyper-parameters.
    
    We further propose a novel Independent-Adaptively Normalized Stochastic Gradient Descent (IAN-SGD) algorithm tailored for stochastic generalized-smooth nonconvex optimization. Specifically, IAN-SGD leverages normalized gradient updates with independent sampling and gradient clipping to reduce bias and enhance algorithm stability. Consequently, we are able to establish convergence of IAN-SGD with $\mathcal{O}(\epsilon^{-4})$ sample complexity under a relaxed assumption on the approximation error of stochastic gradient and constant-level batch size. This makes the algorithm well-suited for solving large-scale problems.
    %, we conduct pilot study and obtain similar implications between learning rate choice of IAN-SGD and function geometry.
    
    We compare the numerical performance of our IAN-SGD algorithm with other state-of-the-art stochastic algorithms in applications of nonconvex phase retrieval, distributionally-robust optimization and training deep neural networks, all of which are generalized-smooth nonconvex problems. Our results demonstrate the efficiency of IAN-SGD in solving generalized-smooth nonconvex problems.
    %and our proposed algorithm outperforms than previous methods in practice.

\section{Related Work}
\textbf{Generalized-Smoothness.}
The concept of generalized-smoothness was first introduced by \citet{zhang2020gradientclippingacceleratestraining} with the $(L_0, L_1)$-smooth condition, which allows a function to either have an affine-bounded Hessian norm or be locally $L$-smooth within a specific region. This definition was extended by \citet{pmlr-v202-chen23ar}, who proposed the $\mathcal{L}_{asym}^*(\alpha)$ and $\mathcal{L}_{sym}^*(\alpha)$ conditions, controlling gradient changes globally with both a constant term and a gradient-dependent term associated with power $\alpha$, thus applying more broadly. Later, \citet{li2023convexnonconvexoptimizationgeneralized} introduced $\ell$-smoothness, which use a non-decreasing sub-quadratic polynomial to control gradient differences. \citet{directionalsmoothness} proposed directional smoothness, which preserves $L$-smoothness along specific directions. 
%these results either require very large batch size or involve large constant factors,

\noindent\textbf{Algorithms for Generalized-Smooth Optimization.} Motivated by achieving comparable lower bounds presented in \citet{lowerbound} under standard assumptions,
algorithms for solving generalized-smooth problems can be categorized into two main series. The first series focuses on gradient descent methods with constant learning rate. \citet{li2023convexnonconvexoptimizationgeneralized} proved that GD, SGD converge with $\mathcal{O}(\epsilon^{-2})$ and $\mathcal{O}(\epsilon^{-4})$ complexity under generalized-smoothness. To ensure convergence, \citet{li2023convexnonconvexoptimizationgeneralized} adopted assumptions such as gradient upper-bound and bounded variance assumption.

Another series of work focuses on adaptive methods. In the nonconvex deterministic settings, \citet{zhang2020gradientclippingacceleratestraining,genralclip,petergensmooth,L0L1newstudy} showed that clipped GD can achieve an iteration complexity of $\mathcal{O}(\epsilon^{-2})$ under mild assumptions. Later, \citet{pmlr-v202-chen23ar} proposed $\beta$-GD that achieves $\mathcal{O}(\epsilon^{-2})$ iteration complexity. Specifically, in the convex setting, \citet{L0L1newstudy} analyzed clipped and normalized gradient descent under mild assumptions. \citet{petergensmooth} studied smoothed gradient clipping, gradient descent with Polyak step-size rule, and triangle method by varying the learning rates. Both works achieve the best-known $\mathcal{O}({\epsilon}^{-1})$ convergence rate for generalized-smooth convex optimization. In the nonconvex stochastic setting, \citet{zhangLsmoothclip} proved that clipped SGD achieves an $\mathcal{O}(\epsilon^{-4})$ sample complexity under the standard $L$-smoothness and bounded variance assumptions. Also, \citet{zhang2020gradientclippingacceleratestraining} established the same convergence rate for clipped SGD under generalized smoothness and almost surely bounded noise assumption. Later, \citet{reisizadeh2023variancereducedclippingnonconvexoptimization, usitchclip} analyzed the convergence of clipped SGD under the bounded variance assumption. Notably, \citet{usitchclip} pointed out that, with bounded variance assumption, clipped SGD converges to $\epsilon$-stationary point with sample complexity $\mathcal{O}(\epsilon^{-5})$ when using $\Omega(1)$ samples per iteration. To improve convergence sample complexity, \citet{reisizadeh2023variancereducedclippingnonconvexoptimization} employed a large batch size $\mathcal{O}(\epsilon^{-2})$ to suppress the impact of the variance term and achieve $\mathcal{O}(\epsilon^{-4})$ complexity. We summarize these results under mild assumptions in Table~\ref{tab: clip-related work}. 

A number of works also investigated the performance of diverse stochastic algorithms—extending beyond SGD and clipped SGD under the generalized-smooth condition.
\citet{adagradnorm,faw2023relax, hongAdagrad} studied AdaGrad \citep{duchiAdagrad} under generalized-smooth and affine variance assumption with different learning rate schemes. They all attain $\mathcal{\tilde{O}}(1/\sqrt{T})$ convergence rate under mild conditions. \citet{trustgensmooth} studied trust-region methods convergence under generalized-smoothness.
For stochastic acceleration methods under the generalized-smoothness condition, \citet{genralclip} proposed a general clipping framework by leveraging momentum clipping and they achieve $\mathcal{O}(\epsilon^{-4})$ sample complexity under almost sure bounded noise assumption; \citet{jin2021nonconvex} studied normalized SGD with momentum under parameter-dependent learning rates schemes, which achieves $\mathcal{O}(\epsilon^{-4})$ sample complexity under generalized-smooth and affine variance assumption; \citet{tuningensmooth} studied normalized SGD with momentum \citep{NSGDm} under parameter-agnostic learning rates schemes, which establishes $\tilde{\mathcal{O}}({\epsilon^{-4}})$ convergence rate under bounded variance assumption.
\citet{pmlr-v202-chen23ar,reisizadeh2023variancereducedclippingnonconvexoptimization} demonstrated that the SPIDER algorithm \citep{fang2018spidernearoptimalnonconvexoptimization} can reach the optimal $\mathcal{O}(\epsilon^{-3})$ sample complexity under affine variance assumption for sum-type functions by employing a large batch size. Furthermore, \citet{qigensmoothadam, huishuaigensmoothadam,huishuaiadam2,adamconvg} explored the convergence of RMSprop \citep{rmsprop} and Adam \citep{kingma2014adam} under different noise assumptions, including coordinate-wise affine variance, expected affine variance, bounded variance and sub-gaussian distribution. The acceleration of sign‑SGD was explored by \citet{signSGDacc,signSGDimprov}, and its generalized‑smooth variants were subsequently analyzed by \citet{mingruisignsgdgs,generalsmoothsignSGD}, where they both obtain $\mathcal{O}(\epsilon^{-4})$ under almost‑sure bound and bounded variance assumptions, respectively.

\begin{table}[htbp]
\resizebox{\textwidth}{!}{
\begin{tabular}{lcccccccc}
\toprule
Method & 
Smoothness$^{1}$ & 
Noise assumption$^{2}$ &
Algorithm$^{3}$ &
Batch Size$^{4}$ &
Convergence$^{5}$ &
Complexity
\\
\midrule
\citet{zhang2020gradientclippingacceleratestraining}       & GS           &    a.s. bound   & Clip-SGD & $\Omega(1)$ & $\epsilon$-stationary & $\mathcal{O}(\epsilon^{-4})$     \\
\citet{zhangLsmoothclip} & LS & var bound & Clip-SGD & $\Omega(1)$ & $\epsilon$-stationary &$\mathcal{O}(\epsilon^{-4})$ \\
\citet{genralclip} & GS & a.s. bound & General-Clip & $\Omega(1)$ & $\epsilon$-stationary & $\mathcal{O}(\epsilon^{-4})$ \\
\citet{li2023convexnonconvexoptimizationgeneralized} & GS & var bound & SGD & $\Omega(1)$ & $\epsilon$-stationary & $\mathcal{O}(\epsilon^{-4})$ \\
\citet{reisizadeh2023variancereducedclippingnonconvexoptimization} & GS & var bound & Clip-SGD & $\Omega(\epsilon^{-2})$ & $\epsilon$-stationary &$\mathcal{O}(\epsilon^{-4})$\\
\citet{usitchclip} & GS & var bound & Clip-SGD & $\Omega(1)$ & $\epsilon$-stationary & $\mathcal{O}(\epsilon^{-5})$ \\
Our Method     &  GS   & a.s. affine & IAN-SGD & $\Omega(1)$ & $\epsilon$-stationary &$\mathcal{O}({\epsilon^{-4}})$    \\
\bottomrule
\end{tabular}}
\caption{Comparison of existing clipping-based algorithms. Denote the stochastic gradient and clipping threshold as $\nabla f_{\xi}(w_t)$ and $c$, respectively. Explanation of abbreviations: (GS) refers to the generalized-smoothness condition, and (LS) refers to the standard $L$-smoothness condition. (a.s. bound) means that for every sample, the stochastic gradient bias is almost surely bounded, i.e., $\| \nabla f_{\xi}(w_t) - \nabla F(w_t) \| \leq \tau_2$; (var bound) denotes the bounded variance condition, i.e., $\mathbb{E}_{\xi}[\| \nabla f_{\xi}(w_t) - \nabla F(w_t) \|^2] \leq \tau_2^2$; (a.s. affine) implies that for every sample, the stochastic gradient bias satisfies the affine bound $\| \nabla f_{\xi}(w_t) - \nabla F(w_t) \| \leq \tau_1 \| \nabla F(w_t) \| + \tau_2$ almost surely, for some $0 \leq \tau_1 < 1$ and $\tau_2 > 0$. The term \textit{Clip-SGD} refers to algorithms that use stochastic gradient clipping of the form $w_{t+1} = w_t - \gamma \min\{1, \frac{c}{\|\nabla f_{\xi}(w_t)\|}\} \nabla f_{\xi}(w_t)$; \textit{SGD} refers to stochastic gradient descent;  \textit{General-Clip} denotes algorithms that incorporate momentum during clipping and updates $w_{t+1} = w_t - ( \nu \min\{c, \frac{\gamma}{\|m_t\|} \} m_t + (1 - \nu) \min\{c, \frac{\gamma}{\|\nabla f_{\xi}(w_t)\|} \} \nabla f_{\xi}(w_t) )$, where $m_{t+1} = (1 - \vartheta) m_t + \vartheta \nabla f_{\xi}(w_t)$. The label $\epsilon$-stationary indicates that the algorithm can converge to an arbitrary target error. For \citet{usitchclip}, we report their convergence result under learning rate $\gamma = \epsilon^{2} / \left(\tau_2^2 (L_0 + cL_1)\right)$ with $c=\tau_2/\epsilon$, where their analysis guarantees convergence to an $\epsilon$-stationary point. }
\label{tab: clip-related work}
\end{table}

\noindent\textbf{Machine Learning Applications.}
Generalized-smoothness has been studied under various machine learning frameworks. \citet{levy2020large,jin2021nonconvex} studied the dual formulation of regularized DRO problems, where the loss function objective satisfies generalized-smoothness. \citet{metagensmooth} identified their meta-learning objective's smoothness constant increases with the norm of the meta-gradient. \citet{bilevelgensmooth1,bilevelgensmooth2,accbilevel,mingruiclip} explored algorithms for bi-level optimization and federated learning within the context of generalized-smoothness. \citet{qimulti-obj} developed algorithms for multi-task learning problem where the objective is generalized-smooth. \citet{onlinegensmooth} studied online mirror descent when the objective is generalized-smooth.
\citet{ziyiminmax} studied min-max optimization algorithms' convergence behavior under generalized-smooth condition. There is a concurrent work \citep{INSGDcompeteversion} using independent sampling with clipped SGD framework to solve variational inequality problem (SVI). Based on this idea, they also propose stochastic Korpelevich method for clipped SGD.
Under generalized-smooth condition, they proved almost-sure convergence in terms of distance to solution set tailored for solving stochastic SVI problems.

\section{Notations}
Throughout the work, we denote $D$ as a set of training examples, $\xi$ as an example sampled from the training dataset $D$, $f(\cdot)$ as the objective function, $f_{\xi}(\cdot)$ as the objective function associated with data sample $\xi$, $F(w)=\mathbb{E}_{\xi\sim \mathbb{P}}[f_{\xi}(w)]$ as the expected loss over all samples following distribution $\mathbb{P}$, and $\|\cdot\|$ as $\ell_2$-norm over Euclidean space.  Let $w\in \mathbf{R}^d$ denote the parameters of function $f$ (i.e., the weights for linear model and neural networks), $f^*$ denote the infimum of $f$. $L_0, L_1$ and $\alpha$ are positive constants used to characterize the generalized-smoothness condition (see Assumption \ref{assum1}). $\rho, \mu$ are positive constant used to describe generalized P{\L}-condition (see Assumption \ref{assum2}). $\tau_1,\tau_2$ characterize the stochastic gradient noise (see Assumption \ref{assum6}).  $\gamma, \beta$ denote the learning rate and normalization power used in adaptive normalization, respectively. In the Distributionally Robust Optimization (DRO) experiments (see Section~\ref{sec: exp}), we define $L(w, \eta)$ with parameter $w$ and dual variable $\eta$, as the dual objective derived from a regularized $\phi$-divergence primal DRO problem.

\section{Generalized-Smooth nonconvex Optimization}
We first introduce generalized-smooth optimization problems. Consider the following optimization problem.
\begin{align}
    \min_{w \in \mathbf{R}^d} f(w), 
\end{align}
where $f: \mathbf{R}^d \rightarrow \mathbf{R}$ denotes a nonconvex and differentiable function and $w$ corresponds to the model parameters. We assume that function $f$ satisfies the following generalized-smooth condition. 
%introduced in \citep{jin2021nonconvex,pmlr-v202-chen23ar}.
%\footnote{\citet{jin2021nonconvex}considered the special case $\alpha=1$, and \citet{pmlr-v202-chen23ar} defined a symmetric version of \eqref{eq: gs}.}
\begin{assumption}[Generalized-smooth]\label{assum1}
    The objective function $f$ satisfies the following conditions.
    \begin{enumerate}
        \item $f$ is differentiable and bounded below, i.e., $f^* := \inf_{w\in \mathbf{R}^d} f(w) > -\infty$;
        \item There exists constants $L_0, L_1>0$ and $\alpha\in [0,1]$ such that for any $w$, $w' \in \mathbf{R}^d$, it holds that
        %\begin{align}
            %\big \|\nabla f(w')- \nabla f(w) \big \|\leq (L_0 +L_1\cdot max_{\theta\in [0,1]}\big \|\nabla f(w_{\theta}(w',w)) \big \|^{\alpha})\big \|w'-w \big \|,
        %\end{align}
        %where $w_{\theta}(w',w)=\theta w'+(1-\theta)w$
        \begin{align}\label{eq: gs}            \big \|\nabla f(w)- \nabla f(w') \big \|
        \leq\big(L_0+L_1 \big \|\nabla f(w') \big \|^{\alpha}\big)\big \|w - w' \big \|.
        \end{align}
    \end{enumerate}
\end{assumption}
The generalized-smooth condition in Assumption \ref{assum1} is a generalization of the standard smooth condition, which corresponds to the special case of $\alpha=0$. It allows the smoothness parameter to scale with the gradient norm polynomially, and therefore Assumption \ref{assum1} can model functions with highly irregular geometry. Moreover, following the standard proof, it is easy to show that generalized-smooth functions satisfy the following descent lemma. See Proof \ref{proof of descent lemma} in Appendix \ref{Appendix B}.

\begin{restatable}{lemma}{mainlemmaone}
\label{descent lemma}
Under Assumption \ref{assum1}, function $f$ satisfies, for any $w, w' \in \mathbf{R}^d$,
    \begin{align}
        f(w)\leq f(w')+ \langle \nabla f(w'), w - w' \rangle
        + \frac{1}{2} (L_0+L_1 \big \| \nabla f(w')\big \|^{\alpha}) \big \| w - w' \big \|^2.
        \label{origindescent}
    \end{align}
\end{restatable}
 
We note that there are several variants of generalized-smooth conditions proposed by the previous works \citep{zhang2020gradientclippingacceleratestraining,jin2021nonconvex,pmlr-v202-chen23ar}. Below, we briefly discuss the relationship between the generalized-smooth condition in Assumption \ref{assum1} and these existing notions.
\begin{remark}[$(L_0, L_1)$-generalized-smooth condition]
The descent lemma of $(L_0, L_1)$-generalized-smooth condition proposed in \citet{zhang2020gradientclippingacceleratestraining} is given as
\begin{align}
    f(w)\leq f(w')+ \langle \nabla f(w'), w - w' \rangle +\frac{1}{2}(4L_0+5L_1 \| \nabla f(w')\|) \| w - w' \|^2,\nonumber 
\end{align}
which is the same as Lemma \ref{descent lemma} with $\alpha = 1$ up to differences in the constant coefficients. 

\end{remark}
\begin{remark}[Symmetric generalized-smooth condition]
\citet{pmlr-v202-chen23ar} introduced a symmetric version of generalized-smooth condition, by replacing $\|\nabla f(w') \|^{\alpha}$ in \eqref{eq: gs} with $\max_{w_{\theta}}\|f(w_{\theta})\|^{\alpha}$, where $w_{\theta}=\theta w'+(1-\theta)w$. 
We notice that the asymmetric generalized-smooth condition adopted in our Assumption \ref{assum1} can also imply the following similar symmetric generalized-smooth condition,
\begin{align}
     \big \|\nabla f(w)- \nabla f(w') \big \|\leq \big(L_0 +L_1 \big(\frac{\big \|\nabla f(w) \big \|^{\alpha}+\big \|\nabla f(w') \big \|^{\alpha}}{2}\big)\big)\big \|w - w' \big \|.
     \label{symmetric gs}
\end{align}
%Compared \eqref{symmetric gs} with symmetric generalized-smooth proposed in \citep{pmlr-v202-chen23ar}, \eqref{symmetric gs} is easier to justify and it can still characterize many complex machine learning problems.
\end{remark}

In \citet{jin2021nonconvex}, it has been shown that certain regularized nonconvex distributionally robust optimization (DRO) problems satisfy the generalized-smooth condition in \eqref{eq: gs} with $L_0=L$, $L_1 = \frac{2M(G+1)^2}{\lambda}$, $\alpha = 1$. Moreover, the classic nonconvex phase retrieval problem \citep{PhaseRetrieval1,PhaseRetrieval2} can be shown to satisfy the above symmetric generalized-smooth condition in \eqref{symmetric gs} (see Proof \ref{proof of Phase Retrieval} in Appendix \ref{Appendix C}). 

% \begin{restatable}{proposition}{phaseretrievalsmooth}
%     The nonconvex phase retrieval problem (see \eqref{noncvxphaseretrieval}) satisfies the symmetric generalized smooth condition 
%     \begin{align}
%         \big \|\nabla F(w)- \nabla F(w') \big \|\leq \frac{1}{m}\sum_{\xi\in D}(L_0+L_1\frac{\|\nabla f_{\xi}(w)\|^{\alpha}+\| \nabla f_{\xi}(w')\|^{\alpha}}{2})\big \|w-w' \big \|
%         \label{symmetric gs stochastic}
%     \end{align}
%     with $L_0=4y_{\max}a_{\max}^2$, and $L_1=\frac{9}{2}a^{\frac{4}{3}}_{\max}$, $\alpha=\frac{2}{3}$.
% \end{restatable}
% \begin{restatable}{proposition}{DROsmooth}
%     Assume the objective of distributionally robust optimization (see \eqref{dualnoncvxDRO} in experiment section) satisfies that $\Psi^*$ is M-smooth, and $\ell_{\xi}(w)$ is G-Lipschitz continuous and $L$-smooth, then it satisfies \eqref{symmetric gs} with $L_0=L$, $L_1 = \frac{2M(G+1)^2}{\lambda}$, $\alpha = 1$.
% \end{restatable}
% (See Proof \ref{proof of Phase Retrieval} in Appendix \ref{Appendix C})
% Thus, throughout this work, we consider the generalized-smooth condition in Assumption \ref{assum1}, which generalizes the existing notions and simplifies arithmetic of the proof. 

In the following sections, we first consider deterministic generalized-smooth optimization and study the impact of adaptive gradient normalization on the convergence rate of gradient methods. Then, we consider stochastic generalized-smooth optimization and propose a novel independent sampling scheme for improving the convergence guarantee of stochastic gradient methods under relaxed noise assumption.

\section{Adaptive Gradient Normalization for Deterministic Generalized-Smooth Optimization}
In deterministic generalized-smooth optimization, several previous works have empirically demonstrated the faster convergence of normalized gradient descent-type algorithms (e.g., clipped GD) over the standard gradient descent algorithm in various machine learning applications \citep{zhang2020gradientclippingacceleratestraining,pmlr-v202-chen23ar}. On the other hand, theoretically, these algorithms were only shown to achieve the same iteration complexity $\mathcal{O}(\epsilon^{-2})$ as the gradient descent algorithm in generalized-smooth optimization.
In this section, to further advance the theoretical understanding and explain the inconsistency between theory and practice, we explore the advantage of adapting gradient normalization to the special Polyak-{\L}ojasiewicz-type (P{\L}) geometry in generalized-smooth optimization. We aim to show that gradient normalization, when properly adapted to the underlying P{\L} geometry, can help accelerate the convergence rate in generalized-smooth optimization.

Specifically, we consider the class of generalized-smooth problems that satisfy the following generalized P{\L} geometry.
\begin{assumption}[Generalized Polyak-{\L}ojasiewicz Geometry]\label{assum2}
       There exists constants $\mu\in \mathbf{R}_+$ and $0<\rho\leq 2$ such that $f(\cdot)$ satisfies, for all $w\in \mathbf{R}^d$,
    \begin{align}
        \big \|\nabla f(w) \big \|^{\rho} \geq 2\mu( f(w)-f^*).
        \label{gen PL}
    \end{align} 
\end{assumption}
The above generalized P{\L} condition is inspired by the Kurdyka-{\L}ojasiewicz (K\L)-exponent condition proposed in \citet{KLexponent}. When $\rho > 1$, \eqref{gen PL} reduces to the K\L-exponent condition. When $\rho=2$, \eqref{gen PL} reduces to the standard P{\L} condition. Moreover, some recent works have shown that P{\L}-type geometries widely exist in the loss landscape of phase retrieval \citep{PhaseretrievalPL} and over-parametrized deep neural networks \citep{genPL1,genPL2}, and we hope that our analysis based on Assumption \ref{assum2} will allow researchers to rethink the relationship between adaptive normalization and loss landscape geometry.

Here, we consider the adaptively normalized gradient descent (AN-GD) algorithm proposed by \citet{pmlr-v202-chen23ar} for generalized-smooth nonconvex optimization. The algorithm normalizes the gradient update as follows
\begin{align}\label{eq: AN-GD}
    \text{(AN-GD)} \quad w_{t+1} = w_t - \gamma \frac{\nabla f(w_t)}{ \| \nabla f(w_t) \|^{\beta}},
\end{align}
where $\gamma>0$ denotes the learning rate and $\beta$ is a normalization scaling parameter that allows us to adapt the normalization scale of the gradient norm to the underlying function geometry. Intuitively, when the gradient norm is large, a smaller $\beta$ would make the normalized gradient update more aggressive; when the gradient norm is small, normalization can slow down gradient vanishing and improve numerical stability. 

\citet{pmlr-v202-chen23ar} studied AN-GD in generalized-smooth nonconvex optimization, showing that it achieves the standard $\mathcal{O}(\epsilon^{-2})$ iteration complexity lower bound. In the following theorem, we obtain the convergence rate of AN-GD under the generalized P{\L} condition. See Proof \ref{Proof of beta GD convergence} in Appendix \ref{Appendix E}. 

\begin{restatable}[Convergence of AN-GD]{theorem}{mainthmone}\label{beta GD convergence}
    Let Assumptions \ref{assum1} and \ref{assum2} hold. Denote $\Delta_t:=f(w_t) - f^*$ as the function value gap. Choose target error $0< \epsilon\leq \min\{1,\frac{1}{2\mu}\}$ and define learning rate $\gamma = \frac{(2\mu\epsilon)^{\beta/\rho}}{8(L_0+L_1)+1}\leq \frac{2}{\mu}$ for some $\beta\in[\alpha,1]$. Then, the following statements hold,  
    \begin{itemize}[leftmargin = *]
        \item %If $\beta<2-\rho$, %denote $C=\min\Big\{\frac{\gamma(2\mu)^{\frac{2-\beta}{\rho}}}{8},\frac{\rho}{2-\beta-\rho}(2^{\frac{2-\beta-\rho}{2-\beta}}-1)\Delta_0^{\frac{2-\beta-\rho}{\rho}}  \Big\}$.
        %then we have recursion
        %\begin{align}
        %\Delta_{t} = \mathcal{O}\Big(\big(\frac{\rho}{(2-\beta-\rho)\gamma t}\big)^{\frac{\rho}{2-\rho-\beta}}\Big).%previously \mathcal{O}
        %\label{thm1_recursion1}
        %\end{align}
    If $\beta<2-\rho$, in order to achieve $\Delta_T\leq \epsilon$, 
    the total number of iterations must satisfy
    \begin{align}
    T \geq\max \Big\{\frac{8\rho(8(L_0+L_1)+1)}{(2-\beta-\rho)(2\mu)^{2/\rho}\epsilon^{(2-\rho)/{\rho}}}, \frac{1}{(2^{(2-\beta-\rho)/(2-\beta)}-1)\Delta_0^{{(\rho+\beta-2)}/{\rho}}\epsilon^{(2-\beta-\rho)/{\rho}}} \Big\}= \Omega \big((\frac{1}{\epsilon})^{\frac{2-\rho}{\rho}}\big).
    \end{align}
    % \Omega
        \item %If $\beta=2-\rho$ and choose $\epsilon$ such that $\gamma < \frac{2}{\mu}$, then we have recursion 
        %\begin{align}
        %    \Delta_t = \mathcal{O}\Big(\big(1-\frac{\gamma\mu}{2}\big)^t \Big). % \mathcal{O}
            %\label{thm1_recursion2}
        %\end{align}
        If $\beta=2-\rho$, in order to achieve $\Delta_{T}\leq \epsilon$,
        the total number of iterations must satisfy
        \begin{align}
        T\geq\frac{2^{1-\beta/\rho}\cdot(8(L_0+L_1)+1)}{\mu\cdot(\mu\epsilon)^{\beta/\rho}}\cdot\log(\frac{\Delta_0}{\epsilon})=\Omega\big((\frac{1}{\epsilon})^{\frac{\beta}{\rho}}\cdot \log(\frac{\Delta_0}{\epsilon})\big).
        \end{align}
        %\Omega previously
        \item %If $1\ge\beta>2-\rho$, then there exists $T_0\in \mathbf{N}$ satisfying $\Delta_{T_0}\leq {4}^{-\frac{\rho}{\rho+\beta-2}}(2\mu)^{\frac{2-\beta}{\rho+\beta-2}}\gamma^{\frac{\rho}{(\rho+\beta-2)}}$ and, for all $t\geq T_0$, we have 
        %    \begin{align}
        %    \Delta_{t} = \mathcal{O}\Big( \Big(\frac{\Delta_{T_0}}{\gamma^{\frac{\rho}{\rho+\beta-2}}}\Big)^{({\frac{\rho}{2-\beta}})^{t-T_0}}\Big). 
            %\label{2ndstage}
            %previously \mathcal{O}
            %\label{thm1_recursion3}
        %\end{align}
        {If $1\ge\beta>2-\rho$, and for some $T_0$ satisfying $\Delta_{T_0}\leq \mathcal{O}\big((\gamma\mu^{(2-\beta)/(\rho+\beta-2)})^{\frac{\rho}{\rho+\beta-2}}\big)$, then
        the total number of iterations after $T_0$ must satisfy}
        \begin{align}
        T\gtrsim&\Omega\Big(\log\Big(\big(\frac{1}{\epsilon}\big)^{\frac{\beta}{\rho+\beta-2}}\Big)+\log\Big(\log\Big(\frac{(2\mu)^{2/(\rho+\beta-2)}}{(32(L_0+L_1)+4)^{\rho/(\rho+\beta-2)}\epsilon}\Big)\Big)\Big)=\Omega \big(\log\big((\frac{1}{\epsilon})^{\frac{\beta}{\rho+\beta-2}}\big)\big).
        \end{align}
    \end{itemize}
\end{restatable}
Theorem \ref{beta GD convergence} characterizes the convergence rates of AN-GD under different choices of $\rho$ and $\beta$. In particular, when $\alpha=\beta=0$ and $\rho=2$, Theorem \ref{beta GD convergence} reduces to the linear convergence rate achieved by gradient descent under the standard P{\L} and $L$-smooth condition. 
Theorem \ref{beta GD convergence} also guides the choice of gradient normalization parameter $\beta$ under different geometric conditions. For example, if there exists $\beta \in [\alpha, 1]$ such that $\rho+\beta>2$, AN-GD can boost its convergence rate from polynomial to local linear convergence. It also indicates when $\rho$ is very small such that $\rho+\beta<2$, the effects of $\beta$ can be marginal. However, for $\rho+\beta\leq 2$, the iteration complexities depend more on the specific values of $\rho$ and $\beta$.
For example, 
when $\rho=1$ and consider two different choices of gradient normalization parameters $\beta_1=\frac{2}{3},\beta_2=1$, AN-GD achieves the iteration complexities $\mathcal{O}({\epsilon}^{-1})$ and $\tilde{\mathcal{O}}({\epsilon}^{-1})$, respectively. This result matches the empirical observation made in \citep{pmlr-v202-chen23ar} that choosing a smaller $\beta\in [\alpha,1]$  can improve the convergence rate.
%When $\rho=3$ and vary $\beta_1=\frac{2}{3},\beta_2 = 1$, AN-GD achieves asymptotic iteration complexities $\Omega(\log(\epsilon^{-0.4}))$ and $\Omega(\log(\epsilon^{-0.5}))$, respectively. For both cases, choosing $\beta=\frac{2}{3}$ leads to faster convergence. 

We note that several concurrent works also studied normalized gradient descent methods in deterministic convex settings. In \citet{petergensmooth}, the authors consider various update rules such as
$
w_{t+1} = w_t - \gamma \tfrac{\nabla f(w_t)}{L_0 + L_1 \|\nabla f(w_t)\|} \quad \text{and} \quad w_{t+1} = w_t - \gamma \tfrac{f(w_t) - f(w^*)}{\|\nabla f(w_t)\|} \nabla f(w_t).$
For generalized smooth convex optimization, they obtain a two-phase convergence result, similar to the case when $1 \geq \beta > 2 - \rho$ in Theorem~\ref{beta GD convergence}, where the second phase achieves complexity bound $\mathcal{O}(\epsilon^{-1})$. They further employ the triangle method to derive a global $\mathcal{O}(\epsilon^{-1})$ complexity bound. In \citet{L0L1newstudy}, the authors propose a related framework by choosing $\gamma = \frac{1}{\sqrt{T}\|\nabla f(w_t) \|}$, which aligns with our AN-GD when $\beta = 1$. They also obtain $\mathcal{O}(\epsilon^{-1})$ complexity for generalized smoothness convex optimization.


\subsection{Comparison between AN-GD and GD.}
In \citet{li2023convexnonconvexoptimizationgeneralized}, it has been shown that GD achieves $\mathcal{O}(\log({\epsilon}^{-1}))$, $\mathcal{O}({\epsilon}^{-1})$, $\mathcal{O}({\epsilon^{-2}})$ complexity bounds for generalized-smooth strongly convex, convex and non-convex optimization, respectively. 
Compared with \citet{pmlr-v202-chen23ar}, both AN-GD and GD achieve $\mathcal{O}({\epsilon^{-2}})$ complexity bound for generalized non-convex optimization,
and hence are optimal deterministic methods. On the other hand, the following Proposition \ref{GD convergence} further obtains the convergence rate of gradient descent (GD) under the same Assumptions \ref{assum1} and \ref{assum2}. See Proof \ref{Proof of GD convergence} in Appendix \ref{Appendix F}.
\begin{restatable}[Convergence of GD]{proposition}{gdconvergence}\label{GD convergence}
Let Assumptions \ref{assum1} and \ref{assum2} hold. Assume there exists a positive constant $G$ such that $\|\nabla f(x_t) \|\leq G, \forall t\in T$, 
    When $\rho=\alpha=1$ and setting $\gamma \leq \min \{\frac{1}{L_0}, \frac{1}{2L_1G} \}$, gradient descent converges to an $\epsilon$-stationary point within $T\geq\Omega(\frac{G}{\epsilon})$ iterations.
\end{restatable}
 From Theorem~\ref{beta GD convergence}, under the setting $\rho = \alpha = 1$ and with $\beta = 1$, AN-GD converges to a $\epsilon$-stationary point in $\tilde{\mathcal{O}}({\epsilon}^{-1})$ iterations.  However, the convergence guarantee for GD requires an additional assumption that the gradient norm is upper bounded by a constant $G$. This constant typically depends on the function value gap $\Delta_0$ and $\|x_0 - x^*\|$, which can be large in generalized smooth optimization settings. Such large value can be the key reason for the slower empirical convergence of GD observed in \citet{pmlr-v202-chen23ar}.


\section{Independently-and-Adaptively Normalized SGD for Stochastic Generalized-Smooth Optimization}
In this section, we study stochastic generalized-smooth optimization problems, where we denote $f_\xi$ as the loss function associated with the data sample $\xi$, and we assume that the following expected loss function $F(\cdot)$ satisfies the generalized-smooth condition in Assumption \ref{assum1}.
\begin{align}\label{eq: sgs}
    \min_{w\in \mathbf{R}^d} F(w) = \mathbb{E}_{\xi\sim \mathbb{P}}\big[f_{\xi}(w) \big].
\end{align}
Having discussed the superior theoretical performance of AN-GD in the previous section, we aim to leverage adaptive gradient normalization to further develop an adaptively normalized algorithm tailored for stochastic generalized-smooth optimization.
\subsection{Normalized SGD and Its Limitations}
%Correspondingly, we assume $\mathcal{L}_{\text{asym}}^*(\alpha)$ generalized-smooth condition for stochastic programming.
To solve the stochastic generalized-smooth problem in \eqref{eq: sgs}, one straightforward approach is to replace the full batch gradient in the AN-GD update rule, \eqref{eq: AN-GD} with the stochastic gradient $\nabla f_\xi(w_t)$, resulting in the following adaptively normalized SGD (AN-SGD) algorithm.
\begin{align}\label{eq: NSGD}
    \text{(AN-SGD)} \quad w_{t+1} = w_t - \gamma \frac{\nabla f_{\xi}(w_t)}{ \| \nabla f_{\xi}(w_t) \|^{\beta}}.
\end{align}
AN-SGD-type algorithms have attracted a lot of attention recently for solving stochastic
generalized-smooth problems \citep{zhang2020gradientclippingacceleratestraining,genralclip,mingruiclip,jin2021nonconvex}. 
% \begin{wrapfigure}{r}{0.22\textwidth}
%     \vspace{-3mm}
%     \centering
%     \includegraphics[width=0.22\textwidth]{Figure_illustration.pdf}
%     \vspace{-3mm}
%     \caption{Comparison between normalized full gradient (blue) and expected normalized stochastic gradient (red). Here, $\xi_1$ and $\xi_2$ are sampled uniformly at random.}
%     \label{figure illustration}
%     %\vspace{-10pt} % To adjust space below the figure
%     \vspace{-3mm}
% \end{wrapfigure}
In particular, previous works have shown that, when choosing $\beta=1$ and using gradient clipping or momentum acceleration, AN-SGD's variations (e.g., NSGD with momentum \citep{NSGDm, jin2021nonconvex}, clipped SGD \citep{zhang2020gradientclippingacceleratestraining,genralclip}) can achieve a sample complexity of $\mathcal{O}(\epsilon^{-4})$. This result matches the convergence order of the standard SGD \citep{ghadimi2013stochastic} for solving classic stochastic smooth problems. 
However, AN-SGD has several limitations as summarized below. 
\begin{itemize}[leftmargin = *]
    \item {\em Biased gradient estimator:} The normalized stochastic gradient used in \eqref{eq: NSGD} is biased, i.e., 
    $
        \mathbb{E}[\frac{\nabla f_{\xi}(w_t)}{\| \nabla f_{\xi}(w_t) \|^{\beta}} ] \ne \frac{\nabla F(w_t)}{\| \nabla F(w_t)\|^{\beta}}.
    $
    This is due to the dependence between $\nabla f_{\xi}(w_t)$ and  $\| \nabla f_{\xi}(w_t)\|^{\beta}$. In particular, the bias can be huge if the stochastic gradients are diverse, as illustrated in Figure \ref{figure illustration}.
    %needs strong numerical trick to upper bound normalization term by a constant or gradient norm. bounded approximation error or extremely large batch size. 
    %However, those tricks are no longer valid under presence of $L_{\text{asym}}^*(\alpha)$ generalized-smoothness since it's hard to align the order of gradient norm in the presence of $\alpha$. 
    \begin{figure}
        \centering
        \includegraphics[width=0.4\linewidth]{Figure_illustration.pdf}
        \caption{Comparison between normalized full gradient (blue) and expected normalized stochastic gradient (red). Here, $\xi_1$ and $\xi_2$ are sampled uniformly at random.}
        \label{figure illustration}
    \end{figure}
    \item {\em Strong assumptions:} To control the estimation bias and establish theoretical convergence guarantee for ANSGD-type algorithms in generalized-smooth nonconvex optimization, the existing studies need to adopt strong assumptions. For example, \citet{zhang2020gradientclippingacceleratestraining, genralclip} and \citet{mingruiclip} assume that the stochastic approximation error $\|\nabla f_\xi(w) - \nabla F(w)\|$ is bounded by a constant almost surely.
    In practical applications, this constant can be a large numerical number if certain sample $\xi$ happen to be an outlier. 
    On the other hand, \citet{pmlr-v202-chen23ar,reisizadeh2023variancereducedclippingnonconvexoptimization} require AN-SGD type of algorithms such as SPIDER \citep{fang2018spidernearoptimalnonconvexoptimization} to have $\Omega(\epsilon^{-2})$ batch size to ensure fast convergence. It essentially makes it a full-batch algorithm and is impractical in practice.

    %In many of the existing studies, such estimation bias is controlled by either adopting the strong assumption that the stochastic gradient approximation error $\|\nabla f_\xi(w) - \nabla F(w)\|$ is bounded \citep{zhang2020gradientclippingacceleratestraining} or 
    %replace strong assumption by constant or affine upper bound of variance \citep{reisizadeh2023variancereducedclippingnonconvexoptimization,pmlr-v202-chen23ar}. However, to guarantee convergence, all of these algorithms require to choose an extremely large batch size in the order of $\Theta(\epsilon^{-2})$. Such a large batch size essentially makes the stochastic algorithm close to the full batch algorithm, which is impractical. 
\end{itemize}
\subsection{Independently-and-Adaptively Normalized SGD}
%\vspace{-2mm}
To overcome the aforementioned limitations, 
we propose the following
independently-and-adaptively normalized stochastic gradient (IAN-SG) estimator 
\begin{align}
    \text{(IAN-SG estimator)} \quad \frac{\nabla f_{\xi}(w)}{\| \nabla f_{\xi'}(w) \|^{\beta}},
\end{align}
where $\xi$ and $\xi'$ are samples drawn {\em independently} from the underlying data distribution. Intuitively, the independence between $\xi$ and $\xi'$ decorrelates the denominator from the numerator, making it an unbiased stochastic gradient estimator (up to a scaling factor). Specifically, we have that
\begin{align}
    \mathbb{E}_{\xi, \xi'} \bigg[ \frac{\nabla f_{\xi}(w)}{ \| \nabla f_{\xi'}(w) \|^{\beta}} \bigg] =  \mathbb{E}_{\xi'} \bigg[ \frac{\mathbb{E}_{\xi} \big [\nabla f_{\xi}(w)\big]}{ \| \nabla f_{\xi'}(w)\|^{\beta}} \bigg] \propto  \nabla F(w).
\end{align}
Moreover, as we show later under mild assumptions, the scaling factor $\mathbb{E} \big[ \| \nabla f_{\xi'}(w)\|^{-\beta}\big]$ can be roughly bounded by the full gradient norm and hence resembling the full-batch AN-GD update.
Based on this idea, we formally propose the following independently-and-adaptively normalized SGD (IAN-SGD) algorithm (see \eqref{update rule}), where $A$, $\Gamma$,$\delta$ are positive constants, $\nabla f_{\xi}(w_t)$, $\nabla f_{\xi'}(w_t)$ corresponds to the stochastic gradient associated with two independent samples $\xi$, $\xi'$ sampled at time $t$.
\begin{align}
   \text{(IAN-SGD):} ~~ w_{t+1} = w_t - \gamma \frac{\nabla f_{\xi}(w_t)}{h_t^{\beta}}, ~~\nonumber\\
   \text{where} ~~ h_t = \max \Big\{1, \Gamma\cdot\Big({A\big \|\nabla f_{\xi'}(w_t)\big \|+ \delta}\Big) \Big \}\label{update rule}.
\end{align}
The above IAN-SGD algorithm adopts a clipping strategy for the normalization term $h_t$. Intuitively, when gradient converges in the later optimization phase, the term $h_t=1$ and IAN-SGD reduces to SGD. This is consistent with the fact that the generalized-smooth condition reduces to the standard smooth condition when gradient approaches to 0. Moreover, imposing a constant lower bound $\delta$ on $h_t$ in \eqref{update rule} helps develop the theoretical convergence analysis and avoid numerical instability in practice. As we show in the ablation study presented in the appendix, the convergence speed of IAN-SGD is insensitive to the choice of $\delta$.
We also note that IAN-SGD requires querying two batches of samples in every iteration. However, we show later in the Theorem \ref{Clip SGD converge} and experiments that these batch sizes can be chosen at a constant level in both theory and practice.


\subsection{Convergence Analysis of IAN-SGD}
We adopt the following bias assumption on the stochastic gradient $\nabla f_{\xi}(w)$.
\begin{assumption}[Unbiased stochastic gradient]\label{unbiasness}
The stochastic gradient $\nabla f_{\xi}(w)$ is unbiased, i.e.,\\ $\mathbb{E}_{\xi\sim \mathbb{P}}\big [\nabla f_{\xi}(w)\big] = \nabla F(w)$ for all $w\in\mathbf{R}^d$.
\end{assumption}
\begin{assumption}[Affine approximation Error]\label{assum6}
    There exists  $0\leq\tau_1<1, \tau_2>0$ such that for any $w\in \mathbf{R}^d$,
    \begin{align}
    \big \|\nabla f_{\xi}(w)-\nabla F(w)\| \leq \tau_1 \big \| \nabla F(w)\big \| + \tau_2 ~\text{ a.s. }~ \forall \xi \sim \mathbb{P}.
    \label{affinebound}
    \end{align}
\end{assumption}
%% Add IAN-SGD convergence under expected noise assumption with large batch-size
Assumption~\ref{assum6} allows the approximation error of stochastic gradient scaling with the full gradient norm on a per-sample basis, while only requiring the error to remain bounded at stationary points. Compared to other noise assumptions, Assumption \ref{assum6} is much weaker than the almost sure error bound (i.e., $\tau_1 = 0$) adopted in \citet{zhang2020gradientclippingacceleratestraining,genralclip,mingruiclip} but stronger than expected error bound (i.e., bounded variance assumption $\mathbb{E}_{\xi}\|\nabla f(w)-\nabla F(w) \|^2 \leq \tau_2^2 $) adopted in \citet{usitchclip, reisizadeh2023variancereducedclippingnonconvexoptimization}. Our motivation for studying the convergence of IAN-SGD under this noise assumption is to strike a balance between achieving standard convergence rates (i.e., $\mathcal{O}(\epsilon^{-4})$) under relaxed noise conditions, while keeping the batch size at the $\Omega(1)$-level. As later we show in Theorem~\ref{Clip SGD converge}, IAN-SGD achieves $\mathcal{O}(\epsilon^{-4})$ sample complexity under Assumption~\ref{assum6}. 
\begin{remark}
    Existing work \citep{usitchclip} has shown that when adopting $\Omega(1)$ sample per iteration, clipped SGD attains $\mathcal{O}(\epsilon^{-5})$ sample complexities under generalized-smooth and bounded variance assumptions.
    In order to achieve the standard $\mathcal{O}(\epsilon^{-4})$ sample complexity under such assumption, \cite{reisizadeh2023variancereducedclippingnonconvexoptimization} demonstrated that clipped SGD requires $\Omega(\epsilon^{-2})$ samples per iteration. Through this way, the constant arising from the variance of the approximation error can be controlled at the $\mathcal{O}(\epsilon^2)$-level. This, in turn, relaxes the burden on iteration complexity, requiring only $\mathcal{O}(\epsilon^{-2})$ iterations for averaging during convergence proof. We leave technical discussions on optimizing the batch size while maintaining $\mathcal{O}(\epsilon^{-4})$ total sample complexity of IAN-SGD for future work.
\end{remark}
Based on these assumptions, we first obtain the following upper bounds regarding $\|\nabla F(w_t) \|$, $\|\nabla f_{\xi}(w_t) \|$. See Proof \ref{Proof of Lemma 2} in Appendix \ref{Appendix H}.
\begin{restatable}{lemma}{mainlemmatwo}
\label{stochastic gradient bound}
    
    Let Assumptions \ref{unbiasness} and \ref{assum6} hold. The gradient $\|\nabla F(w_t) \|$ and stochastic gradient $\nabla f_{\xi}(w_t)$ satisfy 
    \begin{align}
        &\|\nabla F(w_t) \|\leq \frac{1}{1-\tau_1}\|\nabla f_{\xi}(w_t) \|+\frac{\tau_2}{1-\tau_1},\text{and}\nonumber\\
        &\|\nabla f_{\xi}(w_t)\|\leq (1+\tau_1)\|\nabla F(w_t) \|+\tau_2.
    \end{align}
        %When $\big \|\nabla F(w) \big \| \geq \frac{\tau_2}{\tau_1}$,
        %\begin{align}
        %    \big \|\nabla f_{\xi}(w) \big \| \ge \frac{1}{2}\big \|\nabla F(w) \big \| - \frac{\tau_2}{2\tau_1}.
            %\leq  2\big \|\nabla F(w) \big \|. 
        %    \label{stochastic bound 1}
        %\end{align}
        %\item When $\big \|\nabla F(w) \big \| \leq \frac{\tau_2}{\tau_1}$, we have
        %\begin{align}
                %\big \| \nabla F(w)\big \|  \leq 2\big \| \nabla f_{\xi'}(w) \big \| + \frac{\tau_2}{\tau_1} 
                %\leq \frac{2\tau_2}{\tau_1}
                %\label{stochastic bound 2}
        %\end{align}
\end{restatable}
%\begin{remark}
%    For universality, Lemma \ref{stochastic gradient bound} only considers the case when $\tau_1, \tau_2>0$. The key insight for introducing Lemma \ref{stochastic gradient bound} is to construct a lower bound for $(2\|\nabla f_{\xi'}(w_t)\|+\delta)$
%    utilizing $\|\nabla F(w) \|$. When $\tau_1=0$, this goal can also be achieved by considering $B=4$ and $\delta=\frac{\tau_2}{2}$, which reduces to similar algorithm settings introduced in \cite{zhang2020gradientclippingacceleratestraining}.
%\end{remark}
Lemma \ref{stochastic gradient bound} indicates that the full gradient norm can be upper bounded by
the stochastic gradient plus a constant. This result suggests that for $h_t$ in \eqref{update rule}, one can choose $A=\frac{1}{\tau_1-1}$ and $\delta = \frac{\tau_2}{1-\tau_1}$, and it is very useful in our convergence analysis to effectively bound the stochastic gradient norm in $h_t$. We obtain the following convergence result of IAN-SGD. See Proof \ref{Proof of IAN-SGD} in Appendix \ref{Appendix G}

\begin{restatable}[Convergence of IAN-SGD]{theorem}{mainthmtwo}
\label{Clip SGD converge}
    Let Assumptions \ref{assum1}, \ref{unbiasness} and \ref{assum6} hold. For the IAN-SGD algorithm, choose learning rate 
     $\gamma = \min \{ \frac{1}{4L_0(2\tau_1^2+1)},\frac{1}{4L_1(2\tau_1^2+1)}, \frac{1}{\sqrt{T}}, \frac{1}{8L_1(2\tau_1^2+1)(2\tau_2/(1-\tau_1))^{\beta}}\}$, and $A=\frac{1}{1-\tau_1}$ $\delta = \frac{\tau_2}{1-\tau_1}$, $\Gamma=(4L_1\gamma(2\tau_1^2+1))^{\frac{1}{\beta}}$. Denote $\Lambda := F(w_0)-F^* + \frac{1}{2}(L_0+L_1)(1+ 4\tau_2^2)^2$. Then, {with probability at least $\frac{1}{2}$,} IAN-SGD produces a sequence satisfying $\min_{t \le T}\|\nabla F(w_t) \|\leq \epsilon$ if the total number of iteration $T$ satisfies  
     \begin{align}
                T \geq \Lambda \max \Big\{\frac{256\Lambda}{\epsilon^4}, \frac{64L_1(2\tau_1^2+1)(2+2\tau_1)^{\beta}}{(1-\tau_1)^{\beta}\epsilon^{2-\beta}}, (2\tau_1^2+1)\cdot\frac{64(L_0+L_1)+128L_1(2\tau_2/(1-\tau_1))^{\beta}}{\epsilon^2} \Big\}.
          \label{stochastic iteration complexity}
    \end{align}
\end{restatable}
\begin{remark}
Existing work \citep{zhang2020gradientclippingacceleratestraining} studied the iteration complexity of clipped SGD (i.e., $\beta=1$ and $\xi' = \xi$, without independence sampling) in the special bounded noise setting $\tau_1 = 0$, and obtained the complexity 
\begin{align}
    T\geq \Lambda \max\Big\{ \frac{4 \Lambda}{\epsilon^4},\frac{128 L_1}{\epsilon}, \frac{80 L_0 + 512 L_1 \tau_2}{\epsilon^2} \Big\},
\end{align}
where $\Lambda  = \big( F(x_0) - F^* + (5L_0 + 2L_1 \tau_2) \tau_2^2 + 9 \tau_2{L_0^2}/{L_1} \big).
$
As a comparison, the above Theorem \ref{Clip SGD converge} establishes a comparable complexity result under the more general noise model ($\tau_1 \geq 0$) in Assumption \ref{assum6}, by leveraging the adaptive gradient normalization and independent sampling of IAN-SGD. In particular, in the special setting of $\tau_1=0, \beta=1$, these two complexity results match up to absolute constant coefficients. 

Another existing work \citet{genralclip} analyzed a generalized clipping framework that combines momentum acceleration and gradient clipping, namely
$
w_{t+1} = w_t - \big( \nu \min\big\{c, \tfrac{\gamma}{\|m_t\|} \big\} m_t + (1 - \nu) \min\big\{c, \tfrac{\gamma}{\|\nabla f_{\xi}(w_t)\|} \big\} \nabla f_{\xi}(w_t) \big),
$
under a bounded noise setting (i.e., $\tau_1 = 0$). They establish a complexity of $\mathcal{O}\big((F(w_0) - F^*) \tau_2^2 L_0 / \epsilon^4\big)$ when $\gamma / c = \Theta(\tau_2)$, where $m_t$ denotes the momentum term and $c$ is the clipping threshold. This framework enables them to construct a recursion based directly on $\|\nabla F(w_t)\|$, whereas our analysis is built around recursive bounds involving $\|\nabla F(w_t)\|^2$ or $\|\nabla F(w_t)\|^{2 - \beta}$ at each step. Consequently, while both approaches yield the same order of convergence in terms of $\epsilon$, the complexity expressions slightly differ in the numerators due to different recursions.


%\begin{align}
%m_{t+1} &= (1 - \beta) m_t + \beta \nabla f_{\xi}(w_t)\nonumber\\ 
%w_{t+1} &= w_t - ( \nu \min\{c, \frac{\gamma}{\|m_t\|} \} m_t + (1 - \nu) \min\{c, \frac{\gamma}{\|\nabla f_{\xi}(w_t)\|} \} \nabla f_{\xi}(w_t) ).\nonumber
%\end{align}
%And they obtain $T=\mathcal{O}(\frac{(F(w_0)-F^*)\tau_2^2L_0}{\epsilon^{4}})$ under special case of Assumption \ref{assum6} $\tau_1=0$. The denominator $\tau_2^2$ can be large


    %When $\tau_1=0, \beta=1$, Theorem \ref{Clip SGD converge} implies by choosing $\gamma = \min\{\frac{1}{2L_0}, \frac{1}{2L_1}, \frac{1}{\sqrt{T}}, \frac{1}{8\tau_2} \}$, $A=1$, $\delta=\tau_2$ and denote $\Gamma = (2L_1\gamma)^{\frac{1}{\beta}}$, and $\Lambda = F(w_0)-F^*+\frac{1}{2}(L_0+L_1)(1+4\tau_2^2)$, then the time complexity of producing $\min_{t\leq T}\|\nabla F(w_t) \|\leq \epsilon$ with probability at least $\frac{1}{2}$ becomes
    % \begin{align}
    %     T\geq \Lambda \max\Big\{\frac{256 \Lambda}{\epsilon^{4}}, \frac{128L_1}{\epsilon}, \frac{32L_0+32L_1(1+4\tau_2)}{\epsilon^2}
    %     \Big\}.\nonumber
    % \end{align}
    % Under this scenario, the algorithm hyper-parameter setting and obtained $T$ are same as results in \citet{zhang2020gradientclippingacceleratestraining} up to differences in constant coefficients. As a contrast, Theorem \ref{Clip SGD converge} provides general settings of algorithms' hyper-parameter and characterizes time complexity under the existence of affine bounded noise.

    %When $\beta=0$, Theorem \ref{Clip SGD converge}'s iteration complexity $T$ reduces to 
    %\begin{align}
    %    T\geq \Lambda \max \Big\{ \frac{256\Lambda}{\epsilon^4}, \frac{(2\tau_1^2+1)(32L_0+96L_1)}{\epsilon^2}\Big\}.
    %\end{align}
    %This recovers the standard SGD $\mathcal{O}(\epsilon^{-4})$ convergence rate under $L$-smooth setting. 
\end{remark}
%The choices of $B, B' = \mathcal{O}(\tau_1^2)$ are mainly to simplify the symbolic operation during the proof. 
%By deploying normalizing during data pre-processing, the value of $\tau_1$ can be approximately controlled as $\mathcal{O}(1)$ in practice.
Compared to the existing studies \citep{zhang2020gradientclippingacceleratestraining,zhangLsmoothclip, reisizadeh2023variancereducedclippingnonconvexoptimization,genralclip, usitchclip} on gradient clipping algorithms, Theorem \ref{Clip SGD converge} establishes convergence guarantees without the bounded approximation error assumption and the use of extremely large batch sizes that depend on the target error $\epsilon$.
% in neither requires 
% using extremely large batch sizes dependent on target error $\epsilon$ nor depending on the bounded approximation error assumption. 
%Theorem \ref{Clip SGD converge} indicates that IAN-SGD achieves a sample complexity in the order of $\mathcal{O}(\epsilon^{-4})$ with batch sizes independent against target error.  
Through numerical experiments in Section \ref{sec: exp} and ablation study in \ref{sec: effects of Batch Size} later, we show that it suffices to query a small number of independent samples for IAN-SGD in practice.
\subsubsection{Proof outline and novelty} 
The independent sampling strategy adopted by IAN-SGD naturally decouples stochastic gradient from gradient norm normalization, making it easier to achieve the desired optimization progress in generalized-smooth optimization under relaxed conditions.
%is very useful to  the randomness induced from stochastic mini-batch samples $B$ and $B'$ hierarchically. 
By the descent lemma, we have that
\begin{align}
    &\mathbb{E}_{\xi,\xi'}\big [F(w_{t+1}) -F(w_t)|w_t\big]
    \nonumber\\
    \overset{(i)}{\leq}&\mathbb{E}_{\xi'}\Big[-\frac{ \gamma \big \| \nabla F(w_t)\big \|^2}{h_t^{\beta}}+\frac{\gamma^2(L_0
    + L_1\big \|\nabla F(w_t) \big \|^\alpha)}{2h_t^{2\beta}} \mathbb{E}_{\xi}\big [\big \|\nabla f_{\xi}(w_t)\big \|^2 \big]|w_t\Big] \nonumber \\
    \overset{(ii)}{\leq}& \mathbb{E}_{\xi'}\Big[ \Big(\frac{\gamma}{h_t^{\beta}}\big(-1+\gamma(2\tau_1^2+1) \frac{L_0+L_1 \big \|\nabla F(w_t) \big \|^{\alpha}}{h_t^{\beta}}\big)\Big) \|\nabla F(w_t) \big \|^2 +\gamma^2\frac{2\tau_2^2(L_0+L_1 \big \|\nabla F(w_t) \big \|^{\alpha})}{h_t^{2\beta}}|w_t\Big],
    \label{stochastic descent lemma main}
\end{align}
where (i) follows from the descent lemma (conditioned on $w_t$) and the independence between $\xi$ and $\xi'$; (ii) leverages Assumption \ref{assum6} to bound the second moment $\mathbb{E}_{\xi}[\|\nabla f_{\xi}(w_t) \|^2 ]$ by 
$(4\tau_1^2+2) \|\nabla F(w_t) \|^2+ 4\tau_2^2$.
Then, for the first term in \eqref{stochastic descent lemma main}, we leverage the clipping structure of $h_t$ to bound
the coefficient $\gamma(2\tau_1^2+1)(L_0+L_1 \|\nabla F(w) \|^{\alpha})/h_t^{\beta}$ by $\frac{1}{2}$. For the second term in \eqref{stochastic descent lemma main}, we again leverage the clipping structure of $h_t$ and consider two complementary cases: when $\|\nabla F(w_t) \|\leq \sqrt{1+4\tau_2^2/(2\tau_1^2+1})$, this term can be upper bounded by $\frac{1}{2}\gamma^2(L_0+L_1)(1+4\tau_2^2)^2$; when $\|\nabla F(w_t) \|> \sqrt{1+4\tau_2^2/(2\tau_1^2+1)}$, this term can be upper bounded by $\frac{\gamma}{4h_t^{\beta}}\|\nabla F(w_t) \|^2$. Summing them up gives the desired bound. We refer to Lemma \ref{upper bound for term 1} in the appendix for more details.
Substituting these bounds into \eqref{stochastic descent lemma main} and rearranging the terms yields that
\begin{align}
   \mathbb{E}_{\xi'}\big[\frac{\gamma}{4h_t^{\beta}} \big \|\nabla F(w_t) \big \|^2\big]  \leq& \mathbb{E}_{\xi,\xi'}\big[F(w_t) - F(w_{t+1})\big] + \frac{1}{2}(L_0+L_1 )\gamma^2(1+4\tau_2^2)^2.
    \nonumber 
\end{align}
Furthermore, by leveraging the clipping structure of $h_t^{\beta}$ and Assumption \ref{assum6}, the left hand side of above inequality can be lower bounded as 
$\frac{\gamma\| \nabla F(w_t) \|^2}{h_t^{\beta}} \geq \min \{\gamma \|\nabla F(w_t) \|^2, \frac{\|\nabla F(w_t) \|^{2-\beta}}{4((2\tau_2+2)/1-\tau_1)^{\beta}L_1(2\tau_1^2+1)}\}.$
%Through this way, the LHS and RHS of descent lemma don't contain randomness from $B$ and $B'$ anymore. 
Finally, telescoping the above inequalities over $t$ and taking expectation leads to the desired bound in \eqref{stochastic iteration complexity}. 

As a comparison, in the previous works on clipped SGD \citep{zhang2020gradientclippingacceleratestraining,genralclip}, their stochastic gradient and normalization term $h_t$ depend on the same mini-batch of samples, and therefore cannot be treated separately in the analysis. For example, their analysis proposed the following decomposition.
%$\mathbb{E}_{\xi}\big[\frac{\|\nabla f_{\xi}(w_t) \|^2}{h_t^{2\beta}}\big ]$ 
%and have to introduce additional terms to separate $\|\nabla f_{\xi}(w_t) \|^2$ and $h_t^{2\beta}$, which becomes
\begin{align}
    \mathbb{E}_{\xi} \Big[\frac{\|\nabla f_{\xi}(w_t)\|^2 }{h_t^{2\beta}}\Big] =\mathbb{E}_{\xi}\Big[\Big(\| \nabla F(w_t)\|^2+\|\nabla f_{\xi}(w_t)- \nabla F(w_t) \|^2 +2\langle\nabla F(w_t), \nabla f_{\xi}(w_t)-\nabla F(w_t) \rangle\Big)/{h_t^{2\beta}}\Big]. \nonumber 
\end{align}
Hence their analysis needs to assume a constant upper bound for the approximation error $\| \nabla f_{\xi}(w_t)-\nabla F(w_t)\| $ in order to obtain a comparable bound to ours.



\subsubsection{IAN-SGD under generalized P{\L} condition}
Under Assumption \ref{assum2}, we obtain the following recursion for IAN-SGD. See Appendix \ref{Appendix J} for Proof details.
\begin{restatable}[Descent inequality of IAN-SGD under generalized P{\L} condition]{lemma}{mainlemmathree}\label{stochastic PL descent}
        Let Assumptions \ref{assum1}, \ref{assum2}, \ref{unbiasness} and \ref{assum6} hold. For the IAN-SGD algorithm, choose target error $0<\epsilon \leq \min\{1,1/2\mu\}$ and learning rate 
     $\gamma =(2\mu\epsilon)^{\frac{4-2\beta}{\rho}}\cdot \min \{ \frac{1}{4L_0(2\tau_1^2+1)},\frac{1}{4L_1((2\tau_1+2)/(1-\tau_1))(2\tau_1^2+1)}, \frac{1}{8L_1(2\tau_1^2+1)(2\tau_2/(1-\tau_1))^{\beta}},\frac{L_1(2\tau_1^2+1)}{16((L_0+L_1)(1+4\tau_2^2)+L_1(\tau_1^2+1/2))^2}\}$, and $A=\frac{1}{1-\tau_1}$ $\delta = \frac{\tau_2}{1-\tau_1}$, $\Gamma=(4L_1\gamma(2\tau_1^2+1))^{\frac{1}{\beta}}$. Then, we have the following descent inequality 
     \begin{align}
             \mathbb{E}_{w_{t+1}}[\Delta_{t+1}]\leq \mathbb{E}_{w_t}[\Delta_t - \frac{\gamma^{3/2}(L_1(2\tau_1^2+1))^{1/2}(2\mu\Delta_t)^{(2-\beta)/\rho}}{4}+ \frac{\gamma^{3/2}(L_1(2\tau_1^2+1))^{1/2}(2\mu\epsilon)^{(2-\beta)/\rho}}{8}].
             \label{eq: IAN-SGD PL recursion}
     \end{align}
\end{restatable}
We notice that \eqref{eq: IAN-SGD PL recursion} is almost the same as the recursion of AN-GD obtained in Appendix \ref{Appendix D}. The main difference between \eqref{eq: IAN-SGD PL recursion} and \eqref{descentPL} is the usage of $\gamma^{\frac{3}{2}}$ instead of $\gamma$ to characterize the recursion. As a result, the convergence rate of IAN-SGD is expected to be much slower than AN-GD, i.e., IAN-SGD converge with rate $\tilde{\mathcal{O}}({\epsilon}^{-\frac{3\beta}{\rho}})$ when $\rho+\beta=2$; IAN-SGD converge with rate $\mathcal{O}({\epsilon}^{-\frac{3(2-\rho)}{\rho}})$ when $\rho+\beta<2$. We omit the discussions of proof as its convergence analysis exactly follows the proof presented in Appendix \ref{Appendix E}.

\section{Experiments}\label{sec: exp}
We conduct numerical experiments to compare IAN-SGD with other state-of-the-art stochastic algorithms, including the standard SGD \citep{ghadimi2013stochastic}, normalized SGD (NSGD), clipped SGD \citep{zhang2020gradientclippingacceleratestraining}, SPIDER \citep{fang2018spidernearoptimalnonconvexoptimization}, normalized SGD with momentum \citep{NSGDm}(NSGDm) etc. Compared with previous works that study clipped SGD \citep{zhang2020gradientclippingacceleratestraining,genralclip,zhangLsmoothclip,usitchclip,reisizadeh2023variancereducedclippingnonconvexoptimization}, the main change of IAN-SGD lies in its usage of independent sampling and adaptive normalization. To empirically validate the effect of adaptive normalization in accordance with Theorem~\ref{beta GD convergence}, our experiments focus on nonconvex phase retrieval, distributionally robust optimization (DRO), and deep neural network training. Prior works \citep{PhaseretrievalPL,jin2021nonconvex,genPL1,genPL2} have shown that phase retrieval and deep neural networks satisfy generalized P{\L} conditions under mild assumptions. Accordingly, we evaluate the convergence of all algorithms in terms of total sample complexity (for phase retrieval and DRO) and number of epochs (for deep neural networks). Furthermore, to align the DRO formulation with our theoretical assumptions, we focus on regression tasks and adopt the $\chi^{2}$‑divergence to model distributional uncertainty. All experiments were conducted on a PC computer equipped with a $24$‑core CPU, $32$GB of RAM, and a single NVIDIA RTX$4090$ GPU, running Python 3.8.


\subsection{Nonconvex Phase Retrieval}

The phase retrieval problem arises in optics, signal processing, and quantum mechanics \citep{PhaseRetrieval1}. The goal is to recover a signal from measurements where only the intensity is known, and the phase information is missing or difficult to measure. 
%For example, in \citep{PhaseRetrieval1}, phase retrieval is utilized to reconstruct the structure of molecules from the records of intensity.
Specially, denote the underlying object as $w\in \mathbf{R}^d$. Suppose we take $m$ intensity measurements $y_r = |a_r^Tw |^2 + n_r$ for $r=1\cdots m$, where $a_r$ denotes the measurement vector and $n_r$ is the additive noise. We aim to reconstruct $w$ by solving the following regression problem, 
\begin{align}
    \min_{w\in \mathbf{R}^d}f(w)=\frac{1}{2m}\sum_{r=1}^{m} \big(y_r-\big|a_r^Tw \big|^2 \big)^2.
    \label{noncvxphaseretrieval}
\end{align}
Such nonconvex function satisfies generalized-smooth with parameter $\alpha = \frac{2}{3}$. 
In this experiment, we generate the initialization $w_0\sim \mathcal{N}(1,6)$ and the underlying true object $w^*\sim \mathcal{N}(0,0.5)$ with dimension $d=100$. We take 
$m=3$k measurements with {$a_r\sim \mathcal{N}(0,0.5)$ } and $n_r\sim \mathcal{N}(0,16)$. 
We implement all of the stochastic algorithms in original form as described in previous literature. We set batch size $|B|=64$, and for IAN-SGD, we choose a small independent batch size $|B'|=4$. We fine-tuned learning rate for all algorithms, i.e., $\gamma = 5e^{-5}$ for SGD, $\gamma = 0.2$ for NSGD and NSGD with momentum, $\gamma = 0.6$ for clipped SGD, and $\gamma=0.25$ for SPIDER and IAN-SGD. We set the maximal gradient clipping constant as $20$, $\delta = 1e^{-3}$ for both clipped SGD and IAN-SGD. And we set normalization parameter $\beta=\frac{2}{3}$. Figure \ref{PhaseRetrieval+DRO} (left) shows the comparison of objective value versus sample complexity. It can be observed that IAN-SGD consistently converges faster than other baseline algorithms. 
\begin{figure}[htbp]
    \centering
    \begin{subfigure}[b]{0.46\textwidth}
        \centering
        \includegraphics[width=1.0\linewidth]{Fig/PhaseRetrieval.png}
    \end{subfigure}
    \hfill
    \begin{subfigure}[b]{0.46\textwidth}
        \centering
        \includegraphics[width=1.0\linewidth]{Fig/DRO.png}
    \end{subfigure}
    \caption{Experimental Results on Phase Retrieval and DRO}
    \label{PhaseRetrieval+DRO}
\end{figure}
\subsection{Distributionally-Robust Optimization}
Distributionally-robust optimization (DRO) is a popular approach to enhance robustness against data distribution shift. We consider the regularized DRO problem 
$
    \min_{w\in \mathbf{R}^d} f(w)=\sup_{\mathbb{Q}} \big\{\mathbb{E}_{\xi\sim \mathbb{Q}}[\ell_{\xi}(w)]-\lambda \phi(\mathbb{P};\mathbb{Q})\big \},
$
where $\mathbb{Q}$, $\mathbb{P}$ represents the underlying distribution and the nominal distribution respectively. $\lambda$ denotes a regularization hyper-parameter and $\phi$ denotes $\phi$-divergence metric. Under mild technical assumptions, \citet{jin2021nonconvex} showed that such a problem has the following equivalent dual formulation
\begin{align}
    \min_{w\in \mathbf{R}^d,\eta\in \mathbf{R}} L(w,\eta) = \lambda \mathbb{E}_{\xi \sim P} \phi^*\Big(\frac{\ell_{\xi}(w)-\eta}{\lambda}\Big)+\eta,
    \label{dualnoncvxDRO}
\end{align}
where $\phi^*$ denotes the convex conjugate function of $\phi$. In particular, such dual objective function satisfies the generalized-smooth condition \ref{eq: gs} with parameter $\alpha = 1$ \citep{jin2021nonconvex,pmlr-v202-chen23ar}. In this experiment, we use the Life Expectancy data {\citep{life_expectancy_data}} for regression task. For objective \eqref{dualnoncvxDRO}, we set $\lambda = 0.01$, select $\phi^*(t)= \frac{1}{4}(t+2)_+^2-1$, which is the convex conjugate of $\chi^2$-divergence. We adopt the regularized loss $\ell_{\xi}({w})=\frac{1}{2}(y_{\xi}-{x}_{\xi}^{\top}{w})^2+0.1 \sum_{j=1}^{34} \ln (1+|w^{(j)}|)$, where $x_{\xi},y_{\xi}$ represent input and output data, respectively. And we optimize the primal and dual variable of $L(w,\eta)$ in parallel following conclusion drawn from \citet{jin2021nonconvex}.

For the stochastic momentum moving average parameter used for acceleration methods, we set it as $0.1$ and $0.25$ for NSGD with momentum and SPIDER respectively. For stochastic algorithms without usage of multiple mini-batches, i.e., SGD, NSGD, NSGD with momentum and clipped SGD, we set their batch sizes as $|B|=128$. For SPIDER, we set $|B|=128$ and $|B'|=2313$, where the algorithm will conduct a full-gradient computation after every $15$ iterations. For IAN-SGD, we set the batch size for two batch samples as $|B| = 128$ and $|B'| = 8$. We fine-tuned the learning rate for each algorithm, i.e., $\gamma = 4e^{-5}$ for SGD, $\gamma = 5e^{-3}$ for NSGD, NSGD with momentum and SPIDER, $\gamma = 0.11$ for clipped SGD and IAN-SGD. We set the $\delta = 1e^{-1}$, maximal gradient clipping constant as $30$, $25$ for clipped SGD and IAN-SGD respectively. And we set normalization parameter $\beta=\frac{2}{3}$. Figure \ref{PhaseRetrieval+DRO} (right) shows the comparison of objective value versus sample complexity. It can be observed that objective value optimized by IAN-SGD consistently converges faster than other baselines algorithms.
\subsection{Deep Neural Networks}

\begin{figure}[htbp]
\centering
\begin{subfigure}{0.45\textwidth}
  \includegraphics[width=1\linewidth]{Fig/ResNet18.png}
  %\caption{image3}
%  \label{fig:3}
\end{subfigure}
% TBD after DRO is finished
%\medskip
\hfill
\begin{subfigure}{0.45\textwidth}
  \includegraphics[width=1\linewidth]{Fig/ResNet50.png}
  %\caption{image3}
%  \label{fig:3}
\end{subfigure}
\caption{Experimental Result on training ResNet18, ResNet50.}
\label{AddDNN1}
\vspace{-2mm}
\end{figure}
According to \citet{zhang2020gradientclippingacceleratestraining}, generalized-smooth condition \ref{eq: gs} has been observed to hold in deep neural networks.
To further demonstrate the effectiveness of IAN-SGD algorithm, we conduct experiments for training deep neural networks. Specially, we train ResNet18 and ResNet50 \citep{resnet} on CIFAR10 Dataset \citep{CIFAR10} from scratch. We resize images as $32\times 32$ and normalize images with standard derivation equals to $0.5$ on each dimension.
At the beginning of each algorithm, we fix the random seed and initialize model parameters using the Kaiming initialization \citep{kaiminginitialization}. We compare our algorithm with several baseline methods, including SGD \citep{SGD}, Adam \citep{kingma2014adam}, Adagrad \citep{Adagrad}, NSGD, NSGD with momentum \citep{NSGDm} and clipped SGD \citep{zhang2020gradientclippingacceleratestraining}. 

For SGD, Adam and Adagrad, we utilize PyTorch built-in optimizer \citep{pytorchbuiltin} to implement training pipelines. We implement training pipelines for NSGD, NSGD with momentum, clipped SGD and IAN-SGD. The normalization constant is computed through all model parameters at each iteration. The detailed algorithm settings are as follows. For batch size, all algorithms use $B=128$, and $B'=32$ for IAN-SGD. For stochastic momentum moving average parameter, we set it as $0.9$, $0.99$ for Adam, and 0.25 for normalized SGD with momentum. For clipping threshold used in clipped SGD and IAN-SGD, we set them as 2 and $\delta = 1e^{-1}$. The normalization power used for IAN-SGD is $\beta = \frac{2}{3}$. We fine-tuned learning rate for all algorithms, i.e., $\gamma = 1e^{-3}$ for SGD, Adam and Adagrad, $\gamma = 1e^{-1}$ for NSGD and NSGD with momentum, $\gamma = 2e^{-1}$ for clipped SGD and IAN-SGD. We trained ResNet18, ResNet50 on CIFAR10 dataset for $60,80$ epochs and plot the training loss in Figure \ref{AddDNN1}. Figure \ref{AddDNN1} (left) shows the training loss of ResNet18, Figure\ref{AddDNN1} (right) shows the training loss of ResNet50.
As we can see from these figures, the (pink) loss curve optimized by IAN-SGD  indicates fast convergence rate comparable with several baselines, including SGD, NSGD NSGDm, clipped SGD, which demonstrate the effectiveness of IAN-SGD framework.

Based on our experiments on non-convex phase retrieval, distributionally robust optimization (DRO), deep neural networks, and an ablation study \ref{sec: effects of beta}, we found that IAN-SGD performs best when the adaptive normalization parameter is set to $\beta = \frac{2}{3}$. When $\beta$ is too small, i.e., $\beta\leq\frac{1}{2}$, IAN-SGD struggles to converge; when $\beta$ is close to 1, its performance resembles that of clipped SGD. To achieve fast convergence, the learning rate for IAN-SGD should lie between $1e^{-1}$ and $2e^{-1}$. Users should carefully fine-tune this parameter, as overly large learning rates may lead to divergence.

Regarding the numerical stabilization term $\delta$, our ablation study (see Appendix \ref{sec: effects of delta}) shows that IAN-SGD is empirically insensitive to its value; small value such as $1e^{-3}$ is sufficient to ensure convergence. Finally, to balance stable performance and fast convergence, ablation results (see Appendix \ref{sec: effects of Batch Size}) suggest that user should set batch size $B\geq 16$ and the independent sample batch size $|B'|$ should be at least one-quarter of batch size $|B|$.


\section{Conclusion}
In this work, we provide theoretical insights on how normalization interplays with function geometry, and their overall effects on convergence. We then propose independent normalized stochastic gradient descent for stochastic setting, achieving same sample complexity under relaxed assumptions. Our results extend the existing boundary of first-order nonconvex optimization and may inspire new
developments in this direction. In the future, it is interesting to explore if the popular acceleration method such as stochastic momentum and variance reduction can be combined with independent sampling and normalization to improve the sample complexity.

\section{Acknowledgements}
%Acknowledgements starts here
We are grateful to the anonymous reviewers and Dr. Ziyi Chen for insightful comments that led to several improvements in the presentation and technical details.
This research was supported in part by the Air Force Research Laboratory Information Directorate (AFRL/RI), through AFRL/RI Information Institute®, contract number FA8750-20-3-1003. The work of Yufeng Yang and Yi Zhou was supported by the National Science Foundation under grants DMS-2134223, ECCS-2237830. The work of Shaofeng Zou was supported in part by National Science Foundation under Grants CCF-2438429 and ECCS-2438392(CAREER).
The work of Yifan Sun was supported in part by the ONR grant W911NF-22-1-0292. 


%\begin{table}[t]
%\caption{Sample table title}
%\label{sample-table}
%\begin{center}
%\begin{tabular}{ll}
%\multicolumn{1}{c}{\bf PART}  &\multicolumn{1}{c}{\bf DESCRIPTION}
%\\ \hline \\
%Dendrite         &Input terminal \\
%Axon             &Output terminal \\
%Soma             &Cell body (contains cell nucleus) \\
%\end{tabular}
%\end{center}
%\end{table}
%\subsubsection*{Broader Impact Statement}
%In this optional section, TMLR encourages authors to discuss possible repercussions of their work,
%notably any potential negative impact that a user of this research should be aware of. 
%Authors should consult the TMLR Ethics Guidelines available on the TMLR website
%for guidance on how to approach this subject.

%\subsubsection*{Author Contributions}
%If you'd like to, you may include a section for author contributions as is done
%in many journals. This is optional and at the discretion of the authors. Only add
%this information once your submission is accepted and deanonymized. 
%\subsubsection*{Acknowledgments}
%Use unnumbered third level headings for the acknowledgments. All
%acknowledgments, including those to funding agencies, go at the end of the paper.
%Only add this information once your submission is accepted and deanonymized. 
\bibliography{main}
\bibliographystyle{tmlr}

\newpage
\appendix
\setcounter{tocdepth}{1}
\renewcommand{\contentsname}{Appendix Table of Contents} % Change ToC title if desired
\tableofcontents
%\renewcommand{\contentsname}{Appendix Table of Contents}
    % This will generate a ToC for the appendix only



\newpage
\section{Ablation Study}\label{sec: ablation study}
In order to have a comprehensive understanding of the performance of IAN-SGD in practical problems, we conduct ablation studies regarding on important components of IAN-SGD, including adaptive normalization $\beta$, batch size $|B'|$ of independent samples and numerical stabilizer $\delta$.

\subsection{\texorpdfstring{Effects of $\beta$}{Effects of beta}}\label{sec: effects of beta}
To justify the advantage of using adaptive normalization, we conduct the following two experiments. 

In the first experiment, we choose a unified normalization parameter $\beta = \frac{2}{3}$ for all the normalized methods, i.e., NSGD, NSGD with momentum, clipped SGD, SPIDER and IAN-SGD. To guarantee convergence, we fine-tuned the learning rate accordingly, i.e., $\gamma = 0.03$ for NSGD and NSGD with momentum, $\gamma = 0.05$ for SPIDER, $\gamma = 0.17$ for both clipped SGD and IAN-SGD. To make a fair comparison, we keep other parameters unchanged. Figure \ref{effects of beta} (left) shows the comparison of objective value versus sample complexity for the Phase Retrieval problem. It can be observed that, by adjusting $\beta = \frac{2}{3}$, the objective value optimized by all algorithms decreases much faster compared with Figure \ref{PhaseRetrieval+DRO}. This indicates that adaptive normalization can accelerate convergence. Moreover, compared with other normalization methods, although IAN-SGD requires additional sampling at each iteration, the training loss still decreases faster than NSGD SGD with momentum and SPIDER. 

Similarly, for DRO, we unify the normalization parameter of all the normalized methods, i.e., NSGD, NSGD with momentum, clipped SGD, SPIDER and IAN-SGD with $\beta = \frac{2}{3}$. To guarantee algorithm convergence, we fine-tuned learning rate accordingly. For NSGD, NSGDm and SPIDER, we keep learning rate unchanged; For clipped SGD, and IAN-SGD, we set $\gamma = 0.08$. To make a fair comparison, we keep other parameters unchanged. Figure \ref{effects of beta} (right) shows the comparison of objective value versus sample complexity. It can be observed that, by setting $\beta = \frac{2}{3}$, the objective value optimized by all normalization methods decreases faster than Figure \ref{PhaseRetrieval+DRO}. This verifies the effectiveness of adaptive normalization.
Although IAN-SGD requires additional sampling, it converges faster than other normalized methods. 

In summary, the results in Figure \ref{effects of beta} indicate that independent sampling and adaptive normalization doesn't increase the sample complexity to find a stationary point, which justifies IAN-SGD framework's effectiveness when dealing with nonconvex geometry characterized by generalized-smooth condition.
\begin{figure}[ht]
    \centering
    \begin{subfigure}[b]{0.45\textwidth}
        \centering
        \includegraphics[width=1.15\linewidth]{Fig/PhaseR_Ablation1.png}
    \end{subfigure}
    \hfill
    \begin{subfigure}[b]{0.45\textwidth}
        \centering
        \includegraphics[width=1.15\linewidth]{Fig/DRO_ablation1.png}
    \end{subfigure}

    \caption{Advantage of using adaptive normalization on normalized first-order algorithms}
    \label{effects of beta}
\end{figure}

In the second experiment, we vary $\beta$ for IAN-SGD and keep other parameters unchanged. Figure \ref{effects of beta ctd} shows the convergence result with different $\beta$ for Phase Retrieval and DRO. It can be observed that decreasing $\beta$  accelerates convergence in general. However, small $\beta$ can make convergence unstable. In the DRO experiment shown in Figure \ref{effects of beta ctd} (right), we observed that when $\beta=\frac{3}{5}$, the objective value curve fluctuates a lot. If $\beta$ is smaller than $\frac{3}{5}$, the algorithm fails to converge. This phenomenon matches the implications of Theorem \ref{beta GD convergence}, where $\beta$ must satisfy $\beta\in [\alpha,1]$ to guarantee convergence.
\begin{figure}[ht]
    \centering
    \begin{subfigure}[b]{0.45\textwidth}
        \centering
        \includegraphics[width=1.15\linewidth]{Fig/beta_variation_PR.png}
    \end{subfigure}
    \hfill
    \begin{subfigure}[b]{0.45\textwidth}
        \centering
        \includegraphics[width=1.15\linewidth]{Fig/beta_variation_DRO.png}
    \end{subfigure}

    \caption{Effects of adaptive normalization on convergence of IAN-SGD}
    \label{effects of beta ctd}
\end{figure}


\subsection{Effects of Independent Samples' Batch Size}\label{sec: effects of Batch Size}

In this section, we study the effects of batch size for independent samples used in IAN-SGD. For Phase retrieval, we keep $|B|=64$ and other parameters same as section \ref{sec: effects of beta}, and we vary independent batch sizes to be $|B'| = \{ 4, 8, 16, 32, 64\} $. 
Similarly, for DRO, we keep $|B|=128$ and other parameters same as section \ref{sec: effects of beta}, and we vary batch size of independent samples $|B'| = \{ 16, 32, 64, 128\}$.

The following Figure \ref{effects of batch size} shows the convergence of IAN-SGD under different batch size choices for Phase Retrieval (left) and DRO (right).
\begin{figure}[ht]
    \centering
    \begin{subfigure}[b]{0.45\textwidth}
        \centering
        \includegraphics[width=1.15\linewidth]{Fig/batch_size_PR.png}
    \end{subfigure}
    \hfill
    \begin{subfigure}[b]{0.45\textwidth}
        \centering
        \includegraphics[width=1.15\linewidth]{Fig/batch_size_DRO.png}
    \end{subfigure}

    \caption{Effects of independent samples' batch size on convergence}
    \label{effects of batch size}
\end{figure}
It can be observed that for Phase retrieval and DRO problem, small batch sizes such as $|B'|=\{ 4, 8\} $ are sufficient to guarantee convergence, whereas large batch sizes of $|B'|$ do not significantly increase the sample complexity, which verifies the effectiveness of independent sampling.

\subsection{\texorpdfstring{Effects of $\delta$}{Effects of delta}}\label{sec: effects of delta}
In Theorem \ref{Clip SGD converge}, we set $\delta = \frac{\tau_2}{1-\tau_1}$ based on Assumption \ref{assum6}, which assumes an affine bound for approximation error $\|\nabla f_{\xi}(w)-\nabla F(w) \|$. This assumption relaxes the strong assumption used in \citet{zhang2020gradientclippingacceleratestraining,mingruiclip,genralclip}, where they assume that the approximation error $\|\nabla f_{\xi}(w)-\nabla F(w) \|$ is upper bounded by a constant. The major weakness is that if certain samples have outlier gradients, such constant upper bound can be very large. Thus, to verify the effectiveness of Assumption \ref{assum6}, we expect convergence of IAN-SGD under a wide range of $\delta$, especially for small $\delta$. 

To study the effect of $\delta$, we keep other parameters the same as mentioned in Section \ref{sec: effects of beta} and only vary $\delta$. For Phase Retrieval, we vary $\delta=\{1e^{-8},1e^{-3}, 1e^{-1}, 1,10 \}$ and Figure \ref{effects of delta} (left) shows the corresponding convergence result. For DRO, we vary $\delta = \{1e^{-8}, 1e^{-3}, 1e^{-1}, 1, 10 \}$ and Figure \ref{effects of delta} (right) shows the convergence result. We observe that IAN-SGD convergence is robust to the choice of $\delta$. But $\delta=1e^{-3}, 1e^{-1}$ demonstrate faster convergence than others for Phase retrieval and DRO experiments. This result indicates that a small $\delta$ is sufficient to guarantee convergence, which verifies the effectiveness of assumption \ref{assum6} and IAN-SGD when dealing with stochastic nonconvex generalized-smooth geometry.
\begin{figure}[ht]
    \centering
    \begin{subfigure}[b]{0.45\textwidth}
        \centering
        \includegraphics[width=1.15\linewidth]{Fig/effects_of_delta_PR.png}
    \end{subfigure}
    \hfill
    \begin{subfigure}[b]{0.45\textwidth}
        \centering
        \includegraphics[width=1.15\linewidth]{Fig/effects_of_delta_DRO.png}
    \end{subfigure}

    \caption{Effects of $\delta$ on convergence}
    \label{effects of delta}
\end{figure}
\subsubsection{Effects of  $\tau_1$, $\tau_2$ on convergence}
In Theorem~\ref{Clip SGD converge}, we set $ A = \frac{1}{1 - \tau_1}$ and $ \delta = \frac{\tau_2}{1 - \tau_1} $, making the sample complexity dependent on both $ \tau_1 $ and $ \tau_2 $. To further investigate the effect of the gradient noise assumption, we perform experiments that allow flexible choices of $\tau_1 $ and $ \tau_2 $. Specifically, we control the stochastic gradient noise by computing gradients over mini-batch samples with $ B = B' = 16 $ at each iteration. We then evaluate IAN-SGD under various noise settings with $ \tau_1 \in \{0, 0.5, 0.8\} $ and $ \tau_2 \in \{10^{-8}, 10^{-3}, 10^{-1}, 1, 10\} $. As a result, the scaling factor becomes $ A \in \{1, 2, 5\} $ and we define $ \delta = A \cdot \{10^{-8}, 10^{-3}, 10^{-1}, 1, 10\}$ accordingly. For phase retrieval, we use a unified learning rate $ \gamma = 1e^{-1} $, and for DRO, we set $ \gamma = 5e^{-3} $. Other parameters, such as the clipping threshold and $ \beta $, remain unchanged. Figure~\ref{PR DRO effects of tau} plots the results, where the first row shows IAN-SGD's convergence for phase retrieval, and the second row shows IAN-SGD's convergence for DRO.

\begin{figure}[ht]
    \centering
    \begin{subfigure}[b]{0.3\textwidth}
        \centering
        \includegraphics[width=1.0\linewidth]{Fig/PR_Abla_tau_3.png}
    \end{subfigure}
    %\hfill
    \begin{subfigure}[b]{0.3\textwidth}
        \centering
        \includegraphics[width=1.0\linewidth]{Fig/PR_Abla_tau_1.png}
    \end{subfigure}
    \begin{subfigure}[b]{0.3\textwidth}
        \centering
        \includegraphics[width=1.0\linewidth]{Fig/PR_Abla_tau_2.png}
    \end{subfigure}
    \hfill
    \begin{subfigure}[b]{0.3\textwidth}
        \centering
        \includegraphics[width=1.0\linewidth]{Fig/DRO_vary_tau_3.png}
    \end{subfigure}
    %\hfill
    \begin{subfigure}[b]{0.3\textwidth}
        \centering
        \includegraphics[width=1.0\linewidth]{Fig/DRO_vary_tau_1.png}
    \end{subfigure}
    \begin{subfigure}[b]{0.3\textwidth}
        \centering
        \includegraphics[width=1.0\linewidth]{Fig/DRO_vary_tau_2.png}
    \end{subfigure}

    \caption{Effects of $\tau_1, \tau_2$ for IAN-SGD convergence}
    \label{PR DRO effects of tau}
\end{figure}
As we can see from these figures, IAN-SGD converges slightly slower as $\tau_1$ increases for both the phase retrieval and DRO problems. This observation aligns with Theorem~\ref{Clip SGD converge}, which shows that when $\tau_1 > 0$, the resulting complexity is slightly worse than in the case $\tau_1 = 0$. Regarding $\tau_2$, although the final convergence intervals remain approximately the same, we find that smaller values lead to less stable convergence for phase retrieval. For
DRO, the variation in $\tau_2$ does not significantly affect convergence when learning rate is unified. This might be because for the DRO experiment, the model initialization and data processing have reduced the noise of the stochastic gradients, enabling convergence with flexible choices of $\tau_2$.

\subsection{Ablation Study on Extreme Small Batch Size $|B|, |B'|$}
In this subsection, we study the effects of batch size in IAN-SGD. Specially, we want to justify if IAN-SGD works with $B,B'=\Omega(1)$ in practice. For phase retrieval, we set $B=B'$ and keep other parameters the same as in section \ref{sec: effects of beta}. For comparison, we vary batch sizes to be $\{ 1, 2, 4, 8, 16, 32\}$. To guarantee algorithm converging to same level of magnitude, we fine-tune learning rates accordingly, i.e., for each batch size, we set $\gamma=\{5e^{-4},5e^{-3}, 2e^{-2},4e^{-2}, 0.10, 0.10\}$, respectively. 
Similarly, for DRO, we set $|B|=|B'|$ and keep other parameters the same as section \ref{sec: effects of beta}. For comparison, we vary batch size to be $ \{ 1,2,4, 8, 16, 32\}$, and we fine-tuned learning rates accordingly, i.e., for each batch size, we set $\gamma=\{5e^{-4}, 1e^{-3},1e^{-3}, 5e^{-3}, 5e^{-3}, 1e^{-2}\}$, respectively. Figure \ref{effects of batch size} plots the convergence of IAN-SGD under different batch size choices for Phase Retrieval (left) and DRO (right). We conclude IAN-SGD converges under constant-level batch size $B, B'=\Omega(1)$.
\begin{figure}[ht]
    \centering
    \begin{subfigure}[b]{0.45\textwidth}
        \centering
        \includegraphics[width=1.15\linewidth]{Fig/ablation_small_batch_PR.png}
    \end{subfigure}
    \hfill
    \begin{subfigure}[b]{0.45\textwidth}
        \centering
        \includegraphics[width=1.15\linewidth]{Fig/ablation_small_batch_DRO.png}
    \end{subfigure}

    \caption{Effects of batch size on convergence for IAN-SGD}
    \label{effects of extreme small batch size}
\end{figure}

Next, we compare IAN-SGD with other baseline methods under small batch size. To make a trade-off between convergence and stability, we set $|B|=2$ for SGD, NSGD with momentum, Clip-SGD and SPIDER; For SPIDER algorithm \citep{fang2018spidernearoptimalnonconvexoptimization}, we set epoch size to be $20$. For phase retrieval, we fine-tuned learning rate for all algorithms, i.e., $\gamma=1e^{-5}$ for SGD, $\gamma = 1e^{-2}$ for NSGD, NSGD with momentum, $\gamma = 5e^{-3}$ for clipped SGD and IAN-SGD, $\gamma = 5e^{-2}$ for SPIDER. Similarly, for DRO, we set $|B|=2$ for SGD, NSGD with momentum, clipped SGD and SPIDER. We fine-tuned learning rate for all algorithms, i.e., $\gamma=1e^{-6}$ for SGD, $\gamma = 1e^{-3}$ for NSGD, NSGD with momentum, clipped SGD and IAN-SGD, $\gamma = 5e^{-3}$ for SPIDER. We keep other parameters the same as in Section \ref{sec: effects of beta}. The first row of Figure \ref{Comparisons of IAN-SGD under extreme small batch size} plots the comparison results of Phase Retrieval (left) and DRO (right). We find in both problems, normalized methods such as NSGD, SPIDER, clipped SGD are slow to converge under small batch size.

To further justify the effectiveness of IAN-SGD under small batch size, we also test IAN-SGD and other baseline methods using small batch size to train ResNet \citep{resnet} over CIFAR10 \citep{CIFAR10}. For all algorithms, we unify the batch size $B=32$. We fine-tuned learning for all algorithms, i.e., $\gamma = 1e^{-5}$ for SGD, $\gamma = 1e^{-4}$ for Adam, $\gamma = 1e^{-3}$ for Adagrad, $\gamma = 5e^{-2}$ for NSGD and NSGD with momentum, $\gamma = 8e-2$ for clipped SGD and IAN-SGD, the second row of Figure \ref{Comparisons of IAN-SGD under extreme small batch size} plots the corresponding results for training ResNet18 (left) and ResNet50 (right). As a result, we found IAN-SGD still outperforms than most baselines. This verifies the effectiveness of using IAN-SGD with small batch size in practice.

\begin{figure}[ht]
    \centering
    \begin{subfigure}[b]{0.45\textwidth}
        \centering
        \includegraphics[width=1.15\linewidth]{Fig/IAN_SGD_vs_baseline_small_batch_PR.png}
    \end{subfigure}
    \hfill
    \begin{subfigure}[b]{0.45\textwidth}
        \centering
        \includegraphics[width=1.15\linewidth]{Fig/IAN-SGD_vs_baseline_small_batch_DRO.png}
    \end{subfigure}
    \hfill
    \begin{subfigure}[b]{0.45\textwidth}
        \centering
        \includegraphics[width=1.15\linewidth]{Fig/ResNet18_small_batch.png}
    \end{subfigure}
    \hfill
    \begin{subfigure}[b]{0.45\textwidth}
        \centering
        \includegraphics[width=1.15\linewidth]{Fig/ResNet50_small_batch.png}
    \end{subfigure}
    \caption{Comparison of IAN-SGD with other baselines under small batch size}
    \label{Comparisons of IAN-SGD under extreme small batch size}
\end{figure}


\newpage
\section{Proof of Descent Lemma \ref{descent lemma} }\label{Appendix B}
\mainlemmaone*
\begin{proof}\label{proof of descent lemma}
    Denote $w_{\theta} = \theta w'+(1-\theta)w$. By the fundamental theorem of calculus we have $f(w')-f(w)= \int_0^1 \langle \nabla f(w_{\theta}), w'-w\rangle \mathrm{d}\theta$, which implies that
    \begin{align}
        &f(w')-f(w)-\langle \nabla f(w), w'-w \rangle = \int_{0}^{1}\langle \nabla f(w_{\theta}), w'-w\rangle \mathrm{d} \theta - \int_{0}^{1} \langle \nabla f(w), w'-w\rangle \mathrm{d} \theta  \nonumber,
    \end{align}
    %Since the integration integrates over $w_{\theta}$, integrating $\langle \nabla f(w), w'-w \rangle$ doesn't affect the result.
    Substituting the generalized-smooth condition in Assumption \ref{assum1} into the above equation, we have
    \begin{align}
         &f(w')-f(w)-\langle \nabla f(w), w'-w \rangle \nonumber \\ 
        =& \int_{0}^{1}\langle \nabla f(w_{\theta}), w'-w\rangle \mathrm{d} \theta - \int_{0}^{1} \langle \nabla f(w), w'-w\rangle \mathrm{d} \theta  \nonumber \\
        =& \int_{0}^{1}\langle \nabla f(w_{\theta})-\nabla f(w), w'-w\rangle \mathrm{d} \theta \nonumber\\
        \overset{(i)}{\leq}& \int_{0}^{1} \big \|\nabla f(w_{\theta})-\nabla f(w) \big \| \big \|w'-w \big \|\mathrm{d}\theta \nonumber \\
        \overset{(ii)}{\leq}& \int_{0}^{1} \theta (L_0+L_1 \big \|\nabla f(w) \big \|^{\alpha})\big \|w'-w \big \|^2 \mathrm{d}\theta \nonumber \\
        =& \frac{1}{2}(L_0+L_1\big \|\nabla f(w) \big \|^{\alpha})\big \|w'-w \big \|^2,
    \end{align}
    where (i) applies Cauchy-schwarz inequality on $\langle \nabla f(w_{\theta})-\nabla f(w), w_{\theta}-w\rangle$, and  (ii) is due to Assumption \ref{assum1}, $\|\nabla f(w_{\theta})-\nabla f(w) \big \|\leq (L_0+L_1 \|\nabla f(w) \|^{\alpha})\|w_{\theta}-w\|=\theta(L_0+L_1 \|\nabla f(w) \|^{\alpha})\|w'-w\|$. Re-organizing the above inequality gives the desired result.
\end{proof}
\section{Proof for Nonconvex Phase Retrieval}\label{Appendix C}
%\subsection{Proof for nonconvex Phase Retrieval}
%\phaseretrievalsmooth*
The proof of generalized-smooth property for nonconvex Phase retrieval problem is similar as proof in \citet{pmlr-v202-chen23ar} with minor changes. We present the proof details here for completeness.  
\begin{proof}\label{proof of Phase Retrieval}
The objective function of phase retrieval problem can be expressed in the following form
\begin{align}
    F(w)=\frac{2}{m}\sum_{\xi \in D}f_{\xi}(w)=\frac{1}{2m}\sum_{\xi \in D}\big(y_{\xi}-\big|a_{\xi}^{\top} w\big|^2\big)^2,
    \label{phase retrieval objective}
\end{align}
where $f_{\xi}(w)=\frac{1}{4}(y_{\xi}-|a_{\xi}^Tw|^2)^2$.

We prove that for every sample $\xi$, $\nabla f_{\xi}(w)=-(y_{\xi}-|a_{\xi}^{\top}w|^2)(a_{\xi}{a_{\xi}}^{\top})w$ satisfies \eqref{symmetric gs} with $\alpha=\frac{2}{3}$. And thus $F(w)$ satisfies finite-sum versions of \eqref{symmetric gs}, namely
\begin{align}
    \frac{2}{m}\sum_{\xi\in D}\big\| f_{\xi}(w)-f_{\xi}(w')\big\|\leq \big \| w-w'\big \|\cdot\Big(L_0+\frac{L_1}{m}\sum_{\zeta\in D}\frac{\|\nabla f_{\xi}(w)\|^{\alpha}+\|\nabla f_{\xi}(w')\|^{\alpha}}{2}\Big)
    \label{symmetric gs stochastic}
\end{align}

We construct a lower bound of $\|\nabla f_{\xi}(w)\|$ given $w$ and $\xi$ as following
\begin{align}
\big\|\nabla f_{\xi}(w)\big\|^{\frac{2}{3}} 
& =\big\|\big(\big|a_{\xi}^{\top} w\big|^2-y_{\xi}\big)\big(a_{\xi} a_{\xi}^{\top}\big) w\big\|^{\frac{2}{3}} \nonumber\\
& \overset{}{=}  \big\|\big(\big|a_{\xi}^{\top} w\big|^2\cdot a_{\xi}^{\top}w-y_{\xi}a_{\xi}^{\top}w\big)a_{\xi}\big\|^{\frac{2}{3}}\nonumber \\
& \overset{(i)}{=} \Big| \big| a_{\xi}^{\top} w\big|^3-y_{\xi}\big|a_{\xi}^{\top} w\big|\Big|^{\frac{2}{3}}\big\|a_{\xi}\big\|^{\frac{2}{3}} \nonumber\\
& \overset{(ii)}{\geq} \big(\big|a_{\xi}^{\top} w\big|^2-\Big|y_{\xi} \big| a_{\xi}^{\top} w \big| \Big|^{\frac{2}{3}}\big)\big\|a_{\xi}\big\|^{\frac{2}{3}} \nonumber\\
& \overset{(iii)}{\geq} \frac{2}{3}\big(\big|a_{\xi}^{\top} w\big|^2-\big|y_{\xi}\big|\big)\big\|a_{\xi}\big\|^{\frac{2}{3}},
\label{PRgradient}
\end{align}
where (i) extracts scalar $\Big| \big| a_{\xi}^{\top} w\big|^3-y_{\xi}\big|a_{\xi}^{\top} w\big|\Big|^{\frac{2}{3}}$ out; (ii) uses inequality $|a-b|^{\frac{2}{3}}\geq |a|^{\frac{2}{3}}-|b|^{\frac{2}{3}}$ holding for any $a,b\in \mathbf{R}$ (sub-additive property of function $f(x)=x^{2/3}$); (iii)  applies young's inequality $|y_{\xi}|^{\frac{2}{3}}|a_{\xi}^{\top}w|^{\frac{2}{3}}\leq \frac{1}{3}\|a_{\xi}^{\top}w\|^2+\frac{2}{3}|y_{\xi}|$.

Then, for any $w, w'\in \mathbf{R}^d$, we have
\begin{align}
& \big\|\nabla f_{\xi}\big(w^{\prime}\big)-\nabla f_{\xi}(w)\big\| \nonumber \\
=&\big\|\big(\big|a_{\xi}^{\top} w^{\prime}\big|^2-y_{\xi}\big)\big(a_{\xi} a_{\xi}^{\top}\big) w^{\prime}-\big(\big|a_{\xi}^{\top} w\big|^2-y_{\xi}\big)\big(a_{\xi} a_{\xi}^{\top}\big) w\big\| \nonumber \\
=& \frac{1}{2}\big\|2\big(\big|a_{\xi}^{\top} w^{\prime}\big|^2-y_{\xi}\big)\big(a_{\xi} a_{\xi}^{\top}\big) w^{\prime}-2\big(\big|a_{\xi}^{\top} w\big|^2-y_{\xi}\big)\big(a_{\xi} a_{\xi}^{\top}\big) w +\big(\big|a_{\xi}^{\top}w'\big|^2-y_{\xi}\big)\big(a_{\xi}a_{\xi}^{\top}\big)w-\big(\big|a_{\xi}^{\top}w\big|^2-y_{\xi}\big)(a_{\xi}a_{\xi}^{\top})w' \nonumber\\
&\quad -\big(\big|a_{\xi}^{\top}w'\big|^2-y_{\xi}\big)\big(a_{\xi}a_{\xi}^{\top}\big)w+\big(\big|a_{\xi}^{\top}w\big|^2-y_{\xi}\big)(a_{\xi}a_{\xi}^{\top})w'
\big\|\nonumber\\
=& \frac{1}{2}\big \|\big(\big|a_{\xi}^{\top} w'\big|^2-y_{\xi}\big)\big(a_{\xi} a_{\xi}^{\top}\big) w'-\big(\big|a_{\xi}^{\top}w'\big|^2-y_{\xi}\big)\big(a_{\xi}a_{\xi}^{\top}\big)w+\big(\big|a_{\xi}^{\top}w\big|^2-y_{\xi}\big)(a_{\xi}a_{\xi}^{\top})w'
-\big(\big|a_{\xi}^{\top} w\big|^2-y_{\xi}\big)\big(a_{\xi} a_{\xi}^{\top}\big) w \nonumber\\
&\quad+\big(\big|a_{\xi}^{\top} w'\big|^2-y_{\xi}\big)\big(a_{\xi} a_{\xi}^{\top}\big) w'+\big(\big|a_{\xi}^{\top}w'\big|^2-y_{\xi}\big)\big(a_{\xi}a_{\xi}^{\top}\big)w -\big(\big|a_{\xi}^{\top}w\big|^2-y_{\xi}\big)(a_{\xi}a_{\xi}^{\top})w'
-\big(\big|a_{\xi}^{\top} w\big|^2-y_{\xi}\big)\big(a_{\xi} a_{\xi}^{\top}\big) w\big \|\nonumber\\
=& \frac{1}{2}\big\|\big(\big|a_{\xi}^{\top} w^{\prime}\big|^2+\big|a_{\xi}^{\top} w\big|^2-2 y_{\xi}\big)\big(a_{\xi} a_{\xi}^{\top}\big)\big(w^{\prime}-w\big) +\big(\big|a_{\xi}^{\top} w^{\prime}\big|^2-\big|a_{\xi}^{\top} w\big|^2\big)\big(a_{\xi} a_{\xi}^{\top}\big)\big(w^{\prime}+w\big)\big\| \nonumber \\
\overset{(i)}{\leq} & \frac{1}{2}\big\|\big(\big|a_{\xi}^{\top} w^{\prime}\big|^2+\big|a_{\xi}^{\top} w\big|^2-2 y_{\xi}\big)\big(a_{\xi} a_{\xi}^{\top}\big)\big(w^{\prime}-w\big) \big \|
+\frac{1}{2}\big \|\big(\big|a_{\xi}^{\top} w^{\prime}\big|^2-\big|a_{\xi}^{\top} w\big|^2\big)\big(a_{\xi} a_{\xi}^{\top}\big)\big(w^{\prime}+w\big)\big\| \nonumber \\
=& \frac{1}{2}\big\|\big(\big|a_{\xi}^{\top} w^{\prime}\big|^2+\big|a_{\xi}^{\top} w\big|^2-2 y_{\xi}\big)\big(a_{\xi} a_{\xi}^{\top}\big)\big(w^{\prime}-w\big) \big \|
+\frac{1}{2}\big \|\big(\big|a_{\xi}^{\top} w^{\prime}\big|^2-\big|a_{\xi}^{\top} w\big|^2\big)a_{\xi}\big(a_{\xi}^{\top}w^{\prime}+a_{\xi}^{\top}w\big)\big\| \nonumber \\
\overset{(ii)}{\leq}&  \frac{1}{2}\big\|\big(\big|a_{\xi}^{\top} w^{\prime}\big|^2+\big|a_{\xi}^{\top} w\big|^2-2 y_{\xi}\big)\big(a_{\xi} a_{\xi}^{\top}\big)\big(w^{\prime}-w\big) \big \|
+\frac{1}{2}\big \|\big(\big|a_{\xi}^{\top} w^{\prime}\big|^2-\big|a_{\xi}^{\top} w\big|^2\big)a_{\xi}\big(\big|a_{\xi}^{\top}w^{\prime}\big|+\big| a_{\xi}^{\top}w\big|\big)\big\| \nonumber \\
\overset{(iii)}{\leq}& \frac{1}{2}\big\|a_{\xi}\big\|^2\big(\big|a_{\xi}^{\top} w^{\prime}\big|^2+\big|a_{\xi}^{\top} w\big|^2+2\big|y_{\xi}\big|\big)\big\|w^{\prime}-w\big\| +\frac{1}{2}\big\|a_{\xi}\big\|^2\big(\big|a_{\xi}^{\top} w^{\prime}\big|+\big|a_{\xi}^{\top} w\big|\big)^2\big\|w^{\prime}-w\big\| \nonumber \\
\overset{(iv)}{\leq}& \frac{1}{2}\big\|w^{\prime}-w\big\|\big\|a_{\xi}\big\|^2\big(3\big|a_{\xi}^{\top} w^{\prime}\big|^2+3\big|a_{\xi}^{\top} w\big|^2+2\big|y_{\xi}\big|\big) \nonumber\\
\leq& \frac{1}{2}\big\|w^{\prime}-w\big\|\big\|a_{\xi}\big\|^{\frac{4}{3}}\big\|a_{\xi}\big\|^{\frac{2}{3}}\cdot\big(3\big|a_{\xi}^{\top} w^{\prime}\big|^2+3\big|a_{\xi}^{\top} w\big|^2-3\big|y_{\xi}\big|-3\big|y_{\xi}\big|+8\big|y_{\xi}\big|\big)\nonumber\\
\overset{(v)}{\leq}&\big\|w^{\prime}-w\big\|\big(\frac{9}{4} a_{\max }^{\frac{4}{3}}\big\|\nabla f_{\xi}\big(w^{\prime}\big)\big\|^{\frac{2}{3}}+\frac{9}{4} a_{\max }^{\frac{4}{3}}\big\|\nabla f_{\xi}(w)\big\|^{\frac{2}{3}}+4 y_{\max } a_{\max }^2\big),
\end{align}
where (i) uses triangular inequality; (ii) uses the fact $a_{\xi}^Tw \leq | a_{\xi}^Tw|$; (iii) uses the fact $\|a_{\xi}a_{\xi}^{\top}\|=\|a_{\xi} \|^2$ and  \eqref{technical ineq}; (iv) uses $(|a_{\xi}^{\top}w'|+|a_{\xi}^{\top}w|)^2\leq 2|a_{\xi}^{\top}w'|^2+2|a_{\xi}^{\top}w|^2$; (v) uses \eqref{PRgradient} and denotes $y_{\max }=\max\|y_{\xi}\|$ and $a_{\max }=\max\big\|a_{\xi}\big\|$. 

The inequality used in (iii) is constructed as following
\begin{align}
\big|\big|a_{\xi}^{\top} w^{\prime}\big|^2-\big|a_{\xi}^{\top} z\big|^2\big|=&\big(\big|a_{\xi}^{\top} w^{\prime}\big|+\big|a_{\xi}^{\top} w\big|\big)\big(\big|a_{\xi}^{\top} w^{\prime}\big|-\big|a_{\xi}^{\top} w\big|\big)\nonumber\\ \leq&\big(\big|a_{\xi}^{\top} w^{\prime}\big|+\big|a_{\xi}^{\top} w\big|\big)\big\|a_{\xi}^{\top}\big(w^{\prime}-w\big)\big\|\nonumber\\ \leq&\big\|a_{\xi}^{\top}\big\|\big(\big|a_{\xi}^{\top} w^{\prime}\big|+\big|a_{\xi}^{\top} w\big|\big)\big\|w^{\prime}-w\big\|.
\label{technical ineq}
\end{align}
Thus, $f_{\xi}(w)=\frac{1}{4}(y_{\xi}-|a_{\xi}^Tw|^2)^2$ satisfies \eqref{symmetric gs} with $\alpha=\frac{2}{3}$, $L_0=4y_{\max}a_{\max}^2$, and $L_1=\frac{9}{2}a^{\frac{4}{3}}_{\max}$. Thus, $F(w)$ satisfies \eqref{symmetric gs stochastic} with $\alpha=\frac{2}{3}$, $L_0=8y_{\max}a_{\max}^2$, and $L_1={9}a^{\frac{4}{3}}_{\max}$.

\end{proof}

% \subsection{Proof for DRO}
% \DROsmooth*
% The proof of generalized-smooth property for Distributionally Robust Optimization is exactly same as \citet{jin2021nonconvex,pmlr-v202-chen23ar}. We refer readers to check Appendix D.2 in \citet{pmlr-v202-chen23ar} or Appendix C in \citet{jin2021nonconvex} for more details, where they have shown the objective satisfies $\|\nabla F(w)-\nabla F(w') \|\leq(L_0+L_1 \|\nabla F(w')\|^{\alpha})\| w-w'\|$. Thus, DRO problem also satisfies induced inequality \ref{symmetric gs} with $\alpha=1$.




\section{Proof of Descent Lemma under Generalized P{\L} condition}\label{Appendix D}
\textbf{Lemma 5 in} \citet{pmlr-v202-chen23ar}.\label{technical lemma}
    For any $x\geq 0$, $C\in [0,1]$, $\Delta>0$, and $0\leq w\leq w'$ such that $\Delta\geq w'-w$, we have the following inequality hold
    \begin{align}
        Cx^{w}\leq x^{w'}+C^{\frac{w'}{\Delta}}.
        \label{Technical Lemma's inequality}
    \end{align}
    
The proof details for this lemma can be found at \citet{pmlr-v202-chen23ar}, Lemma 5 at Appendix.

\begin{lemma}[Descent Lemma under Generalized PL condition]
Let Assumption \ref{assum1} and \ref{assum2} hold. Apply AN-GD with $\beta\in [\alpha,1]$. Choose target error $0 < \epsilon \leq \min\{1, 1/2\mu\}$.
Let step size $\gamma = \frac{(2\mu\epsilon)^{\beta/\rho}}{8(L_0+L_1)+1}$.
Denote $\Delta_t = f(w_t)-f^*$, then we have 
\begin{align}
  \Delta_{t+1} \leq \Delta_t - \frac{\gamma (2\mu)^{\frac{2-\beta}{\rho}}}{2} \Delta_t^{\frac{2-\beta}{\rho}}+\frac{\gamma}{4}(2\mu\epsilon)^{\frac{2-\beta}{\rho}}, \quad \forall t.
\end{align}
\end{lemma}

\begin{proof}
We consider two complementary cases $\alpha>0$ and $\alpha=0$.

When $\alpha>0$, we have $\beta\ge\alpha>0$, and by the descent lemma \ref{descent lemma}, we have that
\begin{align}
& f\big(w_{t+1}\big)-f\big(w_t\big) \nonumber\\
\overset{(i)}{\leq}& \nabla f\big(w_t\big)^{\top}\big(w_{t+1}-w_t\big)+\frac{1}{2}\big(L_0+L_1\big\|\nabla f\big(w_t\big)\big\|^\alpha\big)\big\|w_{t+1}-w_t\big\|^2 \nonumber\\
\overset{(ii)}{=}&-\gamma\big\|\nabla f\big(w_t\big)\big\|^{2-\beta}+\frac{\gamma}{4}\big(2 L_0 \gamma \cdot\big\|\nabla f\big(w_t\big)\big\|^{2-2 \beta}+2 L_1 \gamma \cdot\big\|\nabla f\big(w_t\big)\big\|^{2+\alpha-2 \beta}\big) \nonumber\\
\overset{(iii)}{\leq}&-\gamma\big\|\nabla f\big(w_t\big)\big\|^{2-\beta} +\frac{\gamma}{4}\big(2\big\|\nabla f\big(w_t\big)\big\|^{2-\beta}+\big(2 L_0 \gamma\big)^{\frac{2}{\beta}-1}+\big(2 L_1 \gamma\big)^{\frac{2}{\beta}-1}\big) \nonumber\\
\overset{(iv)}{\leq}&-\frac{\gamma}{2}\big\|\nabla f\big(w_t\big)\big\|^{2-\beta}+\frac{\gamma^{\frac{2}{\beta}}}{4}\big(2 L_0+2 L_1\big)^{\frac{2}{\beta}-1} \nonumber\\
\overset{(v)}{\leq}&-\frac{\gamma}{2}\big\|\nabla f\big(w_t\big)\big\|^{2-\beta}+\frac{(2\mu\epsilon)^{\frac{2}{\rho}}}{(8(L_0+L_1)+1)^{\frac{2}{\beta}}}(\frac{1}{4})(8(L_0+L_1)+1)^{\frac{2}{\beta}-1} \nonumber\\
\overset{(vi)}{\leq}& -\frac{\gamma}{2}\big\|\nabla f\big(w_t\big)\big\|^{2-\beta} +\frac{1}{4}\frac{(2\mu\epsilon)^{\frac{\beta}{\rho}}(2\mu\epsilon)^{\frac{2-\beta}{\rho}}}{8(L_0+L_1)+1}\nonumber\\
=& -\frac{\gamma}{2}\big\|\nabla f\big(w_t\big)\big\|^{2-\beta}+\frac{\gamma}{4} (2\mu\epsilon)^{\frac{2-\beta}{\rho}},
\end{align}
where (i) follows from the descent lemma \ref{descent lemma}; (ii) follows from the AN-GD update rule by replacing $w_{t+1}-w_t$ with $-\gamma\frac{\nabla f(w_t)}{\|\nabla f(w_t) \|^{\beta}}$; (iii) follows from the technical lemma 5 above proved in \citet{pmlr-v202-chen23ar} by letting $\omega' = 2-\beta$, $\Delta = \beta$ and substituting \eqref{Technical Lemma's inequality} to $2L_0 \gamma \|\nabla f(w_t) \|^{2-2\beta}$ and $2L_1 \gamma \|\nabla f(w_t) \|^{2+\alpha-2\beta}$ (Let $x=\|\nabla f(w_t) \|$ and $C=2L_0\gamma \text{ or }2L_1\gamma$ , $C\in[0,1]$ because $\epsilon \leq 1/2\mu$);
%Since by choosing $\epsilon \in \min \big \{1, \frac{1}{2\mu} \big \}$, we have $2\mu \epsilon <1$ holds, thus, we must have $\gamma<1$, and $2K_0\gamma<1, 2K_1\gamma<1$.
(iv) follows from the fact that $a^{\tau}+b^{\tau}\leq(a+b)^{\tau}$ with $\tau = {2}/{\beta}-1\geq1$ and $a,b\geq 0$, $(v)$ follows from the step size rule $\gamma = {(2\mu\epsilon)^{\beta/\rho}}/{(8(L_0+L_1)+1)}$; (vi) follows from the facts that $2L_0+2L_1 \leq 8(L_0+L_1)+1$ and $(2L_0+2L_1)^{\frac{2}{\beta}-1} \leq (8(L_0+L_1)+1)^{\frac{2}{\beta}-1}$ when $\beta>0$.

When $\alpha=0$, the same argument applies if $\beta>0$. Next, we consider the case $\alpha =\beta =0$, and we have $\gamma = \frac{1}{8(L_0+L_1)+1}$. We obtain that
 \begin{align}
     & f\big(w_{t+1}\big)-f\big(w_t\big) \nonumber\\
\overset{(i)}{\leq}& \nabla f\big(w_t\big)^{\top}\big(w_{t+1}-w_t\big)+\frac{1}{2}\big(L_0+L_1\big)\big\|w_{t+1}-w_t\big\|^2 \nonumber\\
\overset{(ii)}{=}&-\gamma\big\|\nabla f\big(w_t\big)\big\|^{2}+\frac{\gamma}{4}\big(2 L_0 \gamma \cdot\big\|\nabla f\big(w_t\big)\big\|^{2}+2 L_1 \gamma \cdot\big\|\nabla f\big(w_t\big)\big\|^{2}\big) \nonumber\\
\overset{}{=}&-\gamma\big\|\nabla f\big(w_t\big)\big\|^{2} +\frac{\gamma^2}{4}\big(2L_0+2L_1\big)\big\|\nabla f(w_t) \big\|^2 \nonumber\\
\overset{(iii)}{\leq }&-{\gamma}\big\|\nabla f\big(w_t\big)\big\|^{2}+\frac{\gamma}{4}\frac{1}{8(L_0+L_1)+1}\big(8 (L_0+ L_1)+1\big) \big\|\nabla f(w_t) \big\|^2\nonumber\\
\overset{(iv)}{\leq}&-{\gamma}\big\|\nabla f\big(w_t\big)\big\|^{2}+\frac{\gamma}{2}\big \|\nabla f(w_t) \big \|^2 \nonumber\\
\overset{(v)}{\leq} &-{\gamma}\big\|\nabla f\big(w_t\big)\big\|^{2}+\frac{\gamma}{2}\big \|\nabla f(w_t) \big \|^2+\frac{\gamma}{4} (2\mu\epsilon)^{\frac{2}{\rho}} \nonumber\\
=& -\frac{\gamma}{2}\big\|\nabla f\big(w_t\big)\big\|^{2}+\frac{\gamma}{4} (2\mu\epsilon)^{\frac{2}{\rho}},
\end{align}
where (i) follows from the descent lemma \ref{descent lemma} with $\alpha=0$; (ii) follows from the update rule of AN-GD by replacing $w_{t+1}-w_t$ with $-\gamma \nabla f(w_t)$ (Normalization term vanishes due to $\beta=0$); (iii) follows the fact that $2L_0+2L_1\leq 8(L_0+L_1)+1$;
%Since by choosing $\epsilon \in \min \big \{1, \frac{1}{2\mu} \big \}$, we have $2\mu \epsilon <1$ holds, thus, we must have $\gamma<1$, and $2K_0\gamma<1, 2K_1\gamma<1$.
(iv) follows from the fact $\frac{\gamma}{4}\leq \frac{\gamma}{2}$ and (v) holds by adding additional positive number $\frac{\gamma}{4} (2\mu\epsilon)^{\frac{2}{\rho}}$.

Thus, in conclusion, for any $\beta \in [\alpha,1]$, we always have the following descent lemma
\begin{align}
    f\big(w_{t+1}\big)-f\big(w_t\big)\leq -\frac{\gamma}{2}\big\|\nabla f\big(w_t\big)\big\|^{2-\beta}+\frac{\gamma}{4} (2\mu\epsilon)^{\frac{2-\beta}{\rho}}
    \label{unified descent lemma}
\end{align}
under the step size rule $\gamma = \frac{(2\mu\epsilon)^{\beta/\rho}}{8(L_0+L_1)+1}$. Moreover, by the generalized P{\L}-condition in Definition \ref{gen PL}, we have
\begin{align}
    \big \|\nabla f(w) \big \| \geq (2\mu)^{\frac{1}{\rho}} \big(f(w)-f^*\big)^{\frac{1}{\rho}}, \nonumber
\end{align}
which implies that
\begin{align}
    \big \|\nabla f(w_t) \big \|^{2-\beta} \geq (2\mu)^{\frac{2-\beta}{\rho}} \big(f(w_t)-f^*\big)^{\frac{2-\beta}{\rho}} = (2\mu)^{\frac{2-\beta}{\rho}}\Delta_t^{\frac{2-\beta}{\rho}}.
    \label{equalPL}
\end{align}
Substituting \eqref{equalPL} into \eqref{unified descent lemma}, we have
\begin{align}
    f(w_{t+1}) - f(w_t) \leq -\frac{\gamma}{2}(2\mu)^{\frac{2-\beta}{\rho}}\big(f(w_{t}) - f^* \big)^{\frac{2-\beta}{\rho}} + \frac{\gamma}{4}(2\mu\epsilon)^{\frac{2-\beta}{\rho}}. \nonumber
\end{align}
Subtracting $f^*$ on both sides, we have 
\begin{align}
    &f(w_{t+1})-f^* \leq f(w_t) - f^* -\frac{\gamma}{2}(2\mu)^{\frac{2-\beta}{\rho}}(f(w)-f^*)^{\frac{2-\beta}{\rho}} + \frac{\gamma}{4}(2\mu\epsilon)^{\frac{2-\beta}{\rho}}. \nonumber
\end{align}
Denote $\Delta_t = f(w_t) - f^*$, the above inequality can be written as 
\begin{align}
    \Delta_{t+1} \leq \Delta_t- \frac{\gamma(2\mu)^{\frac{2-\beta}{\rho}}}{2}\Delta_t^{\frac{2-\beta}{\rho}}+\frac{\gamma}{4} (2\mu\epsilon)^{\frac{2-\beta}{\rho}}, \quad \forall t. 
    \label{descentPL}
\end{align}

\end{proof}
\section{Proof of Theorem \ref{beta GD convergence}}\label{Appendix E}
\mainthmone*
\begin{proof}\label{Proof of beta GD convergence}
Given target error $0< \epsilon \leq \min \{ 1,1/{2\mu}\}$,
define stopping time as the first time that $\Delta_t$ satisfies $\Delta_t\leq \epsilon$, i.e., $
T= \inf \{t | \Delta_t \leq \epsilon \}$. In the following proof, we provide an estimate of such stopping time $T$.


For any $t$ satisfies $\Delta_t > \epsilon$, we must have
$\Delta_t^{(2-\beta)/{\rho}}>\epsilon^{(2-\beta)/{\rho}}$ since $\frac{2-\beta}{\rho}>0$. This implies that  $\frac{\gamma}{4}(2\mu\epsilon)^{(2-\beta)/{\rho}}<\frac{\gamma}{4}(2\mu)^{(2-\beta)/{\rho}}\Delta_t^{(2-\beta)/{\rho}}$. Thus, \eqref{descentPL} further reduces to
the following descent inequality
\begin{align}
\Delta_{t+1} \leq \Delta_t - \frac{\gamma (2\mu)^{\frac{2-\beta}{\rho}}}{4} \Delta_t^{\frac{2-\beta}{\rho}}. 
\label{descentPL2}
\end{align}



%We first show that as long as \eqref{descentPL2} holds, the generated sequence is guaranteed to converge to 0 as $t \bigarrow \infty$.
%Thus, $\epsilon$-stationary is reachable in finite time $T$ and $\Delta_t$ is non-increasing.

%By definition of $\Delta_t$, $\mu$, $\beta$, $\rho$, we must have $\Delta_t$, $\frac{(2\mu)^{(2-\beta)/\rho}}{4},\gamma $ are non-negative. 
%It's obvious that $\{\Delta \}_t$ generated by \eqref{descentPL2} is non-increasing and non-negative. Now, Suppose the sequence $\{\Delta_t \}$ converges to a positive constant $\kappa$ as $t\bigarrow \infty$. 
%Then, by definition of sequential convergence and non-increasing property of $\{\Delta_t \}$, we conclude $\Delta_t > \frac{\kappa}{2}$ holds before termination, and
%\begin{align}
%\Delta_{t+1}\leq \Delta_t - \gamma\frac{(2\mu)^{(2-\beta)/\rho}}{4}\Delta_t^{\frac{2-\beta}{\rho}}\leq \Delta_t - \gamma\frac{(2\mu)^{(2-\beta)/\rho}}{4}(\frac{\kappa}{2})^{\frac{2-\beta}{\rho}}. 
%\label{refinedrecursioncase3}
%\end{align}
%However, under this scenario, we can find a finite time $T$ from relation
%\begin{align}
%    T\Big(\gamma\frac{(2\mu)^{(2-\beta)/\rho}}{4}(\frac{\kappa}{2})^{\frac{1}{\omega}} \Big)= \Delta_0-\kappa.
%\end{align}
%Thus, we conclude the sequence $\{\Delta_t \}$ reaches $\kappa$ within finite time. This contradicts the assumption that $\{\Delta_t\}$ converges to $\kappa$. We conclude that as long as \eqref{descentPL2} holds, the sequence $\{ \Delta_t\}$ is non-increasing and can reach $\epsilon$-stationary point under finite time.
For above \eqref{descentPL2}, we provide convergence rate based on the value of $\rho$ and $\beta$.\\
\noindent\textbf{Case I: $\beta<2-\rho$}\\
This is equivalent to $\frac{2-\beta}{\rho}>1$. For simplicity of notation, denote $\theta=\frac{2-\beta}{\rho}$.
When $\theta > 1$, we have following inequalities
\begin{align}
    \Delta_{t+1}\leq& \Delta_t \nonumber\\
    \Delta_{t+1}^{\theta} \leq& \Delta_t^{\theta} \nonumber \\
    \Delta_{t+1}^{-\theta} \geq& \Delta_t^{-\theta}.
    \label{delta_relation}
\end{align}
Now define an auxiliary function $\Phi(t)=\frac{1}{\theta-1}t^{1-\theta}$, where its derivative is $\Phi'(t)=-t^{-\theta}$ 
We divide the last inequality of \eqref{delta_relation} into two different cases for analysis. One is the case when $\Delta_t^{-\theta}\leq\Delta_{t+1}^{-\theta}\leq 2\Delta_t^{-\theta}$, another is the case when $\Delta_{t+1}^{-\theta}\geq 2\Delta_t^{-\theta}$.

When $\Delta_{t}^{-\theta}\leq \Delta_{t+1}^{-\theta }\leq 2 \Delta_t^{-\theta}$, we have
\begin{align}
    \Phi(\Delta_{t+1})-\Phi(\Delta_t) =& \int_{\Delta_t}^{\Delta_{t+1}}\Phi'(t)\mathrm{d}t = \int_{\Delta_{t+1}}^{\Delta_t}t^{-\theta}\mathrm{d}t \nonumber\\
    &\overset{(i)}{\geq} (\Delta_{t}-\Delta_{t+1})\Delta_t^{-\theta} \nonumber\\
    &\overset{(ii)}{\geq} (\Delta_t -\Delta_{t+1})\frac{\Delta_{t+1}^{-\theta}}{2} \nonumber\\
    &\overset{(iii)}{\geq} \frac{\gamma (2\mu)^{\theta}}{4}\Delta_t^{\theta}\frac{\Delta_{t+1}^{-\theta}}{2} \nonumber \\
    &\overset{(iv)}{\geq} \frac{\gamma (2\mu)^{\theta}}{4}\Delta_{t+1}^{\theta}\frac{\Delta_{t+1}^{-\theta}}{2}=\frac{\gamma (2\mu)^{\theta}}{8},\nonumber
\end{align}
where $(i)$ applies mean value theorem such that $\Phi(\Delta_{t+1})-\Phi(\Delta_t) = |\Delta_{t+1}-\Delta_{t} \| \Phi'(\xi)|$ and $\Delta_{t+1}^{-\theta}\geq|\Phi'(\xi)|\geq \Delta^{-\theta}_{t}$ holds for any $\xi \in [\Delta_{t+1}, \Delta_{t}]$; (ii) uses the fact $\Delta_{t+1}^{-\theta}\leq 2\Delta_{t}^{-\theta}$; (iii) utilizes recursion $\Delta_t - \Delta_{t+1} \geq \frac{\gamma (2\mu)^{\theta}}{4}\Delta_t^{\theta}$; (iv) uses the fact that $\Delta_{t}^{\theta} > \Delta_{t+1}^{\theta}$ for all $\theta >0$. 

When $\Delta_{t+1}^{-\theta}> 2 \Delta_t^{-\theta}$, it holds that $\Delta_{t+1}^{1-\theta}=(\Delta_{t+1}^{-\theta})^{\frac{1-\theta}{-\theta}} > 2^{\frac{1-\theta}{-\theta}}\Delta_t^{1-\theta} $.
Then, we have
\begin{align}
    \Phi(\Delta_{t+1})-\Phi(\Delta_t)=&\frac{1}{\theta-1}(\Delta_{t+1}^{1-\theta}-\Delta_t^{1-\theta}) \nonumber\\
    \overset{(i)}{\geq}& \frac{1}{\theta-1}\big((2)^{\frac{\theta-1}{\theta}}-1\big)\Delta_t^{1-\theta} \nonumber\\
    \overset{(ii)}{\geq}& \frac{1}{\theta-1}\big((2)^{\frac{\theta-1}{\theta}}-1\big)\Delta_0^{1-\theta},\nonumber
\end{align}
where (i) is due to the recursion $\Delta_{t+1}^{1-\theta}=(\Delta_{t+1}^{-\theta})^{\frac{1-\theta}{-\theta}} > 2^{\frac{1-\theta}{-\theta}}\Delta_t^{1-\theta} $; (ii) is due to the fact the sequence $\{\Delta_t\}_{t=1}^{T}$ generated by \eqref{descentPL2} is non-increasing. 
Substitute $\theta=\frac{2-\beta}{\rho}$ in and denote 
\begin{align}
    C=\min\Big\{\frac{\gamma(2\mu)^{\frac{2-\beta}{\rho}}}{8},\frac{\rho}{2-\beta-\rho}(2^{\frac{2-\beta-\rho}{2-\beta}}-1)\Delta_0^{\frac{\beta+\rho-2}{\rho}}  \Big\}. \nonumber
\end{align}
We conclude for all $t$, we have
\begin{align}
    \Phi(\Delta_t)\geq \Phi(\Delta_t)-\Phi(\Delta_0)= \sum_{i=0}^{t-1}\Phi(\Delta_{i+1})-\Phi(\Delta_{i})\geq C t, \nonumber
\end{align}
Substituting the expression $\Phi(\Delta_t)=\frac{1}{\theta-1}\Delta_t^{1-\theta}$ into the above inequality, we have
\begin{align}
    \Delta_{t}\leq \Big(\frac{\rho}{(2-\beta-\rho)\cdot Ct}\Big)^{\frac{\rho}{2-\rho-\beta}},
    %= \mathcal{O}\Big(\frac{\rho}{(2-\beta-\rho)\gamma t}\Big)^{\frac{\rho}{2-\rho-\beta}} ,
\end{align}
%We divide the convergence proof into two cases with different values of C.
In order to obtain $\Delta_T\leq \epsilon$,
when $C=\frac{\rho}{2-\beta-\rho}(2^{{(2-\beta-\rho)}/{(2-\beta)}}-1)\Delta_0^{{(\beta+\rho-2)}/{\rho}}$, %take logarithm on both sides, we have
%the step-size rule $\gamma= \Theta(\epsilon^{{\beta}/{\rho}})$ in $\frac{\gamma (2\mu)^{(2-\beta)/{\rho}}}{8}$ doesn't effect the result, so we have
%\begin{align}
%    \frac{\rho}{{2-\rho-\beta}}\cdot \log\Big(\frac{(2-\beta-\rho)\cdot CT}{\rho}\Big)=\log\Big(\frac{1}{\epsilon}\Big), \nonumber
%\end{align}
T should satisfy 
\begin{align}
T= \frac{1}{(2^{(2-\beta-\rho)/(2-\beta)}-1)\Delta_0^{{(\rho+\beta-2)}/{\rho}}\epsilon^{(2-\beta-\rho)/{\rho}}};
\end{align}
When $C= \frac{\gamma(2\mu)^{{2-\beta}/{\rho}}}{8}=\frac{(2\mu)^{2/\rho}\cdot \epsilon^{\beta/\rho}}{8(8(L_0+L_1)+1)}$, 
%taking logarithm on above inequality, in order to obtain $\Delta_T \leq \epsilon$, we have
%\begin{align}
%    \frac{\rho}{2-\beta-\rho}\cdot\log\Big(\frac{(2-\beta-\rho)(2\mu)^{2/\rho}\epsilon^{{\beta}/{\rho}}T}{8\rho(8(L_0+L_1)+1)}\Big)= \log\Big(\frac{1}{\epsilon}\Big). \nonumber
%\end{align}
T should satisfy
\begin{align}
T=\frac{8\rho(8(L_0+L_1)+1)}{(2-\beta-\rho)(2\mu)^{2/\rho}\epsilon^{(2-\rho)/{\rho}}}.
\end{align}

This concludes when $\rho+\beta<2$, T should satisfy 
\begin{align}
    T\geq \max \Big\{\frac{8\rho(8(L_0+L_1)+1)}{(2-\beta-\rho)(2\mu)^{2/\rho}\epsilon^{(2-\rho)/{\rho}}}, \frac{1}{(2^{(2-\beta-\rho)/(2-\beta)}-1)\Delta_0^{{(\rho+\beta-2)}/{\rho}}\epsilon^{(2-\beta-\rho)/{\rho}}} \Big\} = \Omega\big((\frac{1}{\epsilon})^{\frac{2-\rho}{\rho}}\big).
\end{align}
%we have
%\begin{align}
%   \log(T)= \log\big((\frac{1}{\epsilon})^{\frac{\rho}{2-\beta-\rho}}\big) +\log\big((\frac{1}{\epsilon})^{\frac{\beta}{\rho}}\big)+\log(\frac{1}{\widetilde{C}})+\log\Big(\frac{\rho}{2-\beta-\rho}\Big),\nonumber
%\end{align}
%Thus, when $0\leq \rho\leq 2-2\beta$, we have $T=\mathcal{O}(\frac{1}{\epsilon})^{\frac{2-\rho-\beta}{\rho}}$; when $2-2\beta < \rho \leq  2-\beta$, we have $T= \mathcal{O}((\frac{1}{\epsilon})^{\frac{\beta}{\rho}})$.
%Thus,In order obtain an point satisfying $\Delta_T\leq \epsilon$, the algorithm complexity should satisfy
%\begin{align}
    %\mathcal{O}\big((\frac{1}{\epsilon})^{\frac{2-\rho-\beta}{\rho}}\big)\leq T\leq \mathcal{O}\big((\frac{1}{\epsilon})^{\frac{2-\rho-\beta}{\rho}}+(\frac{1}{\epsilon})^{\frac{\beta}{\rho}})
%\end{align}

\noindent\textbf{Case II: $\beta=2-\rho$}

This is equivalent to $\frac{2-\beta}{\rho}=1$, \eqref{descentPL2} reduces to 
\begin{align}
    \Delta_{t+1}\leq \Delta_t - \frac{\gamma \mu}{2} \Delta_t = (1-\frac{\gamma\mu}{2})\Delta_t,\nonumber
\end{align}
which leads to
\begin{align}
    \Delta_t \leq (1-\frac{\gamma\mu}{2})^t \Delta_0 = \mathcal{O}\Big((1-\frac{\gamma \mu}{2})^t\Big). \nonumber
\end{align}
To obtain $\Delta_T\leq \epsilon$, we have
\begin{align}
    \Delta_t \leq (1-\frac{\gamma \mu}{2})^t \Delta_0 \leq \exp(-\frac{\gamma\mu t}{2})\Delta_0.
\end{align}
To obtain $\Delta_t \leq \epsilon$, iteration complexity should satisfy
\begin{align}
T\geq \frac{2}{\gamma \mu}\log(\frac{\Delta_0}{\epsilon})= \frac{2^{1-\beta/\rho}\cdot(8(L_0+L_1)+1)}{\mu\cdot(\mu\epsilon)^{\beta/\rho}}\cdot\log(\frac{\Delta_0}{\epsilon})=\Omega\big((\frac{1}{\epsilon})^{\frac{\beta}{\rho}}\cdot \log(\frac{\Delta_0}{\epsilon})\big)\nonumber.
\end{align}

\noindent\textbf{Case III: $1\geq\beta>2-\rho$} 

This is equivalent to $\frac{2-\beta}{\rho}<1$.

%If $\epsilon \leq \big(\frac{(2\mu)^{\frac{2}{\rho}}}{(48(K_0+K_1+2K_2)+4)}\big)^{\frac{\rho}{\rho-\beta}}$, this condition implies $\epsilon \leq C\gamma$, Thus, the sequence of $\big \{\Delta \big\}_t$ will first reach $C\gamma$, i.e., 
%There exists $T_0$ such that $\Delta_t \geq C\gamma$ for all $0\leq t<T_0$. We first estimate the order of  convergence rate for $\big \{\Delta \big \}_t$ to $C\gamma$. Since \eqref{descentPL2} guarantees non-incresing of $\big\{ \Delta_t\big\}$, we have
%\begin{align}
    %\Delta_0\geq \Delta_0 - \Delta_{t_0} = \sum_{t=0}^{t=T_0-1}\Delta_{t}-\Delta_{t+1}\geq \sum_{t=0}^{t=T_0-1} C\gamma\Delta_t^{\frac{1}{\omega}}\geq t_0 (C\gamma)^{1+\frac{1}{\omega}} \nonumber
%\end{align}
    %where the last inequality uses the fact that $0<\frac{1}{\omega}<1$ and $\Delta_t \geq C\gamma$ for all $t>t_0$.\\
    %Re-organize above inequality, we have $T_0 \leq \frac{\Delta_0}{(C\gamma)^{1+\frac{1}{\omega}}}\sim \mathcal{O}\big(\big(\frac{1}{\epsilon}\big)^{\frac{(2-\beta+\rho)\beta}{\rho^2}}\big)$.\\
For simplicity of notation,  denote $\omega = \frac{\rho}{2-\beta}$ and $C=\frac{ (2\mu)^{\frac{2-\beta}{\rho}}}{4}$. Then,
\eqref{descentPL2} can be rewritten as
\begin{align}
    \Delta_{t+1}\leq \Delta_t - C\gamma\Delta_t^{\frac{1}{\omega}}.
\end{align}
When $\Delta_t$ is small enough,  $\Delta_{t+1}^{{1}/{\omega}}$ will dominate $\Delta_{t+1}$ order-wisely since $1/\omega<1$. This fact leads to
    \begin{align}
        C\gamma \Delta_{t+1}^{\frac{1}{\omega}}\overset{(i)}{\leq} \Delta_{t+1}+ C\gamma \Delta_{t+1}^{\frac{1}{\omega}}\overset{(ii)}{\leq} \Delta_{t+1}+ C\gamma \Delta_{t}^{\frac{1}{\omega}}\overset{(iii)}{\leq} \Delta_t, \nonumber
    \end{align}
    where (i) is because $\Delta_{t+1}\geq 0$; (ii) is because $\Delta_{t+1}\leq \Delta_t$; (iii) is because re-organization of \eqref{descentPL2}. 
    %When $(C\gamma)^{\omega/(\omega-1)}>\epsilon$ holds, that is, there exists some $\epsilon$ within $0<\epsilon<1/2\mu$ satisfying
    %\begin{align}
    %    \epsilon >\frac{(32(L_0+L_1)+4)^{\frac{\rho}{2-\rho}}}{(2\mu)^{\frac{2}{2-\rho}}}(\text{when }\rho<2),\text{ or }
    %    1>\frac{16(L_0+L_1)+2}{\mu}(\text{when }\rho=2).\nonumber   
    %\end{align}
    There exists a time $T_0= \inf \{ t \in \mathbf{N}| \Delta_{t}/{(C\gamma)}^{\omega/(\omega-1)}< 1\}$ such that for any $T_0\leq t\leq T$, we have
    \begin{align}
        \Delta_{t+1}\leq (C\gamma)^{-\omega}\Delta_t^{\omega}=&(C\gamma)^{-\omega-\omega^2-...-\omega^{t-T_0}}\Delta_{T_0}^{\omega^{t-T_0}}\nonumber\\
        =& (C\gamma)^{\frac{\omega(1-\omega^{t-T_0})}{\omega-1}} \Delta_{T_0}^{\omega^{t-T_0}}\nonumber\\
        =& (C\gamma)^{{\omega}/{\omega-1}}\big((C\gamma)^{{\omega}/{\omega-1}}\big)^{\omega^{-{t-T_0}}}\Delta_{T_0}^{\omega^{t-T_0}}.
    \end{align}
    %Since $\epsilon \in (0,1)$, $\rho,\beta>0$ and $\rho+\beta>2$ leads to $\frac{\beta}{\rho}>0$, and$\frac{\omega}{\omega -1}=\frac{\rho}{\rho+\beta -2 }>1$.
    %Since  $(C\gamma)^{\omega/{\omega-1}}$ only affects order of convergence up to a constant. To simplify analysis, 
    To simplify analysis, denote all parameters associated with $\epsilon$ as $\hat{C}=(\frac{C(2\mu)^{{\beta}/{\rho}}}{8(L_0+L_1)+1})^{{\omega}/{\omega-1}}$. We have $(C\gamma)^{{\omega}/{\omega-1}} = \hat{C} \epsilon^{\beta/{(\rho+\beta-2)}}\leq \hat{C}$. 
    This enables us to further reduce the recursion to
    \begin{align}
        \Delta_{t+1}\leq&(C\gamma)^{{\omega}/{\omega-1}}\big((C\gamma)^{{\omega}/{\omega-1}}\big)^{\omega^{-{t-T_0}}}\Delta_{T_0}^{\omega^{t-T_0}} \nonumber\\
        \leq& \hat{C} \Big((C\gamma)^{\frac{\omega}{\omega-1}}\Big)^{{-\omega}^{t-T_0}}\Delta_{T_0}^{{\omega}^{t-T_0}} \nonumber\\
        =& \mathcal{O}\bigg(\Big(\frac{\Delta_{T_0}}{\gamma^{\omega/{\omega-1}}}\Big)^{\omega^{t-T_0}}\bigg).
    \label{descentPL3}
    \end{align} 
    %Now, we estimate the order of convergence from $T_0$ to $T$. For simplicity, we use $t$ to replace $T-T_0$ doing analysis,
    To obtain $\Delta_T\leq \epsilon$,
    above recursion should satisfy
    \begin{align}
        \hat{C} \Big((C\gamma)^{\frac{\omega}{\omega-1}}\Big)^{{-\omega}^{t-T_0}}\Delta_{T_0}^{{\omega}^{t-T_0}}=\epsilon. \nonumber
    \end{align}
    Taking logarithm on both sides and extracting $\epsilon^{\beta/(\rho+\beta-2)}$ from ${(C\gamma)^{{\omega}/{\omega-1}}}/{\Delta_{T_0}}$, we have
    \begin{align}
        \log({\hat{C}}/{\epsilon})= \omega^{t-T_0} \cdot \log\Big({(C\gamma)^{\frac{\omega}{\omega-1}}}/{\Delta_{T_0}}\Big)=\omega^{t-T_0}\cdot\log\Big(\big(\hat{C}/{\Delta_{T_0}}\big)\cdot \epsilon^{\frac{\beta}{\rho+\beta-2}}\Big) \leq \omega^{t-T_0}\big(\hat{C}/{\Delta_{T_0}}\big)\cdot \epsilon^{\frac{\beta}{\rho+\beta-2}}\nonumber.
        %\leq \omega^t \big(\log(\frac{\hat{C}^{\frac{\omega}{\omega-1}}}{\Delta_{T_0}})\big)
        %\leq \omega^t(\frac{(C\gamma)^{\frac{\omega}{\omega-1}}}{\Delta_{T_0}})\nonumber  
    \end{align}
    %By definition of $T_0$, we have ${\Delta_{T_0}}/{(C\gamma)^{\frac{\omega}{\omega-1}}}<1$ in order to make the recursion converge. 
    %If  $\frac{\Delta_{T_0}}{(C\gamma)^{\frac{\omega}{\omega-1}}}>1$ for all $T\geq T_0$, it contradicts the fact $\Delta_t$ will converge to 0 as $t\bigarrow \infty$.
    
    %\begin{align}
    %\log\bigg({(C\gamma)^{\frac{\omega}{\omega-1}}}/{\Delta_{T_0}}\bigg) 
    %= \big(\frac{C(2\mu\epsilon)^{\frac{\beta}{\rho}}}{8(L_0+L_1)+1}\big)^{\frac{\omega}{\omega-1}} /\Delta_{T_0}
    %={{\hat{C}/\Delta_{T_0}}}/{(\frac{1}{\epsilon}^{\frac{\beta}{\rho}})^{\frac{\omega}{\omega-1}}}
    %=&\log\bigg(\big(\hat{C}/{\Delta_{T_0}}\big)\cdot \epsilon^{\frac{\beta}{\rho+\beta-2}}\bigg) \nonumber \\
    %\leq& \big(\hat{C}/{\Delta_{T_0}}\big)\cdot \epsilon^{\frac{\beta}{\rho+\beta-2}}, \nonumber
    %- \frac{\beta
    %}{\rho+\beta-2}\log\bigg(\frac{1}{\epsilon}\bigg)> 0. \nonumber
    %\end{align}
    %where the last inequality is due to the fact $\log(x)\leq x, \forall x>0$.
    %Since $\frac{\Delta_{T_0}}{(C\gamma)^{\frac{\omega}{\omega-1}}}<1$,  we must have $\log(\frac{\hat{C}}{\Delta_{T_0}})- \frac{\beta
    %}{\rho+\beta-2}\log(\frac{1}{\epsilon})>0$
    Taking logarithm again on both sides of the above inequality, we have
    \begin{align}
        t-T_0
        %= \frac{\log(\frac{\hat{C}}{\epsilon})}{\log(\frac{\hat{C}}{\Delta_{T_0}})- \frac{\beta
        %}{\rho+\beta-2}\log(\frac{1}{\epsilon})} 
        \geq& \frac{1}{\log\big(\frac{\rho}{2-\beta})}\cdot\Big[\log\Big(\frac{\beta}{\rho+\beta-2}\cdot\log\Big(\frac{1}{\epsilon}\Big)+\log\Big(\log\Big(\frac{(2\mu)^{2/(\rho+\beta-2)}}{(32(L_0+L_1)+4)^{{\rho}/{(\rho+\beta-2)}}\epsilon}\Big)\Big)\nonumber\\
        &\qquad\qquad\qquad\qquad\qquad\qquad \qquad\qquad+\log\Big(\frac{(32(L_0+L_1)+4)^{{\rho}/{(\rho+\beta-2)}}\Delta_{T_0}}{(2\mu)^{2/(\rho+\beta-2)}}\Big)\Big]\nonumber\\
        \gtrsim  &\Omega\Big(\log\Big(\frac{\beta}{\rho+\beta-2}\cdot\log\Big(\frac{1}{\epsilon}\Big)+\log\Big(\log\Big(\frac{(2\mu)^{2/(\rho+\beta-2)}}{(32(L_0+L_1)+4)^{\rho/(\rho+\beta-2)}\epsilon}\Big)\Big)\Big) \nonumber\\
        =&\Omega \big(\log\big((\frac{1}{\epsilon})^{\frac{\beta}{\rho+\beta-2}}\big)\big). \nonumber
    \end{align}
    %Take logrithm on both sides again, we have
%\begin{align}
%    t = \frac{1}{\log(\frac{\rho}{2-\beta})}\log\big( \frac{\log(\frac{\hat{C}}{\epsilon})}{\log(\frac{\hat{C}}{\Delta_{T_0}})} \big) \sim \mathcal{O}\big(\log(\log(\frac{1}{\epsilon}))\big)
%\end{align}


    %\item if $\epsilon > \big(\frac{(2\mu)^{\frac{2}{\rho}}}{(48(K_0+K_1+2K_2)+4)}\big)^{\frac{\rho}{\rho-\beta}}$, that indicates $\epsilon > C\gamma$.  The iteration will terminates early before reaching the phase II descent derived at \eqref{descentPL3}. Similarly as above analysis, we have
    %\begin{align}
        %\Delta_0\geq \Delta_0 - \Delta_{t_0} = \sum_{t=0}^{t=T_0-1}\Delta_{t}-\Delta_{t+1}\geq \sum_{t=0}^{t=T_0-1} C\gamma\Delta_t^{\frac{1}{\omega}}\geq T_0 C\gamma \epsilon^{\frac{1}{\omega}},\nonumber
    %\end{align}
    %where the last inequality uses the fact $0<\frac{1}{\omega}<1$, and $\Delta_t > \epsilon > C\gamma$.
    %Re-organize above inequality we have $T\sim \mathcal{O}\big(\big(\frac{1}{\epsilon}\big)^{\frac{2}{\rho}}\big)$.
\end{proof}
\section{Proof of Proposition \ref{GD convergence}}\label{Appendix F}
\gdconvergence*
\begin{proof}\label{Proof of GD convergence}
    When $\alpha=1$, putting the update rule of gradient descent
    $w_{t+1}=w_t - \gamma \nabla f(w_t)$ into descent lemma \eqref{descent lemma} yields
    \begin{align}
        &f(w_{t+1})\nonumber\\
        \overset{(i)}{\leq}& f(w_t)-\gamma \big \|\nabla f(w_t) \big \|^2+\frac{\gamma^2}{2}(L_0+L_1\big\| \nabla f(w_t)\big\|)
        \big \|\nabla f(w_t) \big \|^2\nonumber \\
        =&f(w_t)-(\gamma-\frac{L_0\gamma^2}{2})\big\|\nabla f(w_t) \big\|^2+\frac{L_1\gamma^2}{2}\big \|\nabla f(w_t) \big \|^3\nonumber\\
        \overset{(ii)}{\leq}& f(w_t)-\frac{\gamma}{2}\big\|\nabla f(w_t) \big\|^2+\frac{L_1\gamma^2}{2}\big\| \nabla f(w_t)\big\|^3\nonumber\\
        \overset{(iii)}{\leq}& f(w_t) - \frac{\gamma}{4}\big\|\nabla f(w_t) \big\|^2
    \end{align}
    where (i) is due to descent lemma \eqref{descent lemma}; (ii) and (iii) are due to the learning rate design $\gamma \leq \min \{\frac{1}{L_0},\frac{1}{2L_1G}  \}$, which ensures $-(\gamma-\frac{L_0\gamma^2}{2})\leq -\frac{\gamma}{2}$ and $\frac{L_1\gamma^2}{2}\| \nabla f(w_t)\|^3\leq \frac{\gamma}{4}\|\nabla f(w_t) \|^2$.
    
    When applying Assumption \ref{assum2} with $\rho=1$ and denote $f(w_t)-f^*$ as $\Delta_t$, we have
    \begin{align}
        \Delta_{t+1}\leq \Delta_t -\gamma\mu^2\Delta_t^2
        \label{gd recursion}
    \end{align}
    The rest of proof is exactly the same as proof for Theorem \ref{beta GD convergence} Case I, we omit discussion here. As a result, one can show \eqref{gd recursion} converges to a $\epsilon$-stationary point after $\mathcal{O}(\frac{G}{\epsilon})$ iterations.
\end{proof}

\newpage
\section{Proof of Lemma \ref{stochastic gradient bound}}\label{Appendix H}
%Before proving Lemma \ref{stochastic gradient bound}, we first show that for any mini-batch sample $\xi'$, Assumption \ref{assum6} always holds.
%\begin{lemma}
%   For mini-batch stochastic gradient estimator satisfying Assumption \ref{assum6}, denote $\delta_{B}(w)$ as the approximation error $\delta_{B}(w) = \frac{1}{B}\sum_{i=1}^{B}\nabla f_i(w) - \nabla F(w)$, we have the upper bound
%    \begin{align}
%        \big \|\delta_B(w_t) \big \|\leq \tau_1\big \|\nabla F(w_t) \big \|+\tau_2.
%        \label{determinbound}
%    \end{align}   
%\end{lemma}
%\begin{proof}
%    When $B=1$, \eqref{determinbound} directly follows from Assumption \ref{assum6}.

%    When $B>1$, we have
%    \begin{align}
%        \big \|\delta_{B}(w_t)\big \| =&\big \|\frac{1}{B}\sum_{i=1}^{B}\delta_i(w_t)\big \| \leq \frac{1}{B} \sum_{i=1}^{B} \|\delta_i \|\leq \tau_1 \|\nabla F(w_t) \|+\tau_2 \nonumber
%    \end{align}
%    where the first inequality uses triangle inequality and convexity of squared L2-norm; the second inequality uses the assumption \eqref{affinebound}.
%\end{proof}
%The motivation of inducing Lemma \ref{stochastic gradient bound} is to construct a lower bound for the term $A\|\nabla f_{B_i}(w) \|+\delta$ with the proper choices of $A$ and $\delta$.
%\begin{lemma}
%    For mini-batch stochastic gradient estimator satisfying assumption \ref{assum6}, denote $\delta_{B}(w)$ as the approximation error $\delta_{B}(w) = \frac{1}{B}\sum_{i=1}^{B}\nabla f_i(w) - \nabla F(w)$, we have the upper bound
%    \begin{align}
%        \big \|\delta_B(w) \big \|\leq \frac{1}{\sqrt{B}}(\tau_1\big \|\nabla F(w) \big \|+\tau_2).
%        \label{determinbound}
%    \end{align}   
%\end{lemma}
%\begin{proof}
%The proof follows from applying Jensen's inequality for L2 norm.
%    \begin{align}
%        \big \|\delta_{B}(w)\big \| =&\big \|\frac{1}{B}\sum_{i=1}^{B}\delta_i(w)\big \| \nonumber\\
        %=& \Big(\big \|\frac{1}{B}\sum_{i=1}^{B}\delta_i(w)\big \|_2^{2}\Big)^{\frac{1}{2}} \nonumber\\
%        & = \frac{1}{B}\Big(\big \|\sum_{i=1}^{B}\delta_i(w)\big \|_2^{2}\Big)^{\frac{1}{2}} \nonumber \\
%        & \leq \frac{1}{B}\Big(\sum_{i=1}^{B}\big \|\delta_i(w)\big \|_2^{2}\Big)^{\frac{1}{2}} \nonumber \\
%        & \leq \frac{1}{B }
%        \Big( \sum_{i=1}^{B} (\tau_1 \big \| \nabla F(w)\big \|+\tau_2)^2\Big)^{\frac{1}{2}} \nonumber\\
%        =& \frac{1}{B}\cdot \sqrt{B}(\tau_1\big \|\nabla F(w) \big \|+\tau_2) \nonumber \\
%        & = \frac{1}{\sqrt{B}}(\tau_1 \big \|\nabla F(w) \big \|+\tau_2).\nonumber
%    \end{align}
%    where the first inequality uses Jensen's inequality and convexity of squared L2-norm; the second inequality uses the assumption, \eqref{affinebound}.
%\end{proof}
%This fact leads to 
%\begin{align}
%    \big \|\nabla f_{\xi}(w) \big \| \leq (\frac{\tau_1}{\sqrt{B}}+1) \big \|\nabla F(w) \big \|+\frac{\tau_2}{\sqrt{B}}.
%    \label{determinebound2}
%\end{align}
\mainlemmatwo*
\begin{proof}\label{Proof of Lemma 2}
    Based on Assumption \ref{assum6}, we have
    \begin{align}
        \| \nabla F(w_t)-\nabla f_{\xi'}(w_t)\| \leq \tau_1 \|\nabla F(w_t) \|+\tau_2,\nonumber
    \end{align}
    which, by triangle inequality, further implies that
    \begin{align}
        \big| \| \nabla F(w_t) \| -\|\nabla f_{\xi'}(w_t)\|\big|\leq 
        \| \nabla F(w_t)-\nabla f_{\xi'}(w_t)\|\leq \tau_1 \|\nabla F(w_t) \|+\tau_2.\nonumber
    \end{align}
    % Following above result, we must have
    % \begin{align}
    %     \| \nabla F(w_t) \| -\| \nabla f_{\xi'}(w_t)\|\leq \tau_1 \| \nabla F(w_t)\|+\tau_2.
    % \end{align}
    Rearranging the above inequality yields that 
    \begin{align}
        \|\nabla F(w_t) \|\leq \frac{1}{1-\tau_1}\|\nabla f_{\xi'}(w_t) \|+\frac{\tau_2}{1-\tau_1}.\nonumber
    \end{align}
    For the upper bound of $\|\nabla f_{\xi'}(w_t) \|$, it directly follows from Assumption \ref{assum6}.
\end{proof}
%Thus, for IAN-SGD, the proper choices of $A=\frac{1}{1-\tau_1}$, $\delta=\frac{\tau_2}{1-\tau_1}$.
%\begin{proof}\label{Proof of Lemma 2}
%    When $\big \| \nabla F(w)\big \|$ is large such that $\big \| \nabla F(w)\big \|\geq \frac{\tau_2}{\tau_1}$, then \eqref{determinbound} indicates 
%    \begin{align}
%    \big \|\delta_{{B}}(w) \big \|\leq \frac{2\tau_1\big \|\nabla F(w) \big \|}{\sqrt{B}}.\nonumber
%    \end{align}
%    In this case, if we choose $B= 16\tau_1^2$, we have $\big \|\delta_{B}(w) \big \|\leq \frac{1}{2}\big \|\nabla F(w) \big \|$. Since in this case, we assume, $\big \|\nabla F(w) \big \|\geq \frac{\tau_2}{\tau_1}\geq \frac{\tau_2}{2\tau_1}$, we have 
%    \begin{align}
%    \Big | \big \|\nabla f_{\xi}(w_t) \big \|-\big \|\nabla F(w) \big \|\Big | \leq \frac{1}{2}\big \|F(w) \big \|,\nonumber 
     %\leq \frac{3}{2}\big \|\nabla F(w) \big \|+\frac{\tau_2}{2\tau_1} \leq 2\big \|\nabla F(w) \big \|,\nonumber
%    \end{align}
%    which is equivalent as
%    \begin{align}
%        \big \|\nabla f_{\xi}(w) \big \| - \big \|\nabla F(w) \big \|\geq -\frac{1}{2}\big \|\nabla F(w) \big \|\nonumber.
%    \end{align}
%    And this fact leads to
%    \begin{align}
%        \frac{1}{2}\big \|\nabla F(w) \big \|\leq \big \|\nabla f_{\xi}(w) \big \| \leq \big \|\nabla f_{\xi}(w) \big \|+\frac{\tau_2}{2\tau_1}. \nonumber
%    \end{align}
%    Re-organize the term gives us
%    \begin{align}
%        \big \|\nabla f_{\xi}(w) \big \| \ge \frac{1}{2}\big \|\nabla F(w) \big \| - \frac{\tau_2}{2\tau_1}. \nonumber
%    \end{align}
    %now we want to show with proper choice of batch size w.h.p, $\big \|\delta_{B'} \big \|\leq \frac{\tau_1}{\sqrt{B'}} \big \| \nabla F(w)\big \|$. From Vector Berstein inequality, by setting $\epsilon = \tau_1 \big \|\nabla F(w) \big \|$
    %\begin{align}
        %P(\big \|\delta(w) \big \|> \tau_1 \big \|\nabla F(w) \big \|)\leq& \exp(-B' \cdot \frac{\tau_1^2 \big \|\nabla F(w) \big \|^2}{8(\tau_1\big \| \nabla F(w)\big \|+\tau_2)^2}+\frac{1}{4})\\
        %\leq & \exp(-B' \cdot \frac{\tau_1^2 \big \|\nabla F(w) \big \|^2}{8\cdot 4 \tau_1^2\big \| \nabla F(w)\big \|^2}+\frac{1}{4})\\
        %=&\exp(-\frac{B'}{32}+\frac{1}{4})
    %\end{align}
    %In this case, $B' \sim \Theta(1) $\\
%    Similarly, when $\big \|\nabla F(w) \big \|\leq \frac{\tau_2}{\tau_1}$, for single stochastic sample, by assumption \ref{assum6}, we have $\big \|\delta(w) \big \|\leq 2\tau_2$, from \eqref{determinbound}, for mini-batch stochastic gradient estimator, we have
%    \begin{align}
%    \big \|\delta_{B}(w) \big \|\leq \frac{2\tau_2}{\sqrt{B}}. \nonumber
%    \end{align}
%    By setting $B = 16 \tau_1^2$, we have
%     $\big \|\delta_{B}(w) \big \|\leq \frac{\tau_2}{2\tau_1}$. This fact leads to
%     \begin{align}
%         \Big |\big \|\nabla f_{\xi}(w) \big \| - \big \|\nabla F(w) \big \| \Big |\leq \frac{\tau_2}{2\tau_1}, \nonumber
%     \end{align}
%     which is equivalent as
%     \begin{align}
%         \big \|\nabla f_{\xi}(w) \big \| - \big \|\nabla F(w) \big \| \geq -\frac{\tau_2}{2\tau_1} \nonumber.
%     \end{align}
%     Thus, we have
%    \begin{align}
%        \frac{1}{2}\big \|\nabla F(w) \big \|\leq& \big \| \nabla F(w)\big \| \nonumber \\
%        =& \big \| \nabla F(w)\big \| +\frac{\tau_2}{2\tau_1}- \frac{\tau_2}{2\tau_1} \nonumber \\ \nonumber
%        \leq& \big \| \nabla f_{\xi}(w) \big \| + \frac{\tau_2}{2\tau_1}, \nonumber %\leq \big \|\nabla F(w) \big \|+\frac{\tau_2}{2\tau_1}+ \frac{\tau_2}{2\tau_1}\leq \frac{2\tau_2}{\tau_1} \nonumber
%    \end{align}
%    which leads to 
%    \begin{align}
%        \big \|\nabla f_{\xi}(w) \big \| \ge \frac{1}{2}\big \|\nabla F(w) \big \| - \frac{\tau_2}{2\tau_1}. \nonumber
%    \end{align}
%    Combine above, we conclude by choosing $B= 16 \tau_1^2$, we always have 
%    \begin{align}
%        \big \|\nabla f_{\xi}(w) \big \| \ge \frac{1}{2}\big \|\nabla F(w) \big \| - \frac{\tau_2}{2\tau_1}.
%    \end{align}
    %Reorganize this inequality gives us desired result.
    %In this way, we can eliminate the term $\frac{1}{\|\nabla f_{\xi'}(w)\|^{\beta}}$. However, the descent lemma's polynomial order will be changed, which might affect the convergence rate.
%\end{proof}
\section{Lemma \ref{upper bound for term 1} and Proof}
%Before proceeding to the proof of Lemma \ref{upper bound for term 1}. We prove a technical remark for variance of $\delta_{B}(w)$ stated as following.
%\begin{remark}[Variance bound for mini-batch $\delta_{B}(w)$]
%For variance of $\delta(w)$, following the same logic above, we have
%    \begin{align}
%        Var(\big \|\delta_{B}(w) \big \|) =& \mathbb{E}\big \|\delta_{B}(w) \big \|^2 \nonumber\\
%        =& \mathbb{E}\Big((\frac{1}{B}\sum_{i=1}^{B}\delta_i(w))^T(\frac{1}{B}\sum_{i=1}^{B}\delta_j(w))\Big) \nonumber \\
%        =&\frac{1}{B^2} \mathbb{E}\Big[\sum_{i=1}^{B} \big \|\delta_i(w) \big \|^2\Big] \nonumber\\
%        \leq& \frac{1}{B}(\tau_1\big \|\nabla F(w) \big \|+\tau_2)^2 \nonumber\\
%        \leq& \frac{1}{B}(2\tau_1^2\big \|\nabla F(w) \big \|^2+ 4\tau_2^2), \nonumber
%    \end{align}
%where the first inequality is due to \eqref{affinebound}.
%Thus, it is equivalent as
%\begin{align}
%    \mathbb{E}_{\xi}\big[ \big\|\nabla f_{\xi}(w) \big \|^2 \big]\leq (\frac{2\tau_1^2}{B}+1) \big \|\nabla F(w) \big \|^2+\frac{4\tau_2^2}{B}.
%    %\label{2ndmomentumbound}
%\end{align}
%\end{remark}
\begin{lemma}\label{upper bound for term 1}
For $\frac{1}{2}\gamma^2\frac{(L_0+L_1\big \|\nabla F(w_t) \big \|^{\alpha})}{h_t^{2\beta}}\cdot{4\tau_2^2}$, we have the following upper bound
\begin{align}
    \frac{1}{2}\gamma^2\frac{(L_0+L_1\big \|\nabla F(w_t) \big \|^{\alpha})}{h_t^{2\beta}}\cdot{4\tau_2^2}\leq \frac{1}{2}\gamma^2(L_0+L_1)(1+4\tau_2^2)^2+ \frac{\gamma}{4h_t^{\beta}}\big \|\nabla F(w_t) \big \|^2.
    \label{stochastic noise bound}
\end{align}
\end{lemma}
\begin{proof}
        When $\big \|\nabla F(w_t) \big \|\leq \sqrt{1+{4\tau_2^2}/{(2\tau_1^2+1)}}$, we have $\big \|\nabla F(w_t) \big \|^{\alpha}\leq (1+{4\tau_2^2}/{(2\tau_1^2+1)})^{\frac{\alpha}{2}}$ for any $\alpha >0$. Then, we have
        \begin{align}
        &\frac{\gamma^2}{2}\frac{(L_0+L_1\big \|\nabla F(w_t) \big \|^{\alpha})}{h_t^{2\beta}}4\tau_2^2 \nonumber\\
        \leq& \frac{1}{2}\gamma^2(1+\frac{4\tau_2^2}{2\tau_1^2+1})^{\frac{\alpha}{2}}\frac{(L_0+L_1)}{h_t^{2\beta}}{4\tau_2^2}\nonumber \\
        \leq& \frac{1}{2}\gamma^2(L_0+L_1)(1+4\tau_2^2)^{\frac{\alpha}{2}}(1+4\tau_2^2) \nonumber\\
        \leq& \frac{1}{2}\gamma^2(L_0+L_1)(1+{4\tau_2^2})^{2},
        \label{stochastic noise bound 1}
    \end{align}
    where the first inequality uses $\|\nabla F(w_t) \|^{\alpha}\leq (1+{4\tau_2^2}/{(2\tau_1^2+1}))^{\alpha/2}$ and $(1+{4\tau_2^2}/{(2\tau_1^2+1)})^{\alpha/2}\geq1$; the second inequality uses the fact that $\frac{1}{h_t}\leq 1$ (so does $\frac{1}{h_t^{2\beta}}$) and upper bound ${4\tau_2^2}$ by 
    $1+4{\tau_2^2}$; the last inequality uses the fact that thus $(1+4{\tau_2^2})^{1+\alpha/2}\leq(1+4{\tau_2^2})^2 $.

    When $\big \|\nabla F(w_t) \big \| > \sqrt{1+{4\tau_2^2}/{(2\tau_1^2+1)}}$, we have
    %we have $\big \|\nabla F(w_t) \big \|^{2}\geq (1+{4\tau_2^2}/{(2\tau_1^2+1)})^2\geq ({4\tau_2^2}/{(2\tau_1^2+1))^2}$ for any $\alpha >0 $, Thus, for $\frac{\gamma^2}{2} \frac{L_1 \big \|\nabla F(w_t) \big \|^{\alpha}}{h_t^{2\beta}}\cdot {4\tau_2^2}$, we have
    \begin{align}
        &\frac{\gamma^2}{2} \frac{L_1 \big \|\nabla F(w_t) \big \|^{\alpha}}{h_t^{2\beta}}\cdot {4\tau_2^2} \nonumber\\
        %=& \frac{\gamma^2}{2} \frac{L_1 \big \|\nabla F(w_t) \big \|^{\alpha}}{h_t^{\beta}} \frac{1}{h_t^{\beta}}\cdot {4\tau_2^2}\nonumber \\
        =& \frac{\gamma^2}{2h_t^{\beta}} \frac{L_1 \big \|\nabla F(w_t) \big \|^{\alpha}}{h_t^{\beta}}\cdot {4\tau_2^2} \nonumber \\
        \overset{(i)}{\leq}&\frac{\gamma^2}{2h_t^{\beta}} \frac{L_1 \big \|\nabla F(w_t) \big \|^{\alpha}}{4L_1\gamma(2\tau_1^2+1)(\frac{1}{1-\tau_1}\|\nabla f_{\xi'}(w_t) \|+\frac{\tau_2}{1-\tau_1})^{\beta}}\cdot 4\tau_2^2 \nonumber \\
        =& \frac{\gamma^2}{2h_t^{\beta}} \frac{L_1 \big \|\nabla F(w_t) \big \|^{\alpha}}{4L_1\gamma(\frac{1}{1-\tau_1}\|\nabla f_{\xi'}(w_t) \|+\frac{\tau_2}{1-\tau_1})^{\beta}}\cdot \frac{4\tau_2^2}{(2\tau_1^2+1)} \nonumber \\
        \overset{(ii)}{<}& \frac{\gamma^2}{2h_t^{\beta}} \frac{L_1 \big \|\nabla F(w_t) \big \|^{\alpha}}{(4L_1\gamma)(\frac{1}{1-\tau_1}\big \|\nabla f_{\xi'}(w_t) \big \|+\frac{\tau_2}{1-\tau_1})^{\beta}} \cdot \big \|\nabla F(w_t) \big \|^2 \nonumber \\
        \overset{(iii)}{\leq}& \frac{\gamma^2}{2h_t^{\beta}} \frac{L_1 \big \|\nabla F(w_t) \big \|^{\alpha}}{(4L_1\gamma)\big \|\nabla F(w_t) \big \|^{\beta}} \big \| \nabla F(w_t)\big \|^2\nonumber \\
        =& \frac{\gamma^2}{2h_t^{\beta}}\frac{L_1}{4L_1\gamma \big \|\nabla F(w_t) \big \|^{\beta - \alpha}}\big \|\nabla F(w_t) \big \|^2 \nonumber \\
        \overset{(iv)}{<}& \frac{\gamma^2}{2h_t^{\beta}} \frac{L_1}{4L_1\gamma} \big \|\nabla F(w_t) \big \|^{2}\nonumber \\
        =& \frac{\gamma}{8h_t^{\beta}} \big \|\nabla F(w_t) \big \|^2, 
        %& \overset{(v)}{\leq }\frac{\gamma}{4h_t^{\beta}} \big \|\nabla F(w_t) \big \|^2,
        %&\frac{1}{2}\gamma^2\frac{(L_0+L_1\big \|\nabla F(w_t) \big \|^{\alpha})}{h_t^{2\beta}}\frac{\tau_2^2}{\tau_1^2} \nonumber \\
        %& \leq \frac{1}{2}\gamma^2 \frac{(2L_1 \big \| \nabla F(w_t)\big \|^{\alpha})}{h_t^{2\beta}}\frac{\tau_2^2}{\tau_1^2} \nonumber \\
        %=& \frac{\gamma^2L_1 \big \|\nabla F(w_t) \big \|^{\alpha}}{h_t^{2\beta}} \frac{\tau_2^2}{\tau_1^2}
        \label{stochastic noise bound 2}
    \end{align}
        where (i) uses the fact $\frac{1}{h_t^{\beta}}\leq \frac{1}{4L_1\gamma(2\tau_1^2+1)(\frac{1}{1-\tau_1}\|\nabla f_{\xi'}(w_t) \|+\frac{\tau_2}{1-\tau_1})^{\beta}}$;
        (ii) uses the fact that now $\|\nabla F(w_t) \|^2> 1+{4\tau_2^2}/{(2\tau_1^2+1)}\geq {4\tau_2^2}/{(2\tau_1^2+1)}$;
     (iii) uses the fact $\|\nabla F(w_t) \|\leq \frac{1}{1-\tau_1}\|\nabla f_{\xi'}(w_t) \|+\frac{\tau_2}{1-\tau_1}$; (iv) uses the fact $\big \|\nabla F(w_t) \big \| > \sqrt{1+{4\tau_2^2}/{(2\tau_1^2+1)}}\geq 1$ and thus $\frac{1}{\|\nabla F(w_t)  \|^{\beta-\alpha}}<1$.
    
    Similarly, when $\big \|\nabla F(w_t) \big \| > \sqrt{1+{4\tau_2^2}/{(2\tau_1^2+1})}$,  for $\frac{1}{2h_t^{2\beta}}\gamma^2 L_0 \cdot 4\tau_2^2$, we have
    \begin{align}
       &\frac{1}{2h_t^{2\beta}}\gamma^2 L_0 \cdot 4\tau_2^2 \nonumber\\
       =&\frac{\gamma}{2h_t^{2\beta}}(\gamma L_0) \cdot 4\tau_2^2 \nonumber\\
       \leq& \frac{\gamma}{2h_t^{\beta}}(\gamma L_0) \cdot 4\tau_2^2 \nonumber\\
       \leq& \frac{\gamma}{2h_t^{\beta}} \frac{1}{4\cdot(2\tau_1^2+1)}4\tau_2^2 \nonumber \\
       =& \frac{\gamma}{8h_t^{\beta}} \frac{4\tau_2^2}{2\tau_1^2+1} \nonumber \\
       \leq& \frac{\gamma}{8h_t^{\beta}} \Big(1+\frac{4\tau_2^2}{2\tau_1^2+1} \Big)\nonumber \\
       <& \frac{\gamma}{8h_t^{\beta}} \big \|\nabla F(w_t) \big \|^2,
       \label{stochastic noise bound 3}
    \end{align}
        where the first inequality uses the fact that $\frac{1}{h_t^{\beta}}\leq 1$; the second inequality uses the fact $\gamma \leq \frac{1}{4L_0(2\tau_1^2+1)}$; the third inequality uses the fact $\frac{4\tau_2^2}{2\tau_1^2+1}\leq 1+\frac{4\tau_2^2}{2\tau_1^2+1}$; the fourth inequality uses the fact that
        $\big \|\nabla F(w_t) \big \|^2 > 1+{4\tau_2^2}/{(2\tau_1^2}+1)$.
        %and $\frac{1}{h_t^{\beta}}\leq 1$, second inequality uses the fact $\gamma L_0\leq \frac{1}{4}$.
        Combining \eqref{stochastic noise bound 1}, \eqref{stochastic noise bound 2}, \eqref{stochastic noise bound 3} gives us the desired result.
\end{proof}


\section{Proof of Theorem \ref{Clip SGD converge}}
\label{Appendix G}
\mainthmtwo*
\begin{proof}\label{Proof of IAN-SGD}
By the descent lemma in \eqref{origindescent} and the update rule of IAN-SGD in \eqref{update rule}, we have
\begin{align}
    &F\big(w_{t+1}\big)-F\big(w_t\big)\nonumber\\
    & {\leq} \nabla F\big(w_t\big)^{\top}\big(w_{t+1}-w_t\big) +\frac{1}{2}\big(L_0+L_1\big\|\nabla F\big(w_t\big)\big\|^\alpha\big)\big\|w_{t+1}-w_t\big\|^2 \nonumber\\
    & = -\gamma  \frac{\nabla F\big(w_t\big)^{\top}\nabla f_{\xi}(w_t)}{h_t^{\beta} }  + \frac{1}{2}\gamma^2\big(L_0+L_1\big\|\nabla F\big(w_t\big) \big \|^{\alpha} \big)\frac{\big \|\nabla f_{\xi}(w_t) \big \|^2}{h_t^{2\beta}}.
    \label{stochastic descent}
\end{align}
Taking expectation over $\xi$,$\xi'$ and $w_t$, and by the independence between $\xi$ and $\xi'$, we have
\begin{align}
    &\mathbb{E}_{w_t}\Big[\big[\mathbb{E}_{\xi,\xi'}\big[F(w_{t+1})-F(w_t)|w_t \big] \Big]\nonumber\\
    \leq& \mathbb{E}_{w_t,\xi'}\Big[ \frac{ -\gamma \big \| \nabla F(w_t)\big \|^2 }{h_t^{\beta}}  + \frac{1}{2}\gamma^2(L_0 + L_1\big \|\nabla F(w_t) \big \|^\alpha) \frac{\mathbb{E}_{\xi}\big [\big \|\nabla f_{\xi}(w_t)\big \|^2|w_t \big]}{h_t^{2\beta}} \Big].\nonumber
\end{align}
%When the expectation is conditioned on $w_t$,  we simplify $\mathbb{E}_{\xi}\big [\| \nabla F(w_t) \|^2 |w_t\big]$ into $\|\nabla F(w_t) \|^2$ since $\nabla F(w_t)$ doesn't include randomness over $\xi$.

From Assumption \ref{assum6}, we have
\begin{align}
    \mathbb{E}_{\xi}\big [\big \|\nabla f_{\xi}(w_t)\big \|^2 |w_t \big]=&\mathbb{E}_{\xi}\big [\big \|\nabla f_{\xi}(w_t)-\nabla F(w_t)+\nabla F(w_t)\big \|^2 |w_t \big]\nonumber\\
    \overset{(i)}{\leq} & 2\mathbb{E}_{\xi}\big [\big \|\nabla f_{\xi}(w_t)-\nabla F(w_t)\big \|^2 |w_t \big]+2\big \| \nabla F(w_t) \big \|^2 \nonumber\\
    \overset{(ii)}{\leq} & 2(\tau_1\big \|\nabla F(w_t) \big \|+\tau_2)^2+2\big \| \nabla F(w_t)\big\|^2\nonumber \\
    \overset{(iii)}{\leq}& ({4\tau_1^2}+2) \big \|\nabla F(w_t) \big \|^2 + {4\tau_2^2},
    \label{second momentum}
\end{align}
where inequalities (i) and (iii) are direct applications of $(a+b)^2\leq 2a^2+2b^2$, and (ii) follows from Assumption \ref{assum6}.
%Let $B = 1$, above inequality reduces to
%\begin{align}
%     \mathbb{E}_{\xi}\big [\big \|\nabla f_{\xi}(w_t)\big \|^2 | w_t \big]\leq (2\tau_1^2+1) \big \|\nabla F(w) \big \|^2 +2{\tau_2^2}.\label{second momentum}
%\end{align}
Substituting \eqref{second momentum} into the above descent lemma, we have
\begin{align}
    &\mathbb{E}_{w_t, \xi'}\big[\mathbb{E}_{\xi}\big[F(w_{t+1})-F(w_t)|w_t \big]\big] \nonumber \\
    \leq& \mathbb{E}_{w_t, \xi'}\Big[ -\gamma \frac{
    \big \| \nabla F(w_t)\big \|^2}{h_t^{\beta}}+ \frac{1}{2}\gamma^2(L_0 + L_1\big \|\nabla F(w_t) \big \|^\alpha) \frac{\mathbb{E}_{\xi}\big [\big \|\nabla f_{\xi}(w_t)\big \|^2|w_t \big]}{h_t^{2\beta}} \Big]\nonumber\\
    \leq& \mathbb{E}_{w_t, \xi'}\Big[ -\gamma \frac{\big \| \nabla F(w_t)\big \|^2}{h_t^{\beta}} + \frac{1}{2}\gamma^2(L_0 + L_1\big \|\nabla F(w_t) \big \|^\alpha) \frac{(4\tau_1^2+2)\big \|\nabla F(w_t) \big \|^2 +4\tau_2^2 }{h_t^{2\beta}} \Big]\nonumber\\
    %& \leq  -\gamma \frac{\big \| \nabla F(w)\big \|^2}{\big(\big \|\nabla f_{\xi'}(w_t) \big \|+\frac{\tau_2}{\tau_1}\big)^{2\beta}} \nonumber\\
    %& \qquad \frac{1}{2}\gamma^2(L_0 + L_1\big \|\nabla F(w_t) \big \|^\alpha) \frac{2\big \|\nabla F(w_t) \big \|^2 +1 }{(\big \|\nabla f_{\xi'}(w_t) \big \|+\frac{\tau_2}{\tau_1})^{2\beta}}\\
    =& \mathbb{E}_{w_t, \xi'}\Big[ \Big(\frac{\gamma}{h_t^{\beta}}\big(-1+\gamma(2\tau_1^2+1)\cdot \frac{L_0+L_1 \big \|\nabla F(w_t) \big \|^{\alpha}}{h_t^{\beta}}\big)\Big) \|\nabla F(w_t) \big \|^2 + \frac{1}{2}\gamma^2\frac{L_0+L_1 \big \|\nabla F(w_t) \big \|^{\alpha}}{h_t^{2\beta}}{4\tau_2^2}\Big]. \label{eq: des_lem}
\end{align}
Next, we provide an upper bound of the term $\gamma(2\tau_1^2+1)\cdot \frac{L_0+L_1 \|\nabla F(w_t) \|^{\alpha}}{h_t^{\beta}}$ in the above equation.
By the clipping structure and step size rule, we have
$\frac{1}{h_t^{\beta}} = \min \big\{1, \frac{1}{4L_1\gamma(2\tau_1^2+1)(\frac{1}{1-\tau_1}\|\nabla f_{\xi'}(w_t)\|+\frac{\tau_2}{1-\tau_1})^{\beta}} \big \}\leq 1$ and $\gamma \leq \frac{1}{4L_0(2\tau_1^2+1)}$, which imply that
%\begin{align}
    %\frac{\gamma}{h_t^{\beta}}<\gamma
    %\label{setp size relation 1}
%\end{align}
\begin{align}
    (2\tau_1^2+1)\cdot\frac{\gamma L_0}{h_t^{\beta}}<(2\tau_1^2+1)\cdot\gamma L_0\leq& \frac{1}{4}, \label{step size relation1}
\end{align}
% \\
%     (2\tau_1^2+1)\cdot \frac{\gamma L_1 \big \|\nabla F(w_t) \big \|^{\alpha}}{h_t^{\beta}}\leq& \frac{1}{4}.
%     \label{step size relation2}
On the other hand, we have
\begin{align}
\allowdisplaybreaks
    \frac{1}{4} h_t^{\beta}\overset{(i)}{\geq}& \frac{1}{4} h_t^{\alpha} \nonumber\\
    \overset{}{=}& \frac{1}{4} (h_t^{\beta})^{\frac{\alpha}{\beta}} \nonumber \\
    \overset{(ii)}{\geq}&\frac{1}{4} (4\gamma L_1(2\tau_1^2+1))^{\frac{\alpha}{\beta}}\big(\frac{1}{1-\tau_1}\big \|\nabla f_{\xi'}(w_t) \big \| + \frac{\tau_2}{1-\tau_1} \big)^{\alpha} \nonumber \\
    \overset{(iii)}{\geq}&\frac{1}{4} (4\gamma L_1(2\tau_1^2+1))\big(\frac{1}{1-\tau_1}\big \|\nabla f_{\xi'}(w_t) \big \| + \frac{\tau_2}{1-\tau_1} \big)^{\alpha} \nonumber \\
    \overset{(iv)}{\geq}& \gamma L_1(2\tau_1^2+1) \big \|\nabla F(w_t) \big \|^{\alpha}, \label{step size relation2}
\end{align}
where (i) utilizes the fact $h_t\geq 1$ and $\beta \geq \alpha$; (ii) utilizes the fact that $h_t^{\beta} \geq 4L_1 \gamma(2\tau_1^2+1) (\frac{1}{1-\tau_1} \big \|\nabla f_{\xi'}(w_t) \big \|+\frac{\tau_2}{1-\tau_1})^{\beta}$; (iii) utilizes the fact that $\gamma \leq \frac{1}{4L_1(2\tau_1^2+1)}$, thus $(4\gamma L_1(2\tau_1^2+1))\leq (4\gamma L_1(2\tau_1^2+1))^{\alpha/\beta}$ since $\beta \in [\alpha,1]$; (iv) follows from Lemma \ref{stochastic gradient bound}.

Combining \eqref{step size relation1} and \eqref{step size relation2}, we conclude that $\gamma(2\tau_1^2+1)\cdot \frac{L_0+L_1 \|\nabla F(w_t) \|^{\alpha}}{h_t^{\beta}}\le \frac{1}{2}$, and \eqref{eq: des_lem} further reduces to 
\begin{align}
    &\mathbb{E}_{w_t}\Big[\mathbb{E}_{\xi, \xi'}\big[F(w_{t+1})-F(w_t)|w_t \big]\Big] \nonumber\\
    & \overset{}{\leq} \mathbb{E}_{w_t, \xi'}\Big[-\frac{\gamma}{2h_t^{\beta}}\big \|\nabla F(w_t) \big \|^2 + \frac{1}{2}\gamma^2 \frac{L_0+L_1 \big \|\nabla F(w_t) \big \|^{\alpha}}{h_t^{2\beta}}{4\tau_2^2}\Big] \nonumber\\
    & \overset{(i)}{\leq} \mathbb{E}_{w_t, \xi'}\Big[ -\frac{\gamma}{2h_t^{\beta}}\big \|\nabla F(w_t) \big \|^2 + \frac{1}{2}\gamma^2 (L_0+L_1) (1+4\tau_2^2)^2+\frac{\gamma}{4h_t^{\beta}}\big \|\nabla F(w_t) \big \|^2 \Big] \nonumber\\
    & = \mathbb{E}_{w_t, \xi'}\Big[ -\frac{\gamma}{4h_t^{\beta}}\big \|\nabla F(w_t) \big \|^2 + \frac{1}{2}\gamma^2 (L_0+L_1) (1+4\tau_2^2)^2 \Big],
    \label{refined descent lemma}
\end{align}
where (i) follows from \eqref{stochastic noise bound} in the auxiliary Lemma \ref{upper bound for term 1}.
Rearranging the above inequality yields that
\begin{align}
   &\mathbb{E}_{w_t, \xi'}\Big[ \frac{\gamma}{4h_t^{\beta}} \big \|\nabla F(w_t) \big \|^2 \Big] \leq \mathbb{E}_{w_t}\Big[\mathbb{E}_{\xi, \xi'
    }\big[F(w_t) - F(w_{t+1})|w_t \big]\Big] + \frac{1}{2}(L_0+L_1 )\gamma^2(1+4\tau_2^2)^2. \quad \forall t \in [T]
    \label{for telescope sum}
\end{align}
Next, we further lower bound $ \frac{ \| \nabla F(w_t) \|^2}{h_t^{\beta}} $ by $\| \nabla F(w_t)\|$ to eliminate the randomness in $h_t$. We have
\begin{align}
    &\gamma\frac{ \big \| \nabla F(w_t) \big \|^2}{h_t^{\beta}} \nonumber\\
    \overset{(i)}{=}&\gamma \min \Big\{1,\frac{1}{4L_1(2\tau_1^2+1)\gamma(\frac{1}{1-\tau_1}\big \|\nabla f_{\xi'}(w_t) \big \|+\frac{\tau_2}{1-\tau_1})^\beta}  \Big \} \big \|\nabla F(w_t) \big \|^2 \nonumber \\
    %&\overset{(iii)}{\geq} \mathbb{E}_{\xi'}\Big[ \gamma \min \Big\{1,\frac{1}{4L_1\gamma\big(\big[2\big \|\nabla f_{\xi'}(w_t) \big \|+\frac{\tau_2}{\tau_1}\big])^\beta}  \Big \} \Big] \big \|\nabla F(w_t) \big \|^2 \nonumber \\
    \overset{(ii)}{\geq}& \gamma \min \Big\{1,\frac{1}{4L_1\gamma(2\tau_1^2+1)\big(\underbrace{\frac{\tau_1+1}{1-\tau_1}}_{C_1}\big \|\nabla F(w_t) \big \|+\underbrace{\frac{2\tau_2}{1-\tau_1}}_{C_2} \big)^{\beta}}  \Big \} \big \|\nabla F(w_t) \big \|^2 \nonumber \\
    \overset{(iii)}{\geq}& \gamma \min \Big\{1,\frac{1}{4L_1\gamma(2\tau_1^2+1) \big(2C_1\big \|\nabla F(w_t) \big \|\big)^{\beta}}, \frac{1}{4L_1\gamma(2\tau_1^2+1)\big( 2C_2\big)^{\beta}}  \Big \} \big \|\nabla F(w_t) \big \|^2 \nonumber \\
    \overset{(iv)}{=}&\min \Big \{\gamma, \frac{1}{4L_1(2\tau_1^2+1)\big(2C_1\big \|\nabla F(w_t) \big \|\big)^{\beta}}, \frac{1}{4L_1(2\tau_1^2+1)(2C_2)^{\beta}} \Big\} \big \|\nabla F(w_t) \big \|^2 \nonumber \\
    \overset{(v)}{=}& \min \Big \{\gamma, \frac{1}{4L_1(2\tau_1^2+1)\big(2C_1\big \|\nabla F(w_t) \big \|\big)^{\beta}} \Big\} \big \|\nabla F(w_t) \big \|^2 \nonumber \\
    \overset{}{=}& \min \Big \{ \gamma \big \|\nabla F(w_t) \big \|^2, \frac{\big \|\nabla F(w_t) \big \|^{2-\beta}}{4(2C_1)^{\beta} L_1(2\tau_1^2+1)}\Big \},
    \label{LHS}
    %& \geq \mathbb{E}_{\xi'}\Big[\gamma \min \Big \{1, \frac{1}{(2\big \|\nabla f_{\xi'}(w_t) \big \|)^{\beta}}, \Big(\frac{\tau_1}{2\tau_2} \Big)^{\beta} \Big \} \Big]\big \|\nabla F(w_t) \big \|^2 \nonumber\\
    %& \geq \mathbb{E}_{\xi'}\Big[\gamma \min \Big \{1, \frac{1}{(2\big \|\nabla f_{\xi'}(w_t) \big \|)^{\beta}}, \Big(\frac{\tau_1}{2\tau_2} \Big)^{\beta} \Big \} \Big]\big \|\nabla F(w_t) \big \|^2
\end{align}
where (i) expands the expression of $\frac{1}{h_t^{\beta}}$; (ii) utilizes the second inequality of Lemma \ref{stochastic gradient bound} to upper bound $\|\nabla f_{\xi'}(w_t) \|$ by $(\tau_1+1) \|\nabla F(w_t) \|+\tau_2$; (iii) utilizes the fact $ \frac{1}{(a+b)^{\beta}}\geq \min \{\frac{1}{(2a)^{\beta}}, \frac{1}{(2b)^{\beta}} \}$ by setting $a=\frac{\tau_1+1}{1-\tau_1}\big \|\nabla F(w_t) \big \|$, $b = \frac{2\tau_2}{1-\tau_1}$, $\beta\geq0$; (iv) substitute $\gamma$ inside the minimum operator. 
Moreover, (v) follows from the step size rule $\gamma= \frac{1}{8L_1(2\tau_1^2+1)(2\tau_2/(1-\tau_1))^{\beta}}\leq \frac{1}{4L_1(2\tau_1^2+1)(4\tau_2/(1-\tau_1))^{\beta}}$. 
% Thus, the third term $\frac{1}{4L_1(2\tau_1^2+1)(4\tau_2/(1-\tau_1))^{\beta}}$ can be directly deleted. This fact leads to lowerbound (v).
Substituting \eqref{LHS} into \eqref{for telescope sum}
and summing over $t$ from $0$ to $T-1$, we obtain that
\begin{align}
        &\sum_{t=0}^{T-1} \mathbb{E}_{w_t}\Big[\min \Big \{ \frac{\gamma}{4} \big \|\nabla F(w_t) \big \|^2, \frac{\big \|\nabla F(w_t) \big \|^{2-\beta}}{16L_1(2C_1)^{\beta}(2\tau_1^2+1)}\Big \} \Big]\nonumber \\
        %=& \mathbb{E}_{w_t}\Big[ \sum_{t=0}^{T-1} \min \Big \{ \frac{\gamma}{4} \big \|\nabla F(w_t) \big \|^2, \frac{\big \|\nabla F(w_t) \big \|^{2-\beta}}{16L_1(\tau_1+1)}\Big \} \Big]\nonumber \\
        \leq& \sum_{t=0}^{T-1}\mathbb{E}_{w_t,\xi'}\Big[ \frac{\gamma}{4h_t^{\beta}} \big \|\nabla F(w_t) \big \|^2 \Big] \nonumber \\
        \leq& \sum_{t=0}^{T-1}\mathbb{E}_{w_t}\Big[\mathbb{E}_{\xi,\xi'}\big[F(w_t) - F(w_{t+1})|w_t \big]\Big]\nonumber+T\frac{1}{2}\gamma^2(L_0+L_1)(1+4\tau_2^2)^2.
\end{align}
By step size rule $\gamma \leq \frac{1}{\sqrt{T}}$, the above inequality becomes
\begin{align}
    &\sum_{t=0}^{T-1}\mathbb{E}_{w_t}\Big[\min \Big \{ \frac{\gamma}{4} \big \|\nabla F(w_t) \big \|^2, \frac{\big \|\nabla F(w_t) \big \|^{2-\beta}}{16L_1(2C_1)^{\beta}(2\tau_1^2+1)}\Big \}\Big]\nonumber  \leq  F(w_0)-F^*+\frac{1}{2}(L_0+L_1)(1+4{\tau_2^2})^2.
\end{align}
Denote $K=\{t |t\in [T] \text{ such that } \gamma \big \|\nabla F(w_t) \big \|^2 \leq \frac{\|\nabla F(w_t) \|^{2-\beta}}{4L_1(2C_1)^{\beta}(2\tau_1^2+1)} \}$, then the above inequality implies that
\begin{align}
    &\sum_{t\in K} \mathbb{E}_{w_t}\Big[\frac{\gamma}{4}\big \|\nabla F(w_t) \big \|^2 \Big]  \leq F(w_0) - F^* + \frac{1}{2}(L_0+L_1)(1+4\tau_2^2)^2,\label{K's summation}
\end{align}
and
\begin{align}
    &\sum_{t\in K^c}\mathbb{E}_{w_t}\Big[ \frac{\big \|\nabla F(w_t) \big \|^{2-\beta}}{16L_1(2C_1)^{\beta}(2\tau_1^2+1)}\Big] \leq F(w_0) - F^* + \frac{1}{2}(L_0+L_1)(1+4\tau_2^2)^2. \label{Kc's summation}
\end{align}
Now denote RHS as $\Lambda := F(w_0)-F^* + \frac{1}{2}(L_0+L_1)(1+4\tau_2^2)^2$, then we have
\begin{align}
    &\mathbb{E}_{w_t}\Big[\min_{t\in T} \big \|\nabla F(w_t) \big \| \Big]\nonumber\\
    \leq& \mathbb{E}_{w_t}\Big[\min \Big \{\frac{1}{|K|} \sum_{t\in K} \big \|\nabla F(w_t) \big \|, \frac{1}{|K^c|}\sum_{t \in K^c}\big \| \nabla F(w_t) \big \| \Big \}\Big] \nonumber\\
    \overset{(i)}{\leq}& \mathbb{E}_{w_t}\Big[ \min \Big \{\sqrt{\frac{1}{|K|} \sum_{t\in K}^{} \big \|\nabla F(w_t) \big \|^2 }, \Big(\frac{1}{|K^c|}\sum_{t\in K^c}^{}\big \| \nabla F(w_t) \big \|^{2-\beta} \Big)^{\frac{1}{2-\beta}}\Big\}\Big] \nonumber\\
    \overset{(ii)}{\leq}&\max\Big\{\mathbb{E}_{w_t}\sqrt{\frac{1}{|K|} \sum_{t\in K}^{} \big \|\nabla F(w_t) \big \|^2 }, \mathbb{E}_{w_t}\Big(\frac{1}{|K^c|}\sum_{t\in K^c}^{}\big \| \nabla F(w_t) \big \|^{2-\beta} \Big)^{\frac{1}{2-\beta}} \Big\} \nonumber\\
    \overset{(iii)}{\leq}&\max\Big\{\sqrt{\frac{1}{|K|} \sum_{t\in K}^{} \mathbb{E}_{w_t}\big \|\nabla F(w_t) \big \|^2 }, \Big(\frac{1}{|K^c|}\sum_{t\in K^c}^{} \mathbb{E}_{w_t}\big \| \nabla F(w_t) \big \|^{2-\beta} \Big)^{\frac{1}{2-\beta}} \Big\} \nonumber\\
    \overset{(iv)}{\leq}& \max \Big \{\sqrt{(4\Lambda) \frac{4(L_0+L_1)(2\tau_1^2+1)+\sqrt{T}+8L_1(2\tau_1^2+1)((2\tau_2)/(1-\tau_1))^{\beta}}{T}},\nonumber\\
    &\qquad \qquad \qquad \qquad \qquad \qquad \quad \big(\Lambda \frac{32L_1(2\tau_1^2+1)((2\tau_1+2)/1-\tau_1)^{\beta}}{T}\big)^{\frac{1}{2-\beta}} \Big \}, \nonumber
\end{align}
where (i) uses the concavity $\psi(\cdot)=y^{\frac{1}{2}}$ and $\psi(\cdot)=y^{\frac{1}{2-\beta}}$ and Jensen's inequality $\psi(\frac{\sum_i x_i}{n})\geq \frac{\sum_{i}\psi(x_i)}{n}$
applied on $\| \nabla F(w_t)\|$; (ii) uses the fact $\mathbb{E}[\min\{A,B\}]\leq \mathbb{E}[A]\leq \max\{\mathbb{E}[A], \mathbb{E}[B]\}$; (iii) applies Jensen's inequality over concave function $\Psi(\cdot)=y^{\frac{1}{2}}$ and $\Psi(\cdot)=y^{\frac{1}{2-\beta}}$;
(iv) is due to re-organize \eqref{K's summation}, \eqref{Kc's summation} and set $K={T}$, $K^c=\frac{T}{2}$ (Since we require either $K$ or $K^c$ should be larger than $\frac{T}{2}$ to guarantee that $K+K_c=T$).

Thus, in order to find a point satisfies
\begin{align}
    Pr(\min_{t\in[T]}\| \nabla F(w_t) \| \geq \epsilon)\leq \frac{1}{2}, \nonumber 
\end{align}
We must have $\mathbb{E}_{w_t}[\min_{t\in [T]} \| \nabla F(w_t) \| ]\leq \frac{\epsilon}{2}$.

This indicates that the RHS of the above inequality should be smaller than $\frac{\epsilon}{2}$, this implies $T$ satisfies
\begin{align}
    T \geq \Lambda \max \Big\{\frac{256\Lambda}{\epsilon^4}, \frac{64L_1(2\tau_1^2+1)(2+2\tau_1)^{\beta}}{(1-\tau_1)^{\beta}\epsilon^{2-\beta}}, (2\tau_1^2+1)\cdot\frac{64(L_0+L_1)+128L_1((2\tau_2)/(1-\tau_1))^{\beta}}{\epsilon^2} \Big\}.
\end{align}
\end{proof}


\section{IAN-SGD under generalized P{\L} condition}\label{Appendix J}
\mainlemmathree*
\begin{proof}
    When chosen target error $0<\epsilon \leq \min\{1,1/2\mu\}$ and
    learning rate satisfy
    \begin{align}
        \gamma \leq (2\mu\epsilon)^{\frac{4-2\beta}{\rho}}\cdot \frac{1}{4L_0(2\tau_1^2+1)}\leq& \frac{1}{4L_0(2\tau_1^2+1)} \nonumber\\
        \gamma \leq(2\mu\epsilon)^{\frac{4-2\beta}{\rho}}\cdot\frac{1}{4L_1((2\tau_1+2)/(1-\tau_1))(2\tau_1^2+1)}\leq& \frac{1}{4L_1(2\tau_1^2+1)} \nonumber\\
        \gamma \leq (2\mu\epsilon)^{\frac{4-2\beta}{\rho}}\cdot \frac{1}{8L_1(2\tau_1^2+1)(2\tau_2/(1-\tau_1))^{\beta}}\leq& \frac{1}{8L_1(2\tau_1^2+1)(2\tau_2/(1-\tau_1))^{\beta}},
    \end{align}
    the derivations from \eqref{stochastic descent} to \eqref{for telescope sum} stated in Proof \ref{Appendix H} are still valid, we omit discussions of detailed derivation here. 

    Next, we lower bound $ \frac{ \| \nabla F(w_t) \|^2}{h_t^{\beta}} $ by $\| \nabla F(w_t)\|$ to eliminate the randomness in $h_t$. For any $\beta>\alpha>0$, we have
\begin{align}
    &\gamma\frac{ \big \| \nabla F(w_t) \big \|^2}{h_t^{\beta}} \nonumber\\
    \overset{(i)}{=}&\gamma \min \Big\{1,\frac{1}{4L_1(2\tau_1^2+1)\gamma(\frac{1}{1-\tau_1}\big \|\nabla f_{\xi'}(w_t) \big \|+\frac{\tau_2}{1-\tau_1})^\beta}  \Big \} \big \|\nabla F(w_t) \big \|^2 \nonumber \\
    %&\overset{(iii)}{\geq} \mathbb{E}_{\xi'}\Big[ \gamma \min \Big\{1,\frac{1}{4L_1\gamma\big(\big[2\big \|\nabla f_{\xi'}(w_t) \big \|+\frac{\tau_2}{\tau_1}\big])^\beta}  \Big \} \Big] \big \|\nabla F(w_t) \big \|^2 \nonumber \\
    \overset{(ii)}{\geq}& \gamma \min \Big\{1,\frac{1}{4L_1\gamma(2\tau_1^2+1)\big(\underbrace{\frac{\tau_1+1}{1-\tau_1}}_{C_1}\big \|\nabla F(w_t) \big \|+\underbrace{\frac{2\tau_2}{1-\tau_1}}_{C_2} \big)^{\beta}}  \Big \} \big \|\nabla F(w_t) \big \|^2 \nonumber \\
    \overset{(iii)}{\geq}& \gamma \min \Big\{1,\frac{1}{4L_1\gamma(2\tau_1^2+1) \big(2C_1\big \|\nabla F(w_t) \big \|\big)^{\beta}}, \frac{1}{4L_1\gamma(2\tau_1^2+1)\big( 2C_2\big)^{\beta}}  \Big \} \big \|\nabla F(w_t) \big \|^2 \nonumber \\
    \overset{(iv)}{=}&\min \Big \{\gamma, \frac{1}{4L_1(2\tau_1^2+1)\big(2C_1\big \|\nabla F(w_t) \big \|\big)^{\beta}}, \frac{1}{4L_1(2\tau_1^2+1)(2C_2)^{\beta}} \Big\} \big \|\nabla F(w_t) \big \|^2 \nonumber \\
    \overset{(v)}{=}& \min \Big \{\gamma, \frac{1}{4L_1(2\tau_1^2+1)\big(2C_1\big \|\nabla F(w_t) \big \|\big)^{\beta}} \Big\} \big \|\nabla F(w_t) \big \|^2 \nonumber \\
    \overset{}{=}& \min \Big \{ \gamma \big \|\nabla F(w_t) \big \|^2 ,\frac{\big \|\nabla F(w_t) \big \|^{2-\beta}}{4(2C_1)^{\beta} L_1(2\tau_1^2+1)}\Big \}\nonumber\\
    \overset{(vi)}{\geq}& \min\{\gamma (L_1\gamma(2\tau_1^2+1))^{\beta/2}\|\nabla F(w_t) \|^{2-\beta}-\gamma^2L_1(2\tau_1^2+1),\frac{\big \|\nabla F(w_t) \big \|^{2-\beta}}{4(2C_1)^{\beta} L_1(2\tau_1^2+1)}\} \nonumber\\
    \overset{(vii)}\geq& \min\{\gamma^{{3}/{2}} (L_1(2\tau_1^2+1))^{1/2}\|\nabla F(w_t) \|^{2-\beta}-\gamma^2L_1(2\tau_1^2+1),\frac{\big \|\nabla F(w_t) \big \|^{2-\beta}}{4(2C_1)^{\beta} L_1(2\tau_1^2+1)}\} \nonumber\\
    \overset{(viii)}{\geq}& \min\{\gamma^{{3}/{2}} (L_1(2\tau_1^2+1))^{1/2}\|\nabla F(w_t) \|^{2-\beta},\frac{\big \|\nabla F(w_t) \big \|^{2-\beta}}{4(2C_1)^{\beta} L_1(2\tau_1^2+1)}\}-\gamma^2L_1(2\tau_1^2+1)\nonumber \\
    \overset{(x)}{\geq}&\gamma^{{3}/{2}} (L_1(2\tau_1^2+1))^{1/2}\|\nabla F(w_t) \|^{2-\beta}-\gamma^2L_1(2\tau_1^2+1)
    \nonumber\\
    \overset{(ix)}{\geq}&\gamma^{{3}/{2}} (L_1(2\tau_1^2+1))^{1/2}(2\mu \Delta_t)^{(2-\beta)/\rho}-\gamma^2L_1(2\tau_1^2+1)
\end{align}
where (i)-(v) follows the same arguments stated in Proof \ref{Appendix H}. (vi) applies \eqref{Technical Lemma's inequality} by setting $C=(L_1\gamma(2\tau_1^2+1)^{\beta/2})$ (since $L_1\gamma(2\tau_1^2+1)<1$), $x=\|\nabla F(w_t) \|$, $\Delta=\beta$, $\omega' = 2,\omega = 2-\beta$; (vii) follows from $(L_1\gamma(2\tau_1^2+1))^{\beta/2}\geq (L_1\gamma(2\tau_1^2+1))^{1/2}$ as $(L_1\gamma(2\tau_1^2+1))<1$; (viii) follows from fact for any $b\geq 0$, $\min \{a-b,c\}\geq \min\{a,c\}-b$;(x) follows from fact $\gamma \leq \frac{1}{4L_1((2\tau_1+2)/(1-\tau_1))(2\tau_1^2+1)}\leq \frac{1}{4L_1(2C_1)^{\beta}(2\tau_1^2+1)}$. Thus, $L_1\gamma(2\tau_1^2+1)<1$ and $\gamma^{3/2}(L_1(2\tau_1^2+1))^{1/2}\leq \gamma\leq  \frac{1}{4(2C_1)^{\beta}(2\tau_1^2+1)}$ always holds; (ix) substitutes assumption \ref{assum2} to replace $\|\nabla F(w_t) \|$ by $(2\mu\Delta_t)^{1/\rho}$.

When $\alpha=0,\beta>0$, above argument still applies, next, we consider the case when $\alpha=\beta=0$, we have
\begin{align}
    \gamma \|\nabla F(w_t) \|^2\geq& \gamma(L_1\gamma(2\tau_1^2+1))^{1/2} \|\nabla F(w_t) \|^2 \nonumber\\
    \geq & \gamma (L_1\gamma(2\tau_1^2+1))^{1/2}\|\nabla F(w_t) \|^2 - \gamma^2 L_1(2\tau_1^2+1)\nonumber \\
    \geq &\gamma (L_1\gamma(2\tau_1^2+1))^{1/2}(2\mu\Delta_t)^{2/\rho} - \gamma^2 L_1(2\tau_1^2+1)\nonumber\\
    = & \gamma^{3/2}(L_1(2\tau_1^2+1))(2\mu\Delta_t)^{2/\rho}-\gamma^2 L_1(2\tau_1^2+1),
\end{align}
where the first inequality is due to $\gamma\leq \frac{1}{L_1(2\tau_1^2+1)}$ and the last inequality substitutes assumption \ref{assum2} to replace $\|\nabla F(w_t) \|$ by $(2\mu\Delta_t)^{1/\rho}$.

Combining two cases, we conclude that for any $\beta\geq 0$, we always have
\begin{align}
    \gamma\frac{ \big \| \nabla F(w_t) \big \|^2}{h_t^{\beta}}\geq \gamma^{3/2}(L_1(2\tau_1^2+1))(2\mu\Delta_t)^{(2-\beta)/\rho}-\gamma^2 L_1(2\tau_1^2+1)
    \label{LHS stochastic PL}
\end{align}

Substituting \eqref{LHS stochastic PL} into \eqref{for telescope sum} and re-arranging terms, we have 
\begin{align}
       &\mathbb{E}_{w_t}\Big[\frac{\gamma^{{3}/{2}} (L_1(2\tau_1^2+1))^{1/2}(2\mu \Delta_t)^{(2-\beta)/\rho}}{4}\Big]\nonumber\\ \leq& \mathbb{E}_{w_t}\Big[\big[F(w_t) - F(w_{t+1})\big]\Big] + \frac{1}{2}(L_0+L_1 )\gamma^2(1+4\tau_2^2)^2+\frac{\gamma^2L_1(2\tau_1^2+1)}{4}.\nonumber\\
       =& \mathbb{E}_{w_t}\Big[\big[F(w_t) - F(w_{t+1})\big]\Big] + \frac{1}{2}\gamma^2\big((L_0+L_1 )(1+4\tau_2^2)^2+\frac{L_1(2\tau_1^2+1)}{2}\big)
\end{align}
Substituting $F^*$ on both sides and re-arranging terms yields that
\begin{align}
    \mathbb{E}_{w_{t+1}}[\Delta_{t+1}]\leq \mathbb{E}_{w_t}[\Delta_t-\frac{\gamma^{3/2}(L_1(2\tau_1^2+1))^{1/2}(2\mu\Delta_t)^{(2-\beta)/\rho}}{4}+\frac{1}{2}\gamma^2\big((L_0+L_1 )(1+4\tau_2^2)^2+{L_1(\tau_1^2+1/2)})].\nonumber
\end{align}
Since $\gamma < \frac{L_1(2\tau_1^2+1)(2\mu\epsilon)^{(4-2\beta)/{\rho}}}{16((L_0+L_1)(1+4\tau_2^2)+L_1(\tau_1^2+1/2))^2}$, the above inequality further implies
\begin{align}
    \mathbb{E}_{w_{t+1}}[\Delta_{t+1}]\leq \mathbb{E}_{w_t}[\Delta_t - \frac{\gamma^{3/2}(L_1(2\tau_1^2+1))^{1/2}(2\mu\Delta_t)^{(2-\beta)/\rho}}{4}+ \frac{\gamma^{3/2}(L_1(2\tau_1^2+1))^{1/2}(2\mu\epsilon)^{(2-\beta)/\rho}}{8}],
    \label{descent PL2 stochastic}
\end{align}
which gives the desired result.
%\eqref{descent PL2 stochastic} is similar as \eqref{descentPL}. And the convergence analysis of \eqref{descent PL2 stochastic} follows naturally from Proof \ref{Appendix E} depending on value of $\rho$, $\beta$. We omit discussions here and refer interested readers to check proof technique stated in Proof \ref{Appendix E}.
\end{proof}

\end{document}
